\documentclass[11pt]{article}
\usepackage[margin=1in]{geometry}
\usepackage{amsmath,amssymb}
\usepackage{enumitem}
\usepackage{titlesec}
\usepackage{fancyhdr}
\usepackage{tikz}
\usepackage{multicol}
\usepackage{booktabs}
\usepackage{array}
\usepackage{tabularx}

\pagestyle{fancy}
\fancyhf{}
\fancyhead[C]{The Hope Science Academy -- Mathematics MCQs \& Questions (Class 9th, Full Book)}
\fancyfoot[C]{\thepage}
\setlength{\headheight}{25.9pt}

\titleformat{\section}{\large\bfseries}{}{0em}{}

% ---- qenv: MCQ environment used for Parts 1,2,4,7 (Q.1-98,99-203,284-359,521-543) ----
\newenvironment{qenv}[1]{%
  \par\medskip\noindent\textbf{#1.}\ \ignorespaces
}{\par}

\newlist{options}{enumerate}{1}
\setlist[options,1]{label=(\Alph*),leftmargin=2.5em,itemsep=0pt,topsep=2pt}

% ---- mcqB: MCQ list used for Part 5 (Q.360-443) ----
\newlist{mcqB}{enumerate}{1}
\setlist[mcqB,1]{label=\textbf{\arabic*.}, leftmargin=*, itemsep=8pt, topsep=6pt}
\newcommand{\opts}[4]{
\begin{enumerate}[label=(\Alph*), itemsep=1pt, topsep=2pt, leftmargin=1.5cm]
\item #1
\item #2
\item #3
\item #4
\end{enumerate}
}

% ---- mcqC: MCQ command used for Part 6 (Q.444-520) ----
\newcounter{qnum}
\setcounter{qnum}{443}
\newcommand{\mcqC}[5]{%
  \refstepcounter{qnum}%
  \noindent\textbf{\theqnum. #1}\par
  \begin{multicols}{2}
  \noindent A) #2 \\
  B) #3 \\
  C) #4 \\
  D) #5
  \end{multicols}
  \vspace{4pt}
}

% ---- shortq: Roman-numeral short questions, locally numbered per block ----
\newlist{shortq}{enumerate}{1}
\setlist[shortq,1]{label=\Roman*.,leftmargin=3.2em,itemsep=6pt,topsep=4pt}

% ---- sq: Roman-numeral short questions, persistent counter (Parts 9, 12) ----
\newlist{sq}{enumerate}{1}
\setlist[sq,1]{label=\textbf{\Roman*.}, leftmargin=*, itemsep=8pt, topsep=6pt}

% ---- wq: word-problem questions, persistent counter (Parts 15, 18) ----
\newlist{wq}{enumerate}{1}
\setlist[wq,1]{label=\textbf{\arabic*.}, leftmargin=*, itemsep=8pt, topsep=6pt}

% ---- detq: detailed questions Q3 (Part 14) ----
\newlist{detq}{enumerate}{1}
\setlist[detq,1]{label=\arabic*.,leftmargin=2.5em,itemsep=10pt,topsep=4pt}

\newcommand{\seg}[1]{\overline{#1}}

\title{\textbf{Mathematics -- Class 9th (Full Book)\\Complete Combined Question Bank}}
\author{The Hope Science Academy, Johar Town, Lahore}
\date{}

\begin{document}
\maketitle
\thispagestyle{fancy}
\tableofcontents
\clearpage

\section*{Part 1 -- MCQs (Q.1--543)}
\addcontentsline{toc}{section}{Part 1 -- MCQs (Q.1--543)}

\section*{Q1. Choose the correct answer.}

\begin{qenv}{1} The Sumerians used which numeral system:
\begin{options} \item Base 10 \item Base 60 \item Base 2 \item Base 12 \end{options} \end{qenv}

\begin{qenv}{2} The concept of zero was first developed by:
\begin{options} \item Egyptians \item Sumerians \item Indians \item Greeks \end{options} \end{qenv}

\begin{qenv}{3} Modern number systems like hexadecimal are based on:
\begin{options} \item Base 10 \item Base 16 \item Base 2 \item Base 60 \end{options} \end{qenv}

\begin{qenv}{4} Which of these is an irrational number:
\begin{options} \item $0.333\ldots$ \item $3.14159$ \item $1.5$ \item $2332$ \end{options} \end{qenv}

\begin{qenv}{5} A recurring decimal is always a:
\begin{options} \item Rational number \item Irrational number \item Integer \item Whole number \end{options} \end{qenv}

\begin{qenv}{6} Which number is rational?
\begin{options} \item $\sqrt{3}$ \item $\sqrt{16}$ \item $\sqrt{5}$ \item $\sqrt{7}$ \end{options} \end{qenv}

\begin{qenv}{7} The decimal expansion of $\dfrac{1}{7}$:
\begin{options} \item Terminating \item Non-terminating and non-recurring \item Non-terminating and recurring \item None of these \end{options} \end{qenv}

\begin{qenv}{8} Which property is illustrated by $a+b=b+a$:
\begin{options} \item Associative \item Identity \item Distributive \item Commutative \end{options} \end{qenv}

\begin{qenv}{9} The additive identity property for real numbers is:
\begin{options} \item $1$ \item $0$ \item $-1$ \item None of these \end{options} \end{qenv}

\begin{qenv}{10} Which property states $a(b+c)=ab+ac$:
\begin{options} \item Associative \item Distributive \item Inverse \item Closure \end{options} \end{qenv}

\begin{qenv}{11} The multiplicative inverse of $5$ is:
\begin{options} \item $-5$ \item $\dfrac{1}{5}$ \item $0$ \item $1$ \end{options} \end{qenv}

\begin{qenv}{12} If $a>b$ and $b>c$, then $a>c$ is called the \underline{\hspace{2cm}} property:
\begin{options} \item Reflexive \item Symmetric \item Transitive \item Trichotomy \end{options} \end{qenv}

\begin{qenv}{13} Which of the following is not an integer:
\begin{options} \item $-3$ \item $0$ \item $4.5$ \item $7$ \end{options} \end{qenv}

\begin{qenv}{14} Which of the following is an example of a rational number?
\begin{options} \item $\pi$ \item $\dfrac{3}{4}$ \item $\sqrt{2}$ \item $e$ \end{options} \end{qenv}

\begin{qenv}{15} What is an example of a recurring decimal:
\begin{options} \item $3.14159265\ldots$ \item $0.3333\ldots$ \item $2.718281828\ldots$ \item $0.49835555$ \end{options} \end{qenv}

\begin{qenv}{16} Which number system does $-3$ belong to:
\begin{options} \item Natural number \item Whole number \item Integers \item Irrational number \end{options} \end{qenv}

\begin{qenv}{17} Which of the following is not a real number:
\begin{options} \item $5$ \item $-2$ \item $\sqrt{7}$ \item $0$ \end{options} \end{qenv}

\begin{qenv}{18} Which number set does $0$ belong to:
\begin{options} \item Whole number \item Natural number \item Irrational number \item None of these \end{options} \end{qenv}

\begin{qenv}{19} Which of the following is an example of an infinite non-recurring decimal?
\begin{options} \item $0.25$ \item $0.333\ldots$ \item $5.75$ \item $\pi$ \end{options} \end{qenv}

\begin{qenv}{20} $0.\overline{6}$ is equivalent to:
\begin{options} \item $\dfrac{2}{3}$ \item $\dfrac{1}{6}$ \item $\dfrac{3}{5}$ \item $\dfrac{6}{99}$ \end{options} \end{qenv}

\begin{qenv}{21} Which is a binomial surd?
\begin{options} \item $\sqrt{3}$ \item $\sqrt{3}+\sqrt{5}$ \item $2\sqrt{3}$ \item $\sqrt{9}$ \end{options} \end{qenv}

\begin{qenv}{22} The conjugate of $\sqrt{7}-\sqrt{2}$:
\begin{options} \item $\sqrt{7}+\sqrt{2}$ \item $\sqrt{7}-\sqrt{2}$ \item $-\sqrt{7}+\sqrt{2}$ \item $-\sqrt{7}-\sqrt{2}$ \end{options} \end{qenv}

\begin{qenv}{23} Rationalize the denominator $\dfrac{1}{\sqrt{5}}$:
\begin{options} \item $\sqrt{5}$ \item $\dfrac{\sqrt{5}}{5}$ \item $\dfrac{1}{5}$ \item $5\sqrt{5}$ \end{options} \end{qenv}

\begin{qenv}{24} Which is not a surd?
\begin{options} \item $\sqrt{8}$ \item $\sqrt{12}$ \item $\sqrt{16}$ \item $\sqrt{18}$ \end{options} \end{qenv}

\begin{qenv}{25} The solution of $2x+5=15$ is:
\begin{options} \item $x=3$ \item $x=10$ \item $x=7.5$ \item $x=5$ \end{options} \end{qenv}

\begin{qenv}{26} If $\sqrt{x}=5$, then $x=?$:
\begin{options} \item $10$ \item $25$ \item $5$ \item $15$ \end{options} \end{qenv}

\begin{qenv}{27} The value of $\dfrac{\sqrt{12}\times\sqrt{3}}{\sqrt{4}}$:
\begin{options} \item $3$ \item $6$ \item $9$ \item $12$ \end{options} \end{qenv}

\begin{qenv}{28} If $a=3$, $b=4$, then $\sqrt{a^2+b^2}$:
\begin{options} \item $5$ \item $7$ \item $12$ \item $25$ \end{options} \end{qenv}

\begin{qenv}{29} Which is equivalent to $\dfrac{\sqrt{3}+\sqrt{2}}{\sqrt{3}-\sqrt{2}}$?
\begin{options} \item $5+2\sqrt{6}$ \item $1$ \item $\sqrt{6}$ \item $5-2\sqrt{6}$ \end{options} \end{qenv}

\begin{qenv}{30} If $x=1+\sqrt{2}$, then $x^2=?$:
\begin{options} \item $2\sqrt{2}$ \item $1+2\sqrt{2}$ \item $3$ \item $3+2\sqrt{2}$ \end{options} \end{qenv}

\begin{qenv}{31} If $a+b=10$ and $ab=21$, then $a^2+b^2=$:
\begin{options} \item $58$ \item $100$ \item $42$ \item $121$ \end{options} \end{qenv}

\begin{qenv}{32} The solution of $2^x=16$ is:
\begin{options} \item $6$ \item $8$ \item $2$ \item $4$ \end{options} \end{qenv}

\begin{qenv}{33} What is the square root of $16$:
\begin{options} \item $4$ \item $-4$ \item Both (A) and (B) \item None of these \end{options} \end{qenv}

\begin{qenv}{34} The value of $\sqrt{25}$:
\begin{options} \item $5$ \item $-5$ \item Both (A) and (B) \item None of these \end{options} \end{qenv}

\begin{qenv}{35} Convert $25^{\circ}\text{C}$ to Fahrenheit using $F=\dfrac{9}{5}C+32$:
\begin{options} \item $77^{\circ}$F \item $68^{\circ}$F \item $90^{\circ}$F \item $45^{\circ}$F \end{options} \end{qenv}

\begin{qenv}{36} A profit of $20\%$ on a cost price of Rs.\ 500 means the selling price is:
\begin{options} \item Rs.\ 600 \item Rs.\ 400 \item Rs.\ 550 \item Rs.\ 700 \end{options} \end{qenv}

\begin{qenv}{37} Compound interest on Rs.\ 2000 at $10\%$ for 2 years is:
\begin{options} \item Rs.\ 420 \item Rs.\ 200 \item Rs.\ 400 \item Rs.\ 240 \end{options} \end{qenv}

\begin{qenv}{38} If three consecutive integers sum to $51$, the largest number is:
\begin{options} \item $15$ \item $16$ \item $17$ \item $18$ \end{options} \end{qenv}

\begin{qenv}{39} A rectangle with sides $\sqrt{3}$ and $\sqrt{12}$ has an area of:
\begin{options} \item $\sqrt{15}$ \item $6$ \item $3\sqrt{3}$ \item $12$ \end{options} \end{qenv}

\begin{qenv}{40} If 5 pens cost Rs.\ 45, the cost of 8 pens is:
\begin{options} \item $100$ \item $60$ \item $90$ \item $72$ \end{options} \end{qenv}

\begin{qenv}{41} A number is increased by $20\%$ and becomes $60$. The original number is:
\begin{options} \item $50$ \item $48$ \item $60$ \item $72$ \end{options} \end{qenv}

\begin{qenv}{42} Two numbers are in the ratio $3:4$. Their sum is $35$. The numbers are:
\begin{options} \item $15, 20$ \item $12, 16$ \item $9, 12$ \item $18, 24$ \end{options} \end{qenv}

\begin{qenv}{43} A train travels $300$ km in $5$ hours. Its speed is:
\begin{options} \item $50$ km/h \item $60$ km/h \item $45$ km/h \item $65$ km/h \end{options} \end{qenv}

\begin{qenv}{44} The area of a square with side $\sqrt{5}$:
\begin{options} \item $5$ \item $10$ \item $25$ \item $20$ \end{options} \end{qenv}

\begin{qenv}{45} The sum of $\sqrt{0.25}+\sqrt{0.04}$:
\begin{options} \item $0.7$ \item $0.5$ \item $0.9$ \item $1.0$ \end{options} \end{qenv}

\begin{qenv}{46} The temperature $100^{\circ}$C is equal to how many Kelvin:
\begin{options} \item $273$ K \item $373$ K \item $473$ K \item $173$ K \end{options} \end{qenv}

\begin{qenv}{47} $\sqrt{7}$ is a/an:
\begin{options} \item Integer \item Rational number \item Irrational number \item Natural number \end{options} \end{qenv}

\begin{qenv}{48} $\pi$ and $e$ are:
\begin{options} \item Natural numbers \item Integers \item Rational numbers \item Irrational numbers \end{options} \end{qenv}

\begin{qenv}{49} If $n$ is not a perfect square, then $\sqrt{n}$ is a/an:
\begin{options} \item Rational number \item Natural number \item Integer \item Irrational number \end{options} \end{qenv}

\begin{qenv}{50} $\sqrt{3}+\sqrt{5}$ is a/an:
\begin{options} \item Whole number \item Integer \item Rational number \item Irrational number \end{options} \end{qenv}

\begin{qenv}{51} For all $x\in\mathbb{R}$, $x=x$ is called the:
\begin{options} \item Reflexive property \item Transitive property \item Symmetric property \item Trichotomy property \end{options} \end{qenv}

\begin{qenv}{52} Let $a,b,c\in\mathbb{R}$, then $a>b$ and $b>c \Rightarrow a>c$ is called the \underline{\hspace{2cm}} property:
\begin{options} \item Trichotomy \item Transitive \item Additive \item Multiplicative \end{options} \end{qenv}

\begin{qenv}{53} $2^x \times 8^x = 64$, then $x=$:
\begin{options} \item $\dfrac{3}{2}$ \item $\dfrac{3}{4}$ \item $\dfrac{5}{6}$ \item $\dfrac{2}{3}$ \end{options} \end{qenv}

\begin{qenv}{54} Let $a,b\in\mathbb{R}$, then $a=b$ and $b=a$ is called the:
\begin{options} \item Reflexive property \item Symmetric property \item Transitive property \item Additive property \end{options} \end{qenv}

\begin{qenv}{55} $\sqrt{75}+\sqrt{27}=$:
\begin{options} \item $9\sqrt{3}$ \item $\sqrt{102}$ \item $5\sqrt{3}$ \item $8\sqrt{3}$ \end{options} \end{qenv}

\begin{qenv}{56} The product of $(3+\sqrt{5})(3-\sqrt{5})$ is a/an:
\begin{options} \item Prime number \item Odd number \item Irrational number \item Rational number \end{options} \end{qenv}

\begin{qenv}{57} Convert $45{,}000$ to scientific notation:
\begin{options} \item $4.5\times10^4$ \item $4.5\times10^3$ \item $45\times10^3$ \item $0.45\times10^5$ \end{options} \end{qenv}

\begin{qenv}{58} The scientific notation of $0.000078$ is:
\begin{options} \item $7.8\times10^{-5}$ \item $78\times10^{-6}$ \item $0.78\times10^{-4}$ \item $7.8\times10^{5}$ \end{options} \end{qenv}

\begin{qenv}{59} Which is not in scientific notation:
\begin{options} \item $3.4\times10^6$ \item $0.5\times10^3$ \item $6.02\times10^{23}$ \item $1.2\times10^{-4}$ \end{options} \end{qenv}

\begin{qenv}{60} Convert $9.01\times10^5$ to ordinary form:
\begin{options} \item $9010$ \item $90100$ \item $901000$ \item $901$ \end{options} \end{qenv}

\begin{qenv}{61} The coefficient in $6.022\times10^{23}$ is:
\begin{options} \item $6$ \item $6.022$ \item $10^{23}$ \item $23$ \end{options} \end{qenv}

\begin{qenv}{62} $73\times10^3$ in scientific notation is:
\begin{options} \item $7.3\times10^2$ \item $7.3\times10^3$ \item $7.3\times10^4$ \item $73\times10^3$ \end{options} \end{qenv}

\begin{qenv}{63} Convert $6.7779\times10^3$ km to standard form:
\begin{options} \item $677.76$ km \item $67.779$ km \item $67779$ km \item $6777.9$ km \end{options} \end{qenv}

\begin{qenv}{64} The exponential form of $\log_5 25=2$ is:
\begin{options} \item $25^2=5$ \item $2^5=25$ \item $5^2=25$ \item $25^5=2$ \end{options} \end{qenv}

\begin{qenv}{65} The logarithmic form of $10^3=1000$ is:
\begin{options} \item $\log_{10}3=1000$ \item $\log_{3}10=1000$ \item $\log_{10}1000=3$ \item $\log_{3}1000=10$ \end{options} \end{qenv}

\begin{qenv}{66} Solve $\log_3 27 = x$:
\begin{options} \item $2$ \item $9$ \item $6$ \item $3$ \end{options} \end{qenv}

\begin{qenv}{67} Convert $7^0=1$ to logarithmic form:
\begin{options} \item $\log_7 1=0$ \item $\log_0 7=1$ \item $\log_7 0=1$ \item $\log_1 7=0$ \end{options} \end{qenv}

\begin{qenv}{68} The base of $\log_{10}100=2$ is:
\begin{options} \item $2$ \item $10$ \item $100$ \item $1$ \end{options} \end{qenv}

\begin{qenv}{69} The characteristic of $\log 0.0078$ is:
\begin{options} \item $-2$ \item $-3$ \item $3$ \item $2$ \end{options} \end{qenv}

\begin{qenv}{70} The mantissa of $\log 345$ (given $\log 345=2.5378$) is:
\begin{options} \item $2$ \item $0.5378$ \item $5378$ \item $2.5378$ \end{options} \end{qenv}

\begin{qenv}{71} If $\log x = \overline{3}.678$, the characteristic is:
\begin{options} \item $3$ \item $-3$ \item $-4$ \item $4$ \end{options} \end{qenv}

\begin{qenv}{72} If $\log x = 2.3010$, then $x=$:
\begin{options} \item $20$ \item $200$ \item $0.2$ \item $2000$ \end{options} \end{qenv}

\begin{qenv}{73} The base of the natural logarithm is:
\begin{options} \item $e$ \item $10$ \item $2$ \item $3$ \end{options} \end{qenv}

\begin{qenv}{74} Simplify $\log(a)+\log(b)-\log(c)$:
\begin{options} \item $\log\left(\dfrac{ab}{c}\right)$ \item $\log(a+b-c)$ \item $\log\left(\dfrac{ac}{b}\right)$ \item $\log\left(\dfrac{a+b}{b}\right)$ \end{options} \end{qenv}

\begin{qenv}{75} Expand $\log(x^2y^3)$:
\begin{options} \item $6\log$ \item $\log x + \log y$ \item $5\log$ \item $2\log x + 3\log y$ \end{options} \end{qenv}

\begin{qenv}{76} $\log_3\left(\dfrac{27}{9}\right) = \ldots\ldots$:
\begin{options} \item $\log_3 27 - \log_3 9$ \item $\log_3 27 + \log_3 9$ \item $\log_3(27\times9)$ \item $\log_3(27+3)$ \end{options} \end{qenv}

\begin{qenv}{77} Simplify $\log_5(2^4)$:
\begin{options} \item $4\log_5 2$ \item $2\log_5 4$ \item $\log_5 8$ \item $\log_5 16$ \end{options} \end{qenv}

\begin{qenv}{78} $\log_4 8 + \log_4 2 = $:
\begin{options} \item $\log_8 2$ \item $\log_4 10$ \item $\log_8 4$ \item $\log_4 16$ \end{options} \end{qenv}

\begin{qenv}{79} Change $\log_7 5$ to base 10:
\begin{options} \item $\log_7 10$ \item $\dfrac{\log_5 7}{\log_5 10}$ \item $\dfrac{\log_{10} 5}{\log_{10} 7}$ \item $\log_{10} 7$ \end{options} \end{qenv}

\begin{qenv}{80} Simplify $\log_6(6^5)$:
\begin{options} \item $5$ \item $6$ \item $1$ \item $0$ \end{options} \end{qenv}

\begin{qenv}{81} $\log_3 81 - \log_3 9 = $:
\begin{options} \item $\log_3 9$ \item $\log_3 72$ \item $\log_3(81/9)$ \item $\log_3(81\times9)$ \end{options} \end{qenv}

\begin{qenv}{82} If $\log_a b = c$, then:
\begin{options} \item $c^a=b$ \item $b^c=a$ \item $a^b=c$ \item $a^c=b$ \end{options} \end{qenv}

\begin{qenv}{83} Convert $\ln(x)=2$ to exponential form:
\begin{options} \item $e^2=x$ \item $10^2=x$ \item $2^{10}=x$ \item $x^2=e$ \end{options} \end{qenv}

\begin{qenv}{84} If $M=\log(A/A_0)$, $A=1000$ and $A_0=10$, find $M$:
\begin{options} \item $4$ \item $3$ \item $2$ \item $5$ \end{options} \end{qenv}

\begin{qenv}{85} The standard form of $5.2\times10^6$ is:
\begin{options} \item $52{,}000$ \item $520{,}000$ \item $5{,}200{,}000$ \item $52{,}000{,}000$ \end{options} \end{qenv}

\begin{qenv}{86} Scientific notation of $0.00034$ is:
\begin{options} \item $3.4\times10^3$ \item $3.4\times10^{-4}$ \item $3.4\times10^4$ \item $3.4\times10^{-3}$ \end{options} \end{qenv}

\begin{qenv}{87} The base of the common logarithm is:
\begin{options} \item $2$ \item $10$ \item $5$ \item $e$ \end{options} \end{qenv}

\begin{qenv}{88} $\log_2(2^3)$:
\begin{options} \item $1$ \item $2$ \item $5$ \item $3$ \end{options} \end{qenv}

\begin{qenv}{89} $\log 100 = \ldots\ldots$:
\begin{options} \item $2$ \item $3$ \item $10$ \item $1$ \end{options} \end{qenv}

\begin{qenv}{90} If $\log 2 = 0.3010$, then $\log 200$ is:
\begin{options} \item $1.3010$ \item $0.6010$ \item $2.3010$ \item $2.6010$ \end{options} \end{qenv}

\begin{qenv}{91} $\log(0) = \ldots\ldots$:
\begin{options} \item Positive \item Negative \item Zero \item Undefined \end{options} \end{qenv}

\begin{qenv}{92} $\log 10{,}000 = \ldots\ldots$:
\begin{options} \item $2$ \item $3$ \item $4$ \item $5$ \end{options} \end{qenv}

\begin{qenv}{93} $\log 5 + \log 3 = \ldots\ldots$:
\begin{options} \item $\log 0$ \item $\log 2$ \item $\log\left(\dfrac{5}{3}\right)$ \item $\log 15$ \end{options} \end{qenv}

\begin{qenv}{94} $3^4=81$ in logarithmic form is:
\begin{options} \item $\log_3 4 = 81$ \item $\log_4 3 = 81$ \item $\log_3 81 = 4$ \item $\log_4 81 = 3$ \end{options} \end{qenv}

\begin{qenv}{95} The sequence $0,1,1,2,3,5,8,13,21,\ldots$ is an example of a/an \underline{\hspace{2cm}} sequence:
\begin{options} \item Roman \item Odd \item Even \item Fibonacci \end{options} \end{qenv}

\begin{qenv}{96} Georg Cantor was a/an \underline{\hspace{2cm}} mathematician:
\begin{options} \item Japanese \item German \item Greek \item Egyptian \end{options} \end{qenv}

\begin{qenv}{97} The Fibonacci sequence formula $F_n = \ldots\ldots$:
\begin{options} \item $F_{n-1}+F_{n+2}$ \item $F_{n-1}+F_{n-2}$ \item $F_{n-1}-F_{n+2}$ \item $F_{n-1}+F_{n}$ \end{options} \end{qenv}

\begin{qenv}{98} The names of sets are represented as:
\begin{options} \item $1,2,3,\ldots$ \item $A,B,C,\ldots$ \item Both (A) and (B) \item None of these \end{options} \end{qenv}




\begin{qenv}{99} How many methods are there to describe a set:
\begin{options} \item $1$ \item $2$ \item $3$ \item $4$ \end{options} \end{qenv}

\begin{qenv}{100} Set $A=\{a,e,i,o,u\}$ is written in \underline{\hspace{2cm}} form:
\begin{options} \item Descriptive \item Tabular \item Builder \item None of these \end{options} \end{qenv}

\begin{qenv}{101} Symbol used for ``member of'' or ``belongs to'' a set:
\begin{options} \item $\subseteq$ \item $\in$ \item $\notin$ \item $\forall$ \end{options} \end{qenv}

\begin{qenv}{102} Set $\{2,3,5,7,11,\ldots\}$ is called a \underline{\hspace{2cm}} set:
\begin{options} \item Natural \item Whole \item Odd \item Prime \end{options} \end{qenv}

\begin{qenv}{103} $Q \cup Q' = \ldots\ldots$:
\begin{options} \item $N$ \item $R$ \item $Q$ \item $Q'$ \end{options} \end{qenv}

\begin{qenv}{104} The set $\{0\}$ is:
\begin{options} \item Singleton \item Empty \item Null \item Prime \end{options} \end{qenv}

\begin{qenv}{105} The set with no element is called an \underline{\hspace{2cm}} set:
\begin{options} \item Singleton set \item Empty \item Power \item Prime \end{options} \end{qenv}

\begin{qenv}{106} Symbol used for the empty set:
\begin{options} \item $\{\phi\}$ \item $\phi$ \item $\{\{\ \}\}$ \item None of these \end{options} \end{qenv}

\begin{qenv}{107} Example of an equal set:
\begin{options} \item $\{a,b,c\},\{1,2,3\}$ \item $\{1,2,3\},\{2,1,3\}$ \item Both (A) and (B) \item None of these \end{options} \end{qenv}

\begin{qenv}{108} Symbol used for equivalent sets $A$ and $B$:
\begin{options} \item $A=B$ \item $A\equiv B$ \item $A+B$ \item $A\cup B$ \end{options} \end{qenv}

\begin{qenv}{109} If $A$ is a subset of $B$, it is written as:
\begin{options} \item $A\subseteq B$ \item $A\supseteq B$ \item $A\in B$ \item $A\notin B$ \end{options} \end{qenv}

\begin{qenv}{110} Number of elements of the power set of $C=\{a,b,c,d\}$:
\begin{options} \item $8$ \item $16$ \item $24$ \item $26$ \end{options} \end{qenv}

\begin{qenv}{111} If $D=\{a\}$ then $P(D) = \ldots\ldots$:
\begin{options} \item $\{a\}$ \item $\phi,\{a\}$ \item $\{\phi,\{a\}\}$ \item None of these \end{options} \end{qenv}

\begin{qenv}{112} $S_1=$ set of odd numbers, $S_2=$ set of even numbers; the intersection of both sets is:
\begin{options} \item Universal set \item Empty set \item Even set \item Odd set \end{options} \end{qenv}

\begin{qenv}{113} $\{x \mid x\in A \lor x\in B\}$ is represented as:
\begin{options} \item $A\cup B$ \item $A\cap B$ \item $A-B$ \item $\phi$ \end{options} \end{qenv}

\begin{qenv}{114} Symbol used for intersection:
\begin{options} \item $\cup$ \item $I$ \item $\cap$ \item $-$ \end{options} \end{qenv}

\begin{qenv}{115} $A\cap B = \ldots\ldots$:
\begin{options} \item $\{x \mid x\in A \lor x\in B\}$ \item $\{x \mid x\in A \land x\in B\}$ \item $\{x \mid x\in A \land x\notin B\}$ \item $\{x \mid x\notin A \land x\in B\}$ \end{options} \end{qenv}

\begin{qenv}{116} $A\cap B=\{1,2,3\}$; if $A=\{1,2,3,7\}$ then $B = \ldots\ldots$:
\begin{options} \item $\{1,2,6\}$ \item $\{1,2,3,4\}$ \item $\{1,3,6\}$ \item $\{1,6,7\}$ \end{options} \end{qenv}

\begin{qenv}{117} If $A\cap B=\phi$, then $A$ and $B$ are called:
\begin{options} \item Subset \item Overlapping \item Disjoint \item None of these \end{options} \end{qenv}

\begin{qenv}{118} Which is an example of overlapping sets:
\begin{options} \item $\{1,2\},\{2,3,4\}$ \item $\{1,2\},\{1,2,3,4\}$ \item $\{1,2\},\{3,4\}$ \item None of these \end{options} \end{qenv}

\begin{qenv}{119} $\{x \mid x\in A \lor x\notin B\}$ is shown as:
\begin{options} \item $A-B$ \item $B-A$ \item $A\cup B$ \item $A+B$ \end{options} \end{qenv}

\begin{qenv}{120} If $A=\{1,2,3,4,5\}$ and $B=\{4,5,6,7,8,9,10\}$, then $A-B = \ldots\ldots$:
\begin{options} \item $\{4,5\}$ \item $\{1,2,3\}$ \item $\{7,8,9,10\}$ \item $\{4,5,6,7,8,9,10\}$ \end{options} \end{qenv}

\begin{qenv}{121} $A'$ is equal to:
\begin{options} \item $U-A$ \item $A-U$ \item $U$ \item $U+A$ \end{options} \end{qenv}

\begin{qenv}{122} $A' = \ldots\ldots$:
\begin{options} \item $U-A$ \item $A-U$ \item $U$ \item $\phi$ \end{options} \end{qenv}

\begin{qenv}{123} $(A\cup B)' = \ldots\ldots$:
\begin{options} \item $A\cap B$ \item $A'\cup B'$ \item $A'\cap B'$ \item $A\cup B$ \end{options} \end{qenv}

\begin{qenv}{124} $A\cap B=B\cap A$ is called the \underline{\hspace{2cm}} property of intersection:
\begin{options} \item Associative \item Commutative \item Distributive \item None of these \end{options} \end{qenv}

\begin{qenv}{125} If $A=\{1,3,5,7,9,11\}$, then $n(A) = \ldots\ldots$:
\begin{options} \item $6$ \item $4$ \item $5$ \item $2$ \end{options} \end{qenv}

\begin{qenv}{126} $n(A)+n(B)-n(A\cap B) = \ldots\ldots$:
\begin{options} \item $n(A\cap B)$ \item $n(A\cup B)$ \item $n$ \item None of these \end{options} \end{qenv}

\begin{qenv}{127} $A\times B$ is written symbolically as:
\begin{options} \item $\{(x,y)\mid x\in A \text{ and } x\in B\}$ \item $\{(x,y)\mid x\in A \text{ and } y\in B\}$ \item $\{(x,y)\mid x\notin A \text{ and } x\in B\}$ \item $\{(x,y)\mid x\in B \text{ and } y\in A\}$ \end{options} \end{qenv}

\begin{qenv}{128} Let $A=\{1,2,3\}$; determine the relation $r=\{(x,y)\mid x<y\}$:
\begin{options} \item $\{(2,1),(3,1),(3,2)\}$ \item $\{(1,2),(1,3),(2,3)\}$ \item $\{(1,1),(2,2),(3,3)\}$ \item None of these \end{options} \end{qenv}

\begin{qenv}{129} If $r=\{(1,2),(1,3),(2,3)\}$, then the domain of $r$ is:
\begin{options} \item $\{1,2\}$ \item $\{1\}$ \item $\{2,3\}$ \item $\{1,2,3\}$ \end{options} \end{qenv}

\begin{qenv}{130} Each ordered pair consists of \underline{\hspace{2cm}} coordinates:
\begin{options} \item $1$ \item $2$ \item $3$ \item $4$ \end{options} \end{qenv}

\begin{qenv}{131} In the ordered pair $(x,y)$, $x$ is called the:
\begin{options} \item Abscissa \item Ordinate \item Both \item None of these \end{options} \end{qenv}

\begin{qenv}{132} If a function $f:A\to B$ is such that range $f \neq B$, then $f$ is called an \underline{\hspace{2cm}} function:
\begin{options} \item Onto \item Into \item One-one \item Bijective \end{options} \end{qenv}

\begin{qenv}{133} If $f(n)=2n-1$, then $f(7)$ is:
\begin{options} \item $12$ \item $13$ \item $14$ \item $15$ \end{options} \end{qenv}

\begin{qenv}{134} If $g(n)=n^2-3$, then $g(4)$ is:
\begin{options} \item $7$ \item $-7$ \item $13$ \item $9$ \end{options} \end{qenv}

\begin{qenv}{135} The set-builder form of the set $\left\{1,\dfrac{1}{3},\dfrac{1}{5},\dfrac{1}{7},\ldots\right\}$:
\begin{options} \item $\left\{x \mid x=\dfrac{1}{n}, n\in W\right\}$ \item $\left\{x \mid x=\dfrac{1}{2n+1}, n\in W\right\}$ \item $\left\{x \mid x=\dfrac{1}{n+1}, n\in W\right\}$ \item $\{x \mid x=2n+1, n\in w\}$ \end{options} \end{qenv}

\begin{qenv}{136} If $A=\{\ \}$, then $P(A)$ is:
\begin{options} \item $\{\ \}$ \item $\{1\}$ \item $\{\{\ \}\}$ \item $\phi$ \end{options} \end{qenv}

\begin{qenv}{137} If $U=\{1,2,3,4,5\}$, $A=\{1,2,3\}$ and $B=\{3,4,5\}$, then $U-(A\cap B)$ is:
\begin{options} \item $\{1,2,4,5\}$ \item $\{2,3\}$ \item $\{1,3,4,5\}$ \item $\{1,2,3\}$ \end{options} \end{qenv}

\begin{qenv}{138} If $A$ and $B$ are overlapping sets, then $n(A-B)$ equals:
\begin{options} \item $n(A)$ \item $n(B)$ \item $A\cap B$ \item $n(A)-n(A\cap B)$ \end{options} \end{qenv}

\begin{qenv}{139} If $A\subseteq B$ and $B-A \neq \phi$, then $n(B-A)$ is equal to:
\begin{options} \item $0$ \item $-n(B)$ \item $n(A)$ \item $n(B)-n(A)$ \end{options} \end{qenv}

\begin{qenv}{140} If $n(A\cup B)=50$, $n(A)=30$ and $n(B)=35$, then $n(A\cap B)=$:
\begin{options} \item $23$ \item $15$ \item $9$ \item $40$ \end{options} \end{qenv}

\begin{qenv}{141} If $A=\{1,2,3,4\}$ and $B=\{x,y,z\}$, then the Cartesian product of $A$ and $B$ contains exactly \underline{\hspace{2cm}} elements:
\begin{options} \item $13$ \item $12$ \item $10$ \item $6$ \end{options} \end{qenv}

\begin{qenv}{142} If $f(x)=x^2-3x+2$, then the value of $f(a+1)$ is equal to:
\begin{options} \item $a+1$ \item $a^2+1$ \item $a^2+2a+1$ \item $a^2-a$ \end{options} \end{qenv}

\begin{qenv}{143} Given that $f(x)=3x+1$, if $f(x)=28$, then the value of $x$ is:
\begin{options} \item $9$ \item $27$ \item $3$ \item $18$ \end{options} \end{qenv}

\begin{qenv}{144} Let $A=\{1,2,3\}$ and $B=\{a,b\}$ be two non-empty sets, and $f:A\to B$ be a function defined as $f=\{(1,a),(2,b),(3,b)\}$; which of the following statements is true?
\begin{options} \item $f$ is injective \item $f$ is surjective \item $f$ is bijective \item $f$ is into only \end{options} \end{qenv}

\begin{qenv}{145} Common factor of $2x-6$ is:
\begin{options} \item $2$ \item $3$ \item $x$ \item None of these \end{options} \end{qenv}

\begin{qenv}{146} Factors of $x^2-3x-10$ are:
\begin{options} \item $(x+2)(x+5)$ \item $(x+2)(x-5)$ \item $(x-2)(x+5)$ \item $(x-2)(x-5)$ \end{options} \end{qenv}

\begin{qenv}{147} Which expression is of the type of a quadratic expression:
\begin{options} \item $ax+b$ \item $ax^2+bx+c$ \item $ax^3+bx^2+c$ \item None of these \end{options} \end{qenv}

\begin{qenv}{148} Which pair is suitable for the factorization of $x^2-11x+24$:
\begin{options} \item $+3,+8$ \item $-3,-8$ \item $-3,+8$ \item $+3,-8$ \end{options} \end{qenv}

\begin{qenv}{149} An expression having degree $2$ is called a \underline{\hspace{2cm}} expression:
\begin{options} \item Cubic \item Linear \item Quadratic \item None of these \end{options} \end{qenv}

\begin{qenv}{150} Factors of $3x^2-4x-4$ are:
\begin{options} \item $(3x+2)(x-2)$ \item $(3x-2)(x+2)$ \item $(3x+2)(x+2)$ \item None of these \end{options} \end{qenv}

\begin{qenv}{151} Which of the following is a prime number:
\begin{options} \item $9$ \item $15$ \item $13$ \item $21$ \end{options} \end{qenv}

\begin{qenv}{152} $a^2-2ab+b^2$ is equal to:
\begin{options} \item $(a+b)^2$ \item $(a-b)^2$ \item $a^2-b^2$ \item $a^2+b^2$ \end{options} \end{qenv}

\begin{qenv}{153} Factors of $8x^3+27$:
\begin{options} \item $(2x+3)(4x^2+6x+9)$ \item $(2x+3)(4x^2-6x+9)$ \item $(2x-3)(4x^2+6x+9)$ \item $(2x-3)(4x^2-6x+9)$ \end{options} \end{qenv}

\begin{qenv}{154} In $2x^2+7x+3$, one factor is $2x+1$, then the other factor is:
\begin{options} \item $(x+2)$ \item $(x+1)$ \item $(x+3)$ \item $(x+4)$ \end{options} \end{qenv}

\begin{qenv}{155} An expression with two terms is called:
\begin{options} \item Monomial \item Binomial \item Trinomial \item None of these \end{options} \end{qenv}

\begin{qenv}{156} $a^3+b^3+3ab(a+b)$ is equal to:
\begin{options} \item $(a+b)^3$ \item $(a-b)^3$ \item $a^3-b^3$ \item $a^3+b^3$ \end{options} \end{qenv}

\begin{qenv}{157} How many methods are there to solve for the H.C.F:
\begin{options} \item $1$ \item $2$ \item $3$ \item $4$ \end{options} \end{qenv}

\begin{qenv}{158} H.C.F of $6x^2y$, $9xy^2$:
\begin{options} \item $3x^2y$ \item $3xy$ \item $6x^2y$ \item $54x^2y^3$ \end{options} \end{qenv}

\begin{qenv}{159} L.C.M $= \ldots\ldots\ldots\ldots\ldots$:
\begin{options} \item $\dfrac{p(x)\times q(x)}{H.C.F}$ \item $\dfrac{H.C.F \times p(x)}{q(x)}$ \item $P$ \item None of these \end{options} \end{qenv}

\begin{qenv}{160} Square root of $a^2-2a+1$:
\begin{options} \item $\pm(a-1)$ \item $\pm(a+1)$ \item $a-1$ \item $a+1$ \end{options} \end{qenv}

\begin{qenv}{161} The square root of $x^2-8x+16$ is:
\begin{options} \item $(x+4)$ \item $-(x+4)$ \item Both (A) and (B) \item None of these \end{options} \end{qenv}

\begin{qenv}{162} The return on investment $R(x)$ is given by $R(x)=x^2+6x-8$. What are the investment levels that result in zero return:
\begin{options} \item $x=2, x=4$ \item $x=1, x=8$ \item $x=2, x=-4$ \item $x=-2, x=4$ \end{options} \end{qenv}

\begin{qenv}{163} The factorization of $12x+36$ is:
\begin{options} \item $12(x+3)$ \item $12(3x)$ \item $12(3x+1)$ \item $x(12+36x)$ \end{options} \end{qenv}

\begin{qenv}{164} The factors of $4x^2-12x+9$:
\begin{options} \item $(2x+3)^2$ \item $(2x-3)^2$ \item $(2x-3)(2x+3)$ \item $(2+3x)(2-3x)^2$ \end{options} \end{qenv}

\begin{qenv}{165} The HCF of $a^2b^2$ and $a^3b^3$:
\begin{options} \item $a^3b^3$ \item $ab^2$ \item $a^4b^5$ \item $a^2b$ \end{options} \end{qenv}

\begin{qenv}{166} The LCM of $16x^2$, $4x$ and $30xy$ is:
\begin{options} \item $480x^3y$ \item $240xy$ \item $240x^2y$ \item $120x^4y$ \end{options} \end{qenv}

\begin{qenv}{167} Product of LCM and HCF $= \ldots\ldots$ of two polynomials:
\begin{options} \item Sum \item Difference \item Product \item Quotient \end{options} \end{qenv}

\begin{qenv}{168} The square root of $x^2-6x+9$:
\begin{options} \item $\pm(x-3)$ \item $\pm(x+3)$ \item $x-3$ \item $x+3$ \end{options} \end{qenv}

\begin{qenv}{169} The LCM of $(a-b)^2$ and $(a-b)^4$ is:
\begin{options} \item $(a-b)^2$ \item $(a-b)^3$ \item $(a-b)^4$ \item $(a-b)^6$ \end{options} \end{qenv}

\begin{qenv}{170} Factorization of $x^3+3x^2+3x+1$:
\begin{options} \item $(x+1)^3$ \item $(x-1)^3$ \item $(x+1)(x^2+x+1)$ \item $(x-1)(x^2-x+1)$ \end{options} \end{qenv}

\begin{qenv}{171} A cubic polynomial has degree:
\begin{options} \item $1$ \item $2$ \item $3$ \item $4$ \end{options} \end{qenv}

\begin{qenv}{172} One of the factors of $x^3-27$:
\begin{options} \item $x-3$ \item $x+3$ \item $x^2-3x+9$ \item Both (A) and (C) \end{options} \end{qenv}

\begin{qenv}{173} Which of the following is a linear equation in one variable:
\begin{options} \item $2x+3y=7$ \item $5x-4=0$ \item $x^2+3x=8$ \item $xy+2=5$ \end{options} \end{qenv}

\begin{qenv}{174} The solution of the equation $3x-7=2$ is:
\begin{options} \item $3$ \item $-3$ \item $\dfrac{5}{3}$ \item $9$ \end{options} \end{qenv}

\begin{qenv}{175} In the inequality $4x+2>10$, the solution for $x$ is:
\begin{options} \item $x>2$ \item $x<2$ \item $x>1$ \item $x<2$ \end{options} \end{qenv}

\begin{qenv}{176} Which symbol represents `less than or equal to':
\begin{options} \item $>$ \item $<$ \item $\leq$ \item $\geq$ \end{options} \end{qenv}

\begin{qenv}{177} The graphical representation of a linear equation in two variables is:
\begin{options} \item A curve \item A straight line \item A parabola \item A circle \end{options} \end{qenv}

\begin{qenv}{178} What is the solution of $2(x-3)=4$:
\begin{options} \item $x=2$ \item $x=5$ \item $x=3$ \item $x=1$ \end{options} \end{qenv}

\begin{qenv}{179} The inequality $x-5 \geq 3$ means:
\begin{options} \item $x$ is greater than $3$ \item $x$ is greater than or equal to $3$ \item $x$ is greater than or equal to $8$ \item $x$ is less than $3$ \end{options} \end{qenv}

\begin{qenv}{180} The equation $x+2y=6$ represents a:
\begin{options} \item Quadratic equation \item Linear equation \item Exponential equation \item None of these \end{options} \end{qenv}

\begin{qenv}{181} The associated equation of the inequality $x+2y \leq 8$:
\begin{options} \item $x+2y=8$ \item $x+2y>8$ \item $x+2y<8$ \item $x+2y \geq 8$ \end{options} \end{qenv}

\begin{qenv}{182} The solution of the system $x-y=2$ and $x+y=6$ is:
\begin{options} \item $(4,2)$ \item $(2,4)$ \item $(3,3)$ \item $(6,0)$ \end{options} \end{qenv}

\begin{qenv}{183} If $x=3$ is a solution of $5x-2=k$, then $k$ is:
\begin{options} \item $11$ \item $13$ \item $15$ \item $17$ \end{options} \end{qenv}

\begin{qenv}{184} What is the graphical representation of a system of linear inequalities:
\begin{options} \item A shaded region \item A single point \item A straight line \item A curve \end{options} \end{qenv}

\begin{qenv}{185} The equation $2x-y=4$ is equivalent to the inequality:
\begin{options} \item $2x-y>4$ \item $2x-y\leq 4$ \item $2x-y\geq 4$ \item None of these \end{options} \end{qenv}

\begin{qenv}{186} Which of the following is an objective function:
\begin{options} \item $x+y$ \item $x^2+y^2$ \item $3x+2y$ \item Both (A) and (C) \end{options} \end{qenv}

\begin{qenv}{187} The feasible region for an inequality system is always:
\begin{options} \item A single point \item A region on the plane \item An empty set \item A line segment \end{options} \end{qenv}

\begin{qenv}{188} A vertical line divides the plane into:
\begin{options} \item Two half-planes \item Four quadrants \item Three sections \item Infinite parts \end{options} \end{qenv}

\begin{qenv}{189} If a system of inequalities has no solution, its feasible region is:
\begin{options} \item Infinite \item A single point \item A line segment \item Empty \end{options} \end{qenv}

\begin{qenv}{190} The intersection of two half-planes is known as:
\begin{options} \item A vertex \item A feasible region \item A line \item A parabola \end{options} \end{qenv}

\begin{qenv}{191} The graph of $x+y \geq 5$ includes the line:
\begin{options} \item $x+y=5$ \item $x+y<5$ \item $x+y>5$ \item None of these \end{options} \end{qenv}

\begin{qenv}{192} Which method is used to find the optimal solution of a linear programming problem:
\begin{options} \item Graphical method \item Algebraic method \item Simplex method \item All of these \end{options} \end{qenv}

\begin{qenv}{193} The equation $x=2$ represents a:
\begin{options} \item Horizontal line \item Vertical line \item Diagonal line \item None of these \end{options} \end{qenv}

\begin{qenv}{194} A point that satisfies all given inequalities is called:
\begin{options} \item Feasible solution \item Optimal solution \item Boundary point \item None of these \end{options} \end{qenv}

\begin{qenv}{195} The line $y=3x+2$ has a slope of:
\begin{options} \item $2$ \item $3$ \item $-3$ \item $-2$ \end{options} \end{qenv}

\begin{qenv}{196} If the inequalities $x\geq 0$ and $y\geq 0$ are given, the solution region lies in the:
\begin{options} \item First quadrant \item Second quadrant \item Third quadrant \item Fourth quadrant \end{options} \end{qenv}

\begin{qenv}{197} The inequality $5x-3y<15$ represents:
\begin{options} \item A parabola \item A straight line \item A shaded region \item A single point \end{options} \end{qenv}

\begin{qenv}{198} The maximum or minimum value of a function occurs at:
\begin{options} \item A random point \item A feasible point \item A corner point \item An undefined point \end{options} \end{qenv}

\begin{qenv}{199} If $x-y \leq 4$ then the region of solution is:
\begin{options} \item Above the line \item Below the line \item On the line \item None of these \end{options} \end{qenv}

\begin{qenv}{200} The intersection of $x+y=4$ and $x-y=2$ is:
\begin{options} \item $(1,3)$ \item $(3,1)$ \item $(2,2)$ \item $(4,0)$ \end{options} \end{qenv}

\begin{qenv}{201} In the following, the linear equation is:
\begin{options} \item $5x>7$ \item $4x-2<1$ \item $2x+1=1$ \item $3+1=4$ \end{options} \end{qenv}

\begin{qenv}{202} Solution of $5x-10=10$ is:
\begin{options} \item $0$ \item $50$ \item $4$ \item $-4$ \end{options} \end{qenv}

\begin{qenv}{203} If $7x+4<6x+6$, then $x$ belongs to the interval:
\begin{options} \item $(2,\infty)$ \item $[2,\infty)$ \item $(-\infty,2)$ \item $(-\infty,2]$ \end{options} \end{qenv}



\begin{enumerate}[label=\textbf{\arabic*.}, start=204, leftmargin=2em]

\item A vertical line divides the plane into:
\begin{enumerate}[label=(\Alph*)] \item Left half plane \item Right half plane \item Full plane \item Two half planes \end{enumerate}

\item The linear equation formed out of the linear inequality is called:
\begin{enumerate}[label=(\Alph*)] \item Linear equation \item Associated equation \item Quadratic equal \item None of these \end{enumerate}

\item $3x + 4 < 0$ is:
\begin{enumerate}[label=(\Alph*)] \item Equation \item Inequality \item Not inequality \item Identity \end{enumerate}

\item Corner point is also called:
\begin{enumerate}[label=(\Alph*)] \item Code \item Vertex \item Curve \item Region \end{enumerate}

\item $(0,0)$ is solution of inequality:
\begin{enumerate}[label=(\Alph*)] \item $4x + 5y > 8$ \item $3x + y > 6$ \item $-2x + 3y < 0$ \item $x + y > 4$ \end{enumerate}

\item The solution region restricted to the first quadrant is called:
\begin{enumerate}[label=(\Alph*)] \item Objective region \item Feasible region \item Solution region \item Constraints region \end{enumerate}

\item A function that is to be maximized or minimized is called:
\begin{enumerate}[label=(\Alph*)] \item Solution function \item Objective function \item Feasible function \item None of these \end{enumerate}

\item In $\angle AOB$ vertex is:
\begin{enumerate}[label=(\Alph*)] \item $A$ \item $O$ \item $B$ \item None of these \end{enumerate}

\item Ray $OA$ written as symbolically:
\begin{enumerate}[label=(\Alph*)] \item $\overrightarrow{OA}$ \item $\overleftrightarrow{OA}$ \item $\overline{OA}$ \item $\overleftarrow{OA}$ \end{enumerate}

\item Co-terminal of $65^\circ$:
\begin{enumerate}[label=(\Alph*)] \item $360^\circ$ \item $-360^\circ$ \item $425^\circ$ \item $-405^\circ$ \end{enumerate}

\item One full revolution of a circle is \rule{1cm}{0.4pt} radian:
\begin{enumerate}[label=(\Alph*)] \item $\pi$ \item $2\pi$ \item $3\pi$ \item $4\pi$ \end{enumerate}

\item Area of a sector is $=$ \rule{1cm}{0.4pt}:
\begin{enumerate}[label=(\Alph*)] \item $\dfrac{1}{2}r^2\theta$ \item $\pi r^2$ \item $4\pi r^2$ \item $\dfrac{1}{2}r^3\theta$ \end{enumerate}

\item Value of $\dfrac{\pi}{4}$ radian:
\begin{enumerate}[label=(\Alph*)] \item $\sqrt{2}$ \item $\sqrt{3}$ \item $\dfrac{1}{\sqrt{2}}$ \item $\dfrac{1}{\sqrt{3}}$ \end{enumerate}

\item What is the measure of a full rotation in degrees:
\begin{enumerate}[label=(\Alph*)] \item $90^\circ$ \item $180^\circ$ \item $270^\circ$ \item $360^\circ$ \end{enumerate}

\item What is the measure of a right angle in radians?
\begin{enumerate}[label=(\Alph*)] \item $\dfrac{\pi}{2}$ \item $\pi$ \item $\dfrac{3\pi}{2}$ \item $2\pi$ \end{enumerate}

\item Which of the following is an obtuse angle:
\begin{enumerate}[label=(\Alph*)] \item $45^\circ$ \item $90^\circ$ \item $120^\circ$ \item $180^\circ$ \end{enumerate}

\item What is the measure of a straight angle in radians?
\begin{enumerate}[label=(\Alph*)] \item $2\pi$ \item $\dfrac{3\pi}{2}$ \item $\pi$ \item $\dfrac{\pi}{2}$ \end{enumerate}

\item What is the measure of 1 radian in degrees:
\begin{enumerate}[label=(\Alph*)] \item $57.2958^\circ$ \item $90^\circ$ \item $180^\circ$ \item $360^\circ$ \end{enumerate}

\item Convert $73.12^\circ$ to degrees, minutes, and seconds:
\begin{enumerate}[label=(\Alph*)] \item $73^\circ12'12''$ \item $73^\circ12'7''$ \item $73^\circ7'7''$ \item $73^\circ7'12''$ \end{enumerate}

\item Convert $45^\circ45'45''$ to decimal degrees:
\begin{enumerate}[label=(\Alph*)] \item $45.1234^\circ$ \item $45.4567^\circ$ \item $45.7625^\circ$ \item $45.9876^\circ$ \end{enumerate}

\item What is the value of $1^\circ$ in seconds:
\begin{enumerate}[label=(\Alph*)] \item $60''$ \item $360''$ \item $3600''$ \item $60'$ \end{enumerate}

\item What is the measure of a full circle in radians?
\begin{enumerate}[label=(\Alph*)] \item $2\pi$ \item $\dfrac{3\pi}{2}$ \item $\pi$ \item $\dfrac{\pi}{2}$ \end{enumerate}

\item Convert $\dfrac{5\pi}{3}$ radians to degrees.
\begin{enumerate}[label=(\Alph*)] \item $300^\circ$ \item $150^\circ$ \item $210^\circ$ \item $330^\circ$ \end{enumerate}

\item Convert $15^\circ$ to radians:
\begin{enumerate}[label=(\Alph*)] \item $\dfrac{\pi}{12}$ \item $\dfrac{\pi}{6}$ \item $\dfrac{\pi}{4}$ \item $\dfrac{\pi}{3}$ \end{enumerate}

\item Convert $1.2$ radians to degrees:
\begin{enumerate}[label=(\Alph*)] \item $90^\circ$ \item $57.2958^\circ$ \item $45^\circ$ \item $68.75^\circ$ \end{enumerate}

\item In right-angled triangle $ABC$, if $m\angle B = 90^\circ$ then which statement is true:
\begin{enumerate}[label=(\Alph*)] \item $a^2+b^2=c^2$ \item $a^2-b^2=c^2$ \item $a^2+b^2=b^2$ \item $a^2+b^2=a^2$ \end{enumerate}

\item Reciprocal of $\sin\theta$:
\begin{enumerate}[label=(\Alph*)] \item $\sec\theta$ \item $\csc\theta$ \item $\cot\theta$ \item $\tan\theta$ \end{enumerate}

\item $\cot\theta = $ \rule{1cm}{0.4pt}
\begin{enumerate}[label=(\Alph*)] \item $\dfrac{1}{\sin\theta}$ \item $\dfrac{1}{\cos\theta}$ \item $\dfrac{1}{\tan\theta}$ \item $1-\cos^2\theta$ \end{enumerate}

\item $\sin\theta = $ \rule{1cm}{0.4pt}
\begin{enumerate}[label=(\Alph*)] \item $\sqrt{1-\cos^2\theta}$ \item $\dfrac{1}{\tan\theta}$ \item $\dfrac{\sin\theta}{\cos\theta}$ \item $1-\cos^2\theta$ \end{enumerate}

\item If $\sin\theta = \dfrac{2}{3}$, then value of $\csc\theta$:
\begin{enumerate}[label=(\Alph*)] \item $\dfrac{2}{3}$ \item $-\dfrac{2}{3}$ \item $\dfrac{3}{2}$ \item $-\dfrac{3}{2}$ \end{enumerate}

\item Fundamental trigonometric ratios are:
\begin{enumerate}[label=(\Alph*)] \item 3 \item 4 \item 5 \item 6 \end{enumerate}

\item What is the value of $\cot 30^\circ$:
\begin{enumerate}[label=(\Alph*)] \item $\sqrt{3}$ \item $\dfrac{1}{\sqrt{3}}$ \item 1 \item 0 \end{enumerate}

\item What is the value of $\sin(90^\circ - \theta)$?
\begin{enumerate}[label=(\Alph*)] \item $\tan\theta$ \item $\sin\theta$ \item $\cos\theta$ \item $\cot\theta$ \end{enumerate}

\item What is the value of $\tan\theta$ in terms of $\sin\theta$ and $\cos\theta$?
\begin{enumerate}[label=(\Alph*)] \item $\dfrac{\sin\theta}{\cos\theta}$ \item $\dfrac{\cos\theta}{\sin\theta}$ \item $\dfrac{1}{\sin\theta}$ \item $\dfrac{1}{\cos\theta}$ \end{enumerate}

\item Which of the following is a Pythagorean identity:
\begin{enumerate}[label=(\Alph*)] \item $\sin^2\theta+\cos^2\theta=1$ \item $\sin\theta+\cos\theta=1$ \item $\tan\theta+\cot\theta=1$ \item $\sec\theta+\csc\theta=1$ \end{enumerate}

\item What is the value of $\tan^2\theta + 1$?
\begin{enumerate}[label=(\Alph*)] \item $\sin^2\theta$ \item $\csc^2\theta$ \item $\cot^2\theta$ \item $\sec^2\theta$ \end{enumerate}

\item What is the value of $1 + \cot^2\theta$:
\begin{enumerate}[label=(\Alph*)] \item $\csc^2\theta$ \item $\sec^2\theta$ \item $\tan^2\theta$ \item $\sin^2\theta$ \end{enumerate}

\item If $\tan\theta = \dfrac{3}{4}$, what is the value of $\sin\theta$?
\begin{enumerate}[label=(\Alph*)] \item $\dfrac{5}{3}$ \item $\dfrac{4}{5}$ \item $\dfrac{3}{5}$ \item $\dfrac{5}{4}$ \end{enumerate}

\item What is the value of $\sin 45^\circ$?
\begin{enumerate}[label=(\Alph*)] \item $\dfrac{1}{\sqrt{2}}$ \item $\dfrac{1}{2}$ \item $\dfrac{\sqrt{3}}{2}$ \item 1 \end{enumerate}

\item What is the value of $\cos 30^\circ$:
\begin{enumerate}[label=(\Alph*)] \item $\dfrac{1}{2}$ \item $\dfrac{1}{\sqrt{2}}$ \item $\dfrac{\sqrt{3}}{2}$ \item 1 \end{enumerate}

\item What is the value of $\tan 60^\circ$?
\begin{enumerate}[label=(\Alph*)] \item $\sqrt{3}$ \item $\dfrac{1}{\sqrt{3}}$ \item 1 \item $\dfrac{\sqrt{3}}{2}$ \end{enumerate}

\item What is the value of $\cot 45^\circ$:
\begin{enumerate}[label=(\Alph*)] \item 1 \item $\dfrac{\sqrt{3}}{2}$ \item $\dfrac{1}{\sqrt{3}}$ \item 0 \end{enumerate}

\item In a right-angled triangle, if one angle is $30^\circ$, what is the other acute angle:
\begin{enumerate}[label=(\Alph*)] \item $60^\circ$ \item $45^\circ$ \item $90^\circ$ \item $120^\circ$ \end{enumerate}

\item In a right-angled triangle, if one side is 3 cm and the other side is 4 cm, what is the length of the hypotenuse:
\begin{enumerate}[label=(\Alph*)] \item 5 cm \item 6 cm \item 7 cm \item 8 cm \end{enumerate}

\item In a right-angled triangle, if one angle is $45^\circ$, what is the length of the hypotenuse if one side is 5 cm?
\begin{enumerate}[label=(\Alph*)] \item $5\sqrt{2}$ cm \item 5 cm \item 10 cm \item $10\sqrt{2}$ cm \end{enumerate}

\item In a right-angled triangle, if the hypotenuse is 10 cm and one side is 6 cm, what is the length of the other side:
\begin{enumerate}[label=(\Alph*)] \item 4 cm \item 6 cm \item 8 cm \item 2 cm \end{enumerate}

\item In a right-angled triangle, if one side is 3 cm and the other side is 4 cm, what is the length of the hypotenuse:
\begin{enumerate}[label=(\Alph*)] \item 5 cm \item 6 cm \item 7 cm \item 8 cm \end{enumerate}

\item The value of $\tan^{-1} 2$:
\begin{enumerate}[label=(\Alph*)] \item $\dfrac{\pi}{2}$ \item $\dfrac{3\pi}{2}$ \item $0.4636\pi$ \item $0.4636$ \end{enumerate}

\item In a right triangle, the hypotenuse is 13 units and one of the angles is $\theta = 30^\circ$:
\begin{enumerate}[label=(\Alph*)] \item 6.5 units \item 7.5 units \item 6 units \item 5 units \end{enumerate}

\item A person standing 50 m away from a building sees the top of the building at an angle of elevation of $45^\circ$. Height of the building is:
\begin{enumerate}[label=(\Alph*)] \item 50 m \item 25 m \item 35 m \item 70 m \end{enumerate}

\item $\sec^2\theta - \tan^2\theta$:
\begin{enumerate}[label=(\Alph*)] \item $\sin^2\theta$ \item 1 \item $\cos^2\theta$ \item $\cot^2\theta$ \end{enumerate}

\item If $\sin\theta = \dfrac{3}{5}$ and $\theta$ is an acute angle, $\cos^2\theta = $ \rule{1cm}{0.4pt}
\begin{enumerate}[label=(\Alph*)] \item $\dfrac{7}{25}$ \item $\dfrac{24}{25}$ \item $\dfrac{16}{25}$ \item $\dfrac{4}{25}$ \end{enumerate}

\item $\dfrac{5\pi}{24}$ rad $= $ \rule{1cm}{0.4pt} degree
\begin{enumerate}[label=(\Alph*)] \item $30^\circ$ \item $37.5^\circ$ \item $45^\circ$ \item $52.5^\circ$ \end{enumerate}

\item $292.5^\circ = $ \rule{1cm}{0.4pt} rad:
\begin{enumerate}[label=(\Alph*)] \item $\dfrac{17\pi}{6}$ \item $\dfrac{17\pi}{4}$ \item $1.6\pi$ \item $1.625\pi$ \end{enumerate}

\item Which of the following is a valid identity:
\begin{enumerate}[label=(\Alph*)] \item $\cos\left(\dfrac{\pi}{2}-\theta\right) = \sin\theta$ \item $\cos\left(\dfrac{\pi}{2}-\theta\right) = \cos\theta$ \item $\cos\left(\dfrac{\pi}{2}-\theta\right) = \sec\theta$ \item $\cos\left(\dfrac{\pi}{2}-\theta\right) = \csc\theta$ \end{enumerate}

\item $\sin 60^\circ = $ \rule{1cm}{0.4pt}
\begin{enumerate}[label=(\Alph*)] \item 1 \item $\dfrac{1}{2}$ \item $\sqrt{(3)^2}$ \item $\dfrac{\sqrt{3}}{2}$ \end{enumerate}

\item $\cos^2 100\pi + \sin^2 100\pi$:
\begin{enumerate}[label=(\Alph*)] \item 1 \item 2 \item 3 \item 4 \end{enumerate}

\item In a plane, horizontal line represents:
\begin{enumerate}[label=(\Alph*)] \item $xox'$ \item $yoy'$ \item $zoz'$ \item None of these \end{enumerate}

\item In a plane $yoy'$ is called \rule{1cm}{0.4pt}:
\begin{enumerate}[label=(\Alph*)] \item Horizontal line \item Vertical line \item Parallel line \item None of these \end{enumerate}

\item The points lying on $Ox$ is:
\begin{enumerate}[label=(\Alph*)] \item $+$ve \item $-$ve \item $+$ve or $-$ve \item None of these \end{enumerate}

\item In $(x, y)$ if $x > 0$ and $y < 0$ then the point lies in which quadrant:
\begin{enumerate}[label=(\Alph*)] \item I \item II \item III \item IV \end{enumerate}

\item The point $P(3, -4)$ lies in which quadrant:
\begin{enumerate}[label=(\Alph*)] \item I \item II \item III \item IV \end{enumerate}

\item The distance between the points $A(2, 3)$ and $B(5, 7)$ is:
\begin{enumerate}[label=(\Alph*)] \item 5 \item 6 \item 7 \item 8 \end{enumerate}

\item The midpoint of the line segment joining $(1, 2)$ and $(5, 8)$ is:
\begin{enumerate}[label=(\Alph*)] \item $(3, 5)$ \item $(3, 4)$ \item $(4, 5)$ \item $(4, 6)$ \end{enumerate}

\item If the distance between $(x, 3)$ and $(4, 7)$ is 5, then $x$ is:
\begin{enumerate}[label=(\Alph*)] \item 1 \item 2 \item 3 \item 4 \end{enumerate}

\item The point $(0, 5)$ lies on which axis:
\begin{enumerate}[label=(\Alph*)] \item $x$-axis \item $y$-axis \item Both axes \item None \end{enumerate}

\item The midpoint of a bridge connecting $(2, 5)$ and $(8, 9)$ is:
\begin{enumerate}[label=(\Alph*)] \item $(5, 7)$ \item $(6, 7)$ \item $(5, 8)$ \item $(6, 8)$ \end{enumerate}

\item The slope of the line passing through $(2, 3)$ and $(4, 7)$ is:
\begin{enumerate}[label=(\Alph*)] \item 1 \item 2 \item 3 \item 4 \end{enumerate}

\item The inclination of a line with slope 1 is:
\begin{enumerate}[label=(\Alph*)] \item $30^\circ$ \item $45^\circ$ \item $60^\circ$ \item $90^\circ$ \end{enumerate}

\item If the slope of a line is undefined, the line is:
\begin{enumerate}[label=(\Alph*)] \item Horizontal \item Vertical \item Parallel to $x$-axis \item Parallel to $y$-axis \end{enumerate}

\item The slope of a line parallel to the $x$-axis is:
\begin{enumerate}[label=(\Alph*)] \item 0 \item 1 \item Undefined \item $-1$ \end{enumerate}

\item The slope of a line perpendicular to a line with slope 2 is:
\begin{enumerate}[label=(\Alph*)] \item 2 \item $-2$ \item $\dfrac{1}{2}$ \item $-\dfrac{1}{2}$ \end{enumerate}

\item The equation of a line with slope 3 and $y$-intercept 2 is:
\begin{enumerate}[label=(\Alph*)] \item $y = 3x + 2$ \item $y = 2x + 3$ \item $y = 3x - 2$ \item $y = 2x - 3$ \end{enumerate}

\item The equation of a line passing through $(1, 2)$ with slope 4 is:
\begin{enumerate}[label=(\Alph*)] \item $y = 4x - 2$ \item $y = 4x + 2$ \item $y = 4x - 6$ \item $y = 4x + 6$ \end{enumerate}

\item The equation of a line passing through $(2, 3)$ and $(4, 7)$ is:
\begin{enumerate}[label=(\Alph*)] \item $y = 2x - 1$ \item $y = 2x + 1$ \item $y = 3x - 1$ \item $y = 3x + 1$ \end{enumerate}

\item The equation of a line with $x$-intercept 3 and $y$-intercept 4 is:
\begin{enumerate}[label=(\Alph*)] \item $4x+3y=12$ \item $3x+4y=12$ \item $4x-3y=12$ \item $3x-4y=12$ \end{enumerate}

\item The equation of a line parallel to $y = 2x + 3$ and passing through $(1, 1)$ is:
\begin{enumerate}[label=(\Alph*)] \item $y = 2x - 1$ \item $y = 2x + 1$ \item $y = 2x - 3$ \item $y = 2x + 3$ \end{enumerate}

\item The equation of a horizontal line passing through $(3, 4)$ is:
\begin{enumerate}[label=(\Alph*)] \item $y = 3$ \item $y = 4$ \item $x = 3$ \item $x = 4$ \end{enumerate}

\item The equation of a vertical line passing through $(5, 6)$ is:
\begin{enumerate}[label=(\Alph*)] \item $y = 5$ \item $y = 6$ \item $x = 5$ \item $x = 6$ \end{enumerate}

\item The slope of the line $2x + 3y = 6$ is:
\begin{enumerate}[label=(\Alph*)] \item $\dfrac{2}{3}$ \item $-\dfrac{2}{3}$ \item $\dfrac{3}{2}$ \item $-\dfrac{3}{2}$ \end{enumerate}

\end{enumerate}




\begin{qenv}{284} If two lines are perpendicular, the product of their slopes is:
\begin{options} \item $1$ \item $-1$ \item $0$ \item $2$ \end{options} \end{qenv}

\begin{qenv}{285} The slope of a line perpendicular to $x=3$ is:
\begin{options} \item $0$ \item $1$ \item $2$ \item $-1$ \end{options} \end{qenv}

\begin{qenv}{286} The equation $2x+3y=6$ in slope-intercept form is:
\begin{options} \item $y=-\dfrac{2}{3}x+2$ \item $y=\dfrac{2}{3}x+2$ \item $y=-\dfrac{2}{3}x-2$ \item $y=\dfrac{2}{3}x-2$ \end{options} \end{qenv}

\begin{qenv}{287} The equation $3x-4y=12$ in intercept form is:
\begin{options} \item $\dfrac{x}{4}+\dfrac{y}{3}=1$ \item $\dfrac{x}{4}-\dfrac{y}{3}=1$ \item $\dfrac{x}{3}+\dfrac{y}{4}=1$ \item $\dfrac{x}{3}-\dfrac{y}{4}=1$ \end{options} \end{qenv}

\begin{qenv}{288} The equation $5x-12y+39=0$ in the form $5x \pm 12y = k$ is:
\begin{options} \item $5x-12y=-39$ \item $5x+12y=39$ \item $5x-12y=39$ \item $5x+12y=-39$ \end{options} \end{qenv}

\begin{qenv}{289} The distance between two towns with coordinates $(3,4)$ and $(7,1)$ is:
\begin{options} \item $5$ \item $6$ \item $7$ \item $8$ \end{options} \end{qenv}

\begin{qenv}{290} The length of side $AB$ of a triangle with vertices $A(2,3)$, $B(5,7)$ and $C(6,2)$ is:
\begin{options} \item $5$ \item $6$ \item $7$ \item $8$ \end{options} \end{qenv}

\begin{qenv}{291} The distance between two cities with coordinates $(10,20)$ and $(30,40)$ is:
\begin{options} \item $20\sqrt{2}$ \item $30\sqrt{2}$ \item $40\sqrt{2}$ \item $50\sqrt{2}$ \end{options} \end{qenv}

\begin{qenv}{292} The distance between two points on a map with coordinates $(12,5)$ and $(8,-4)$ is:
\begin{options} \item $9$ \item $10$ \item $11$ \item $12$ \end{options} \end{qenv}

\begin{qenv}{293} The length of a fence around a rectangular field with corners at $(0,0)$, $(0,5)$, $(8,5)$ and $(8,0)$ is:
\begin{options} \item $40$ \item $30$ \item $36$ \item $26$ \end{options} \end{qenv}

\begin{qenv}{294} The midpoint of a journey from $(40^{\circ}\text{N}, 100^{\circ}\text{W})$ to $(50^{\circ}\text{N}, 80^{\circ}\text{W})$ is:
\begin{options} \item $(50^{\circ}\text{N}, 90^{\circ}\text{W})$ \item $(45^{\circ}\text{N}, 80^{\circ}\text{W})$ \item $(45^{\circ}\text{N}, 90^{\circ}\text{W})$ \item $(40^{\circ}\text{N}, 90^{\circ}\text{W})$ \end{options} \end{qenv}

\begin{qenv}{295} The equation of a straight line in the slope-intercept form is written as:
\begin{options} \item $y=m(x+c)$ \item $y-y_1=m(x-x_1)$ \item $y=mx+c$ \item $ax+by+c=0$ \end{options} \end{qenv}

\begin{qenv}{296} The gradients of two parallel lines are:
\begin{options} \item Equal \item Zero \item Negative reciprocals of each other \item Always undefined \end{options} \end{qenv}

\begin{qenv}{297} If the product of the gradients of two lines is $-1$, then the lines are:
\begin{options} \item Parallel \item Perpendicular \item Collinear \item Coincident \end{options} \end{qenv}

\begin{qenv}{298} Distance between two points $P(1,2)$ and $Q(4,6)$ is:
\begin{options} \item $5$ \item $6$ \item $\sqrt{13}$ \item $4$ \end{options} \end{qenv}

\begin{qenv}{299} The midpoint of a line segment with endpoints $(-2,4)$ and $(6,-2)$ is:
\begin{options} \item $(4,2)$ \item $(2,1)$ \item $(1,1)$ \item $(0,0)$ \end{options} \end{qenv}

\begin{qenv}{300} A line passing through points $(1,2)$ and $(4,5)$ is:
\begin{options} \item $y=x+1$ \item $y=2x+3$ \item $y=3x-2$ \item $y=x+2$ \end{options} \end{qenv}

\begin{qenv}{301} The equation of a line in point-slope form is:
\begin{options} \item $y=m(x+c)$ \item $y-y_1=m(x-x_1)$ \item $y=c+mx$ \item $ax+by+c=0$ \end{options} \end{qenv}

\begin{qenv}{302} $2x+3y-6=0$ in the slope-intercept form is:
\begin{options} \item $y=-\dfrac{2}{3}x+2$ \item $y=\dfrac{2}{3}x-2$ \item $y=\dfrac{2}{3}x+1$ \item $y=-\dfrac{2}{3}x-2$ \end{options} \end{qenv}

\begin{qenv}{303} The equation of a line in symmetric form is:
\begin{options} \item $\dfrac{x}{a}+\dfrac{y}{b}=1$ \item $\dfrac{x-x_1}{1}=\dfrac{y-y_1}{m}=\dfrac{z-z_1}{1}$ \item $ax+by+c=0$ \item $y-y_1=m(x-x_1)$ \end{options} \end{qenv}

\begin{qenv}{304} The equation of a line in normal form is:
\begin{options} \item $y=mx+c$ \item $\dfrac{x}{a}+\dfrac{y}{b}=1$ \item $\dfrac{x-x_1}{\cos\alpha}=\dfrac{y-y_1}{\sin\alpha}$ \item $x\cos\alpha+y\sin\alpha=p$ \end{options} \end{qenv}

\begin{qenv}{305} Which of the following is often related to inductive reasoning:
\begin{options} \item Based on repeated experiments \item If and only if statements \item Statement proven by a theorem \item Based on general principles \end{options} \end{qenv}

\begin{qenv}{306} Which of the following describes deductive reasoning:
\begin{options} \item General conclusion from a limited number of observations \item Based on repeated experiments \item Based on units of information that are accurate \item Draw conclusion from well-known facts \end{options} \end{qenv}

\begin{qenv}{307} Which one of the following statements is true:
\begin{options} \item The set of integers is finite \item The sum of the interior angles of any quadrilateral is always $180^{\circ}$ \item $\dfrac{22}{7} \notin Q'$ \item All isosceles triangles are equilateral triangles \end{options} \end{qenv}

\begin{qenv}{308} Which of the following is best to represent the negation of ``The stove is burning'':
\begin{options} \item The stove is not burning \item The stove is dim \item The stove is turned to low heat \item It is both burning and not burning \end{options} \end{qenv}

\begin{qenv}{309} The conjunction of two statements $p$ and $q$ is true when:
\begin{options} \item Both $p$ and $q$ are false \item Both $p$ and $q$ are true \item Only $q$ is true \item Only $p$ is true \end{options} \end{qenv}

\begin{qenv}{310} A conditional is regarded as false only when:
\begin{options} \item Antecedent is true and consequent is false \item Consequent is true and antecedent is false \item Antecedent is true only \item Consequent is false only \end{options} \end{qenv}

\begin{qenv}{311} Contrapositive of $q \to p$:
\begin{options} \item $q \to \sim p$ \item $\sim q \to p$ \item $\sim p \to \sim q$ \item $\sim q \to \sim p$ \end{options} \end{qenv}

\begin{qenv}{312} The statement ``Every integer greater than 2 is a sum of two prime numbers'' is a:
\begin{options} \item Theorem \item Conjecture \item Axiom \item Postulate \end{options} \end{qenv}

\begin{qenv}{313} The statement ``A straight line can be drawn between any two points'' is a:
\begin{options} \item Theorem \item Conjecture \item Axiom \item Logic \end{options} \end{qenv}

\begin{qenv}{314} The statement ``The sum of the interior angles of a triangle is $180^{\circ}$'' is a:
\begin{options} \item Converse \item Theorem \item Axiom \item Conditional \end{options} \end{qenv}

\begin{qenv}{315} A sentence or mathematical expression which may be true or false but not both is called a:
\begin{options} \item Argument \item Statement \item Proof \item Check \end{options} \end{qenv}

\begin{qenv}{316} The sum of the measures of the interior angles of a triangle is:
\begin{options} \item $90^{\circ}$ \item $180^{\circ}$ \item $270^{\circ}$ \item $360^{\circ}$ \end{options} \end{qenv}

\begin{qenv}{317} The circumference of a circle with radius $r$ is:
\begin{options} \item $2\pi r$ \item $\pi r^2$ \item $4\pi r^2$ \item None of these \end{options} \end{qenv}

\begin{qenv}{318} Symbol used for `or':
\begin{options} \item $\wedge$ \item $\vee$ \item $\sim$ \item $\equiv$ \end{options} \end{qenv}

\begin{qenv}{319} The symbol $\wedge$ stands for:
\begin{options} \item And \item Or \item Not \item Is equivalent \end{options} \end{qenv}

\begin{qenv}{320} If $q$, then the negation is written as:
\begin{options} \item $\sim q$ \item $\wedge q$ \item $\vee q$ \item $\Rightarrow q$ \end{options} \end{qenv}

\begin{qenv}{321} The conjunction of two statements $p$ and $q$ is symbolically written as:
\begin{options} \item $p \to \sim q$ \item $p \vee q$ \item $p \wedge q$ \item $p \to q$ \end{options} \end{qenv}

\begin{qenv}{322} An example of a conjunction is:
\begin{options} \item $3<5 \wedge 7<8$ \item $3<3 \wedge 10<10$ \item $2<5 \wedge 9<6$ \item None of these \end{options} \end{qenv}

\begin{qenv}{323} The statement ``Triangle has two right angles or Lahore is the capital of Sindh'' is called a:
\begin{options} \item Conjunction \item Disjunction \item Negation \item None of these \end{options} \end{qenv}

\begin{qenv}{324} A compound statement of the form $p \to q$ is called a:
\begin{options} \item Conjunction \item Disjunction \item Negation \item Conditional \end{options} \end{qenv}

\begin{qenv}{325} In $p \to q$, ``$p$'' is called the:
\begin{options} \item Hypothesis \item Conclusion \item Both hypothesis and conclusion \item None of these \end{options} \end{qenv}

\begin{qenv}{326} A conditional is regarded as \underline{\hspace{2cm}} when the antecedent is true and the consequent is false:
\begin{options} \item True \item False \item Negation \item None of these \end{options} \end{qenv}

\begin{qenv}{327} The statement $p \to q$ is called a:
\begin{options} \item Conditional \item Biconditional \item Conjunction \item Disjunction \end{options} \end{qenv}

\begin{qenv}{328} $p \to q$ is a given conditional; its converse is:
\begin{options} \item $q \to p$ \item $\sim p \to \sim q$ \item $\sim q \to \sim p$ \item None of these \end{options} \end{qenv}

\begin{qenv}{329} $p \to q$ is a given conditional; then $\sim q \to \sim p$ is called its:
\begin{options} \item Inverse \item Converse \item Contrapositive \item None of these \end{options} \end{qenv}

\begin{qenv}{330} A statement which is true for all possible values of the variable(s) involved is called a:
\begin{options} \item Tautology \item Absurdity \item Contingency \item None of these \end{options} \end{qenv}

\begin{qenv}{331} A statement which is false for all possible values is called a/an:
\begin{options} \item Tautology \item Absurdity \item Contingency \item None of these \end{options} \end{qenv}

\begin{qenv}{332} $\dfrac{a}{b}+\dfrac{c}{d} = \ldots\ldots$
\begin{options} \item $\dfrac{ac}{bd}$ \item $\dfrac{a+c}{b+d}$ \item $\dfrac{ad+bc}{bd}$ \item $\dfrac{ab+cd}{bd}$ \end{options} \end{qenv}

\begin{qenv}{333} $(x-1)^2+(x+1)^2 = \ldots\ldots\ldots\ldots\ldots$:
\begin{options} \item $4(x^2+1)$ \item $3(x^2+1)$ \item $2(x^2+1)$ \item $(x^2+1)$ \end{options} \end{qenv}

\begin{qenv}{334} A closed figure with more than four sides is called a:
\begin{options} \item Circle \item Rectangle \item Triangle \item Polygon \end{options} \end{qenv}

\begin{qenv}{335} \underline{\hspace{2cm}} figures have the same shape but different size:
\begin{options} \item Congruent \item Similar \item Equal \item None of these \end{options} \end{qenv}

\begin{qenv}{336} An example of similar \underline{\hspace{2cm}}:
\begin{options} \item Scalene triangle \item Acute angle triangle \item Equilateral triangle \item Obtuse angle triangle \end{options} \end{qenv}

\begin{qenv}{337} Symbol used for similarity:
\begin{options} \item $\equiv$ \item $\sim$ \item $=$ \item $\neq$ \end{options} \end{qenv}

\begin{qenv}{338} If $\triangle ABC$ and $\triangle DEF$ are similar, it is written symbolically as:
\begin{options} \item $\triangle ABC \equiv \triangle DEF$ \item $\triangle ABC \sim \triangle DEF$ \item $\triangle ABC > \triangle DEF$ \item $\triangle ABC < \triangle DEF$ \end{options} \end{qenv}

\begin{qenv}{339} Area of a parallelogram:
\begin{options} \item Base $+$ altitude \item Base $\times$ altitude \item $\dfrac{1}{2}$(Base $\times$ altitude) \item Length $\times$ width \end{options} \end{qenv}

\begin{qenv}{340} What is the condition for two polygons to be similar:
\begin{options} \item Corresponding angles are equal and corresponding sides are equal \item Corresponding angles are equal and corresponding sides are proportional \item Corresponding angles are proportional and corresponding sides are equal \item Corresponding angles are proportional and corresponding sides are proportional \end{options} \end{qenv}

\begin{qenv}{341} If the ratio of corresponding sides of two triangles is $3:4$, what is the ratio of their areas:
\begin{options} \item $3:4$ \item $9:8$ \item $9:16$ \item $12:16$ \end{options} \end{qenv}

\begin{qenv}{342} In triangles $ABC$ and $DEF$, if $AB=6$ cm, $DE=9$ cm and $AC=8$ cm, what is $DF$ if the triangles are similar:
\begin{options} \item $10$ cm \item $12$ cm \item $14$ cm \item $16$ cm \end{options} \end{qenv}

\begin{qenv}{343} If base $=8$ cm and altitude $=4$ cm, then area of the parallelogram:
\begin{options} \item $12$ cm$^2$ \item $4$ cm$^2$ \item $18$ cm$^2$ \item $32$ cm$^2$ \end{options} \end{qenv}

\begin{qenv}{344} If the area of a smaller polygon is $25$ cm$^2$ and the ratio of corresponding sides is $1:2$, what is the area of the larger polygon:
\begin{options} \item $100$ cm$^2$ \item $75$ cm$^2$ \item $50$ cm$^2$ \item $125$ cm$^2$ \end{options} \end{qenv}

\begin{qenv}{345} Formula used for the exterior angle of a polygon:
\begin{options} \item $\dfrac{360^{\circ}}{n}$ \item $\dfrac{180^{\circ}}{n}$ \item $\dfrac{(n-2)\times180^{\circ}}{n}$ \item None of these \end{options} \end{qenv}

\begin{qenv}{346} Size of each interior angle of a regular polygon:
\begin{options} \item $\dfrac{(n-2)\times180^{\circ}}{n}$ \item $\dfrac{180^{\circ}}{n}$ \item $\dfrac{360^{\circ}}{n}$ \item None of these \end{options} \end{qenv}

\begin{qenv}{347} What is the sum of the interior angles of a hexagon:
\begin{options} \item $720^{\circ}$ \item $540^{\circ}$ \item $900^{\circ}$ \item $1080^{\circ}$ \end{options} \end{qenv}

\begin{qenv}{348} What is the measure of each interior angle of a regular pentagon:
\begin{options} \item $72^{\circ}$ \item $108^{\circ}$ \item $120^{\circ}$ \item $144^{\circ}$ \end{options} \end{qenv}

\begin{qenv}{349} What is the measure of each exterior angle of a regular octagon:
\begin{options} \item $30^{\circ}$ \item $90^{\circ}$ \item $60^{\circ}$ \item $45^{\circ}$ \end{options} \end{qenv}

\begin{qenv}{350} How many diagonals does a regular heptagon have:
\begin{options} \item $7$ \item $21$ \item $14$ \item $28$ \end{options} \end{qenv}

\begin{qenv}{351} What is the sum of the exterior angles of any polygon:
\begin{options} \item $180^{\circ}$ \item $360^{\circ}$ \item $540^{\circ}$ \item $720^{\circ}$ \end{options} \end{qenv}

\begin{qenv}{352} In an equilateral triangle, what is the measure of each angle:
\begin{options} \item $45^{\circ}$ \item $60^{\circ}$ \item $90^{\circ}$ \item $120^{\circ}$ \end{options} \end{qenv}

\begin{qenv}{353} In an isosceles triangle, if the vertex angle is $80^{\circ}$, what is the measure of each of the other two (equal) angles:
\begin{options} \item $50^{\circ}$ \item $60^{\circ}$ \item $70^{\circ}$ \item $80^{\circ}$ \end{options} \end{qenv}

\begin{qenv}{354} In a parallelogram, what is the relationship between opposite sides:
\begin{options} \item Equal and parallel \item Equal but not parallel \item Parallel but not equal \item Neither equal nor parallel \end{options} \end{qenv}

\begin{qenv}{355} What is the measure of each angle in a rectangle:
\begin{options} \item $60^{\circ}$ \item $120^{\circ}$ \item $90^{\circ}$ \item $180^{\circ}$ \end{options} \end{qenv}

\begin{qenv}{356} In a rhombus, what is the relationship between the diagonals:
\begin{options} \item Equal \item Unequal \item Parallel \item Bisect each other at right angles \end{options} \end{qenv}

\begin{qenv}{357} If two polygons are similar, then:
\begin{options} \item Their corresponding angles are equal \item Their areas are equal \item Their volumes are equal \item Their corresponding sides are equal \end{options} \end{qenv}

\begin{qenv}{358} The ratio of the areas of two similar polygons is:
\begin{options} \item Equal to the ratio of their perimeters \item Equal to the square of the ratio of their corresponding sides \item Equal to the cube of the ratio of their corresponding sides \item Equal to the sum of their corresponding sides \end{options} \end{qenv}

\begin{qenv}{359} If the volumes of two similar solids are $125\,\text{cm}^3$ and $27\,\text{cm}^3$, the ratio of their corresponding heights is $\ldots\ldots$:
\begin{options} \item $3:5$ \item $5:3$ \item $25:9$ \item $9:25$ \end{options} \end{qenv}



\begin{mcqB}[start=360]

\item The exterior angle of a regular pentagon is:
\opts{$40^{\circ}$}{$45^{\circ}$}{$60^{\circ}$}{$72^{\circ}$}

\item A parallelogram has an area of $64\ \text{cm}^2$ and a similar parallelogram has an area of $144\ \text{cm}^2$. If a side of the smaller parallelogram is 8 cm, the corresponding side of the larger parallelogram is:
\opts{$10$ cm}{$12$ cm}{$18$ cm}{$16$ cm}

\item The total number of diagonals in a polygon with 9 sides is:
\opts{$18$}{$21$}{$25$}{$27$}

\item Two shapes are similar, and their radii are in the ratio $4:5$. If the surface area of the larger sphere is $500\pi\ \text{cm}^2$, the surface area of the smaller sphere is:
\opts{$256\pi\ \text{cm}^2$}{$320\pi\ \text{cm}^2$}{$400\pi\ \text{cm}^2$}{$405\pi\ \text{cm}^2$}

\item A regular polygon has an exterior angle of $30^{\circ}$. How many diagonals does the polygon have:
\opts{$54$}{$90$}{$72$}{$108$}

\item In a regular hexagon, the ratio of the length of a diagonal to the side length is:
\opts{$\sqrt{3}:1$}{$2:1$}{$3:2$}{$2:3$}

\item A regular polygon has an interior angle of $165^{\circ}$. How many sides does it have:
\opts{$15$}{$16$}{$20$}{$24$}

\item What is the general form of a linear function?
\opts{$f(x)=mx+c$}{$f(x)=ax^2+bx+c$}{$f(x)=ax^3+bx^2+cx+d$}{$f(x)=ke^{2x}$}

\item In the linear function $y=3x-5$, what is the slope $m$:
\opts{$5$}{$3$}{$-3$}{$-5$}

\item What is the y-intercept of the linear function $y=-2x+8$:
\opts{$-2$}{$2$}{$-8$}{$8$}

\item The graph of $y=x^2$ represents:
\opts{A straight line}{A parabola opening upwards}{A parabola opening downwards}{A cubic curve}

\item In the quadratic function $y=-x^2$, what is the direction of the parabola:
\opts{Upwards}{Downwards}{Left}{Right}

\item The standard form of cubic functions:
\opts{$y=ax^3+bx^2+cx+d$}{$y=ax^2+bx+c$}{$y=ax+b$}{None of these}

\item The graph of $y=2x^3-3x^2+x-5$ passes through which point:
\opts{$(0,-5)$}{$(0,0)$}{$(0,5)$}{$(0,-10)$}

\item What is the general form of a reciprocal function?
\opts{$y=ax+b$}{$y=ax^2+bx+c$}{$y=\dfrac{q}{x}$}{$y=x^3$}

\item What type of function is shown (equation):
\opts{Exponential}{Linear}{Quadratic}{Cubic}

\item What does the graph of $y=x^3$ look like:
\opts{A straight line}{A parabola}{A cubic curve passing through the origin}{A hyperbola}

\item The graph of the function $y=5^x$ represents:
\opts{Exponential growth}{Exponential decay}{A line passing through the origin}{A hyperbola}

\item For the quadratic function $y=x^2+2x-3$, what is the vertex:
\opts{$(-1,-3)$}{$(0,3)$}{$(1,2)$}{$(-2,1)$}

\item The graph of $y=-x^2+4$ cuts the y-axis at:
\opts{$(0,4)$}{$(0,-4)$}{$(4,0)$}{$(-4,0)$}

\item What is the slope of the line $y=2x+3$:
\opts{$2$}{$3$}{$-2$}{$0$}

\item What is the slope of the line $y=-3x+2$:
\opts{$3$}{$-3$}{$2$}{$-2$}

\item The equation $y=3x$ represents which type of graph:
\opts{A line with slope 3}{A parabola}{A cubic curve}{Exponential decay}

\item The gradient of the tangent to the curve $y=2x^2-3x+1$ at $x=1$:
\opts{$1$}{$2$}{$3$}{$4$}

\item The equation $y=-2x^2+4x-3$ represents which type of graph:
\opts{Parabola opening downwards}{Parabola opening upwards}{Line}{Exponential growth}

\item In a reciprocal function $y=\dfrac{1}{x}$, as $x$ approaches zero, what happens to $y$?
\opts{It approaches infinity}{It approaches zero}{It becomes undefined}{It remains constant}

\item The general form of a quadratic function is:
\opts{$y=mx+c$}{$y=ax^2+bx+c$}{$y=ax^3+bx^2+cx+d$}{$y=k\cdot e^{2x}$}

\item In the function $y=x^2+4x+4$, what is the value of $y$ when $x=-2$:
\opts{$0$}{$4$}{$-4$}{$8$}

\item In the equation $y=5x-2$, what is the y-intercept:
\opts{$-5$}{$-2$}{$2$}{$5$}

\item What is the break-even point when $C(x)=1200+20x$ and $R(x)=50x$:
\opts{$40$}{$50$}{$60$}{$75$}

\item In the function $y=x^3+3x^2$, what is the slope of the tangent at $x=0$:
\opts{$0$}{$3$}{$6$}{$-3$}

\item The function $y=3x^2+2x-1$ has its minimum value at:
\opts{$x=0$}{$x=-1$}{$x=1$}{$x=2$}

\item In the function $y=x^2+4x$, what is the axis of symmetry:
\opts{$x=0$}{$x=-2$}{$x=2$}{$x=-4$}

\item Which of the following is the domain of the reciprocal function:
\opts{All real numbers except zero}{All real numbers}{Positive real numbers}{Negative real numbers}

\item The graph of $y=2x^3-x^2-5$ is:
\opts{Linear}{Quadratic}{Cubic}{Exponential}

\item The graph of $y=2x^3-x^2-5$ is:
\opts{Linear}{Quadratic}{Cubic}{Exponential}

\item For the function $y=3x-4$, the graph is a:
\opts{Line}{Parabola}{Hyperbola}{Exponential function}

\item The graph of $y=x^3$ passes through the point:
\opts{$(0,0)$}{$(1,0)$}{$(0,1)$}{$(0,-1)$}

\item What is the y-intercept of the equation $y=4x-6$:
\opts{$-6$}{$6$}{$4$}{$-4$}

\item The cost function is $C(x)=1500+10x$. The revenue function is $R(x)=50x$. What is the break-even point:
\opts{$30$ units}{$50$ units}{$100$ units}{$150$ units}

\item The graph of $y=x^2-2x+1$ represents:
\opts{A line}{A parabola}{A cubic curve}{An exponential curve}

\item In the function $y=2x^2+3x+5$, the value of $y$ when $x=-1$ is:
\opts{$4$}{$2$}{$1$}{$3$}

\item The graph of $y=x^2-4$ intersects the x-axis at:
\opts{$(2,0)$ and $(-2,0)$}{$(4,0)$ and $(-4,0)$}{$(1,0)$ and $(-1,0)$}{$(3,0)$ and $(-3,0)$}

\item The function $y=x^4$ is a:
\opts{Linear function}{Quadratic function}{Cubic function}{Quartic function}

\item The graph of $y=x^2+x-2$ represents:
\opts{A straight line}{A hyperbola}{A parabola}{An exponential function}

\item For the function $y=x^2+2x+1$, what is the vertex of the parabola:
\opts{$(0,0)$}{$(-1,0)$}{$(-1,1)$}{$(1,1)$}

\item The graph of $y=x^2+5$ intersects the x-axis at:
\opts{$(0,5)$}{$(5,0)$}{$(-5,0)$}{$(1,0)$}

\item The function $y=3x^2$ represents:
\opts{A line}{A parabola opening upwards}{A parabola opening downwards}{A cubic function}

\item The function $y=3x^2$ represents:
\opts{A line}{A parabola opening upwards}{A parabola opening downwards}{A cubic function}

\item The function $y=2^x$ represents what kind of growth:
\opts{Linear growth}{Exponential growth}{Exponential decay}{Quadratic growth}

\item What is the gradient of the tangent to the curve $y=x^2$ at $x=1$:
\opts{$4$}{$2$}{$0$}{$8$}

\item Which function represents exponential decay:
\opts{$y=2^x$}{$y=e^x$}{$y=2^{-x}$}{$y=x^2$}

\item In the function $y=x^2-3x-2$, what is the value of $y$ when $x=0$:
\opts{$-2$}{$0$}{$2$}{$-3$}

\item The graph of the function $y=x^2-2$ represents a:
\opts{Parabola opening upwards}{Parabola opening downwards}{Hyperbola}{Circle}

\item What is the break-even point when the cost function is $C(x)=90{,}000+600x$ and the revenue function is $R(x)=1{,}800x$:
\opts{$50$ footballs}{$75$ footballs}{$100$ footballs}{$150$ footballs}

\item Which of the following is an exponential growth function:
\opts{$y=x^2$}{$y=3x$}{$y=2^x$}{$y=x^3$}

\item Which of the following represents an exponential decay function:
\opts{$y=2^x$}{$y=e^{-x}$}{$y=x^2$}{$y=3x$}

\item The graph of $y=e^x$ represents:
\opts{Linear growth}{Exponential growth}{Parabolic growth}{Cubic growth}

\item $x=5$ represents:
\opts{The $x$-axis}{The $y$-axis}{A line parallel to the $x$-axis}{A line parallel to the $y$-axis}

\item Slope of the line $y=5x+3$ is:
\opts{$3$}{$-3$}{$5$}{$-5$}

\item The y-intercept of $y=-2x-1$ is:
\opts{$-2$}{$2$}{$-1$}{$1$}

\item The graph of $y=x^3$ cuts the x-axis at:
\opts{$x=0$}{$x=1$}{$x=-1$}{$x=2$}

\item The graph of $3^x$ cuts the x-axis at (behaviour):
\opts{Growth}{Decay}{Both}{A line}

\item The graph of $y=-x^2+5$ opens:
\opts{Upward}{Downward}{Left side}{Right side}

\item The graph of $y=x^2-9$ opens:
\opts{Upward}{Downward}{Left side}{Right side}

\item $y=5^x$ is ...... function:
\opts{Linear}{Quadratic}{Cubic}{Exponential}

\item A reciprocal function is:
\opts{$y=7^x$}{$y=\dfrac{2}{x}$}{$y=2x^2$}{$y=5x^3$}

\item $y=-3x^3+7$ is ...... function:
\opts{Exponential}{Cubic}{Linear}{Reciprocal}

\item The point of intersection of the perpendicular bisectors of the sides of a triangle is called:
\opts{Incenter}{Centroid}{Circumcenter}{Orthocenter}

\item If the measure of one angle in a triangle is greater than $90^{\circ}$, which type of triangle is formed:
\opts{Acute triangle}{Right-angled triangle}{Obtuse triangle}{Isosceles triangle}

\item The point where all the medians of a triangle meet is known as:
\opts{Circumcenter}{Incenter}{Centroid}{Orthocenter}

\item Which locus defines a point that is a fixed distance from a given point:
\opts{Perpendicular bisector}{Circle}{Angle bisector}{Line segment}

\item What is the angle bisector of a triangle:
\opts{A line passing through a vertex}{A line that divides an angle into two equal parts}{A line joining midpoints of sides}{A line perpendicular to a side}

\item The orthocenter of a triangle is where:
\opts{The altitudes of the triangle meet}{The sides are perpendicular}{The circumcenter meets the centroid}{None of these}

\item The intersection of the three angle bisectors of a triangle is called:
\opts{Orthocenter}{Centroid}{Circumcenter}{Incenter}

\item An equilateral triangle is a(n) ...... angled triangle:
\opts{Obtuse}{Right}{Acute}{None of these}

\item The center of a circle that passes through all vertices of a triangle is located at the:
\opts{Incenter}{Circumcenter}{Centroid}{Orthocenter}

\item How many types of a triangle are there w.r.t. angles:
\opts{$2$}{$3$}{$4$}{$5$}

\item A triangle with all sides and angles equal is called a:
\opts{Right triangle}{Isosceles triangle}{Equilateral triangle}{Scalene triangle}

\item The centroid of a triangle is the point of concurrency of the:
\opts{Angle bisectors}{Altitudes}{Perpendicular bisectors}{Medians}

\item When constructing a right-angled triangle, the side opposite the right angle is called the:
\opts{Hypotenuse}{Median}{Altitude}{Base}

\item Which of the following is true for an equilateral triangle:
\opts{It has only one angle of $90^{\circ}$}{All angles are $60^{\circ}$}{It has two equal sides}{The sides are all different}

\item What is the sum of the interior angles of any triangle:
\opts{$180^{\circ}$}{$360^{\circ}$}{$90^{\circ}$}{$120^{\circ}$}

\item A triangle with two equal sides is called a:
\opts{Equilateral triangle}{Scalene triangle}{Isosceles triangle}{Right triangle}

\end{mcqB}



\mcqC{How many types of a triangle w.r.t sides:}{2}{3}{4}{5}

\mcqC{The centroid of a triangle divides each median into two parts in a ratio of:}{1:2}{2:1}{1:3}{3:1}

\mcqC{Which of the following is not a type of triangle based on sides:}{Scalene}{Isosceles}{Equilateral}{Acute}

\mcqC{Perpendicular from a vertex to the opposite side in a triangle is called:}{Altitude}{Median}{Bisector}{Angle}

\mcqC{A triangle with one angle equal to $90^{\circ}$ is called a:}{Scalene triangle}{Equilateral triangle}{Right-angled triangle}{Isosceles triangle}

\mcqC{A triangle is a closed figure having:}{Two sides}{Three sides}{Four sides}{Five sides}

\mcqC{When constructing a triangle, which condition requires the measure of one angle and two sides:}{SAS criterion}{ASA criterion}{SSS criterion}{AAS criterion}

\mcqC{What is the angle formed by two intersecting rays from a vertex in a triangle called:}{Complementary angle}{Opposite angle}{Vertical angle}{Exterior angle}

\mcqC{A triangle can be constructed if the sum of the measure of any two sides is ...... the measure of the third side:}{Less than}{Greater than}{Equal to}{Greater than and equal to}

\mcqC{An equilateral triangle ......:}{Can be isosceles}{Can be right angled}{Can be obtuse angled}{Has each angle equal to $50^{\circ}$}

\mcqC{If the sum of the measure of two angles is less than $90^{\circ}$, then the triangle is ......:}{Equilateral}{Acute angled}{Obtuse angled}{Right angled}

\mcqC{The line segment joining the midpoint of a side to its opposite vertex in a triangle is called ......:}{Median}{Perpendicular bisector}{Angle bisector}{Circle}

\mcqC{The angle bisectors of a triangle intersect at ......:}{One point}{Two points}{Three points}{Four points}

\mcqC{Locus of all points equidistant from a fixed point is ......:}{Circle}{Perpendicular bisector}{Angle bisector}{Parallel lines}

\mcqC{Locus of points equidistant from two fixed points is ......:}{Circle}{Perpendicular bisector}{Angle bisector}{Parallel lines}

\mcqC{Locus of points equidistant from a fixed line is/are ......:}{Circle}{Perpendicular bisector}{Angle bisector}{Parallel lines}

\mcqC{Locus of points equidistant from two intersecting lines is ......:}{Circle}{Perpendicular bisector}{Angle bisector}{Parallel lines}

\mcqC{The set of all points which is farther than 2 km from a fixed point B is a region outside a circle of radius ...... and centre at B:}{1 km}{1.9 km}{2 km}{2.1 km}

\mcqC{What is the process of collecting, organizing, summarizing, analyzing, and interpreting numerical data called:}{Data Collection}{Information Handling}{Data Analysis}{Data Interpretation}

\mcqC{Which of the following is an example of discrete data:}{Height of students}{Weight of students}{Number of books sold}{Temperature of a city}

\mcqC{What is the term for data that can take every possible value in a given interval:}{Discrete data}{Continuous data}{Grouped data}{Ungrouped data}

\mcqC{Which of the following is an example of ungrouped data:}{Marks of students arranged in classes}{Number of toys sold in a month}{Height of students in a class}{Temperature recorded daily}

\mcqC{What is the term for data arranged in systematic order (groups or classes):}{Ungrouped data}{Grouped data}{Discrete data}{Continuous data}

\mcqC{Who used the word ``data handling'' first:}{Newton}{Henry}{Sir Ronald Fisher}{Arthur Cayley}

\mcqC{What is the formula to calculate the size of class limits:}{Size = Range / Number of classes}{Size = Range $\times$ Number of classes}{Size = Range + Number of classes}{Size = Range -- Number of classes}

\mcqC{What is the range of the data if the greatest value is 136 and the smallest value is 30:}{106}{166}{96}{136}

\mcqC{What is the term for the difference between the upper and lower class limits:}{Class size}{Class boundary}{Class interval}{Class frequency}

\mcqC{In class 10-14, the lower class limit is:}{9.5}{10}{12}{146}

\mcqC{What is a histogram:}{A graph of adjacent rectangles}{A graph of connected lines}{A graph of pie charts}{A graph of bar charts}

\mcqC{What is the term for the midpoint of a class interval:}{Class boundary}{Class mark}{Class limit}{Class frequency}

\mcqC{What is the formula to calculate the midpoint of a class interval?}{Midpoint = (Lower limit + Upper limit)/2}{Midpoint = (Lower -- Upper)/2}{Midpoint = (Lower limit $\times$ Upper limit)/2}{None of these}

\mcqC{In class 10-14, frequency is 10; represent it using tally marks:}{Tally marks grouped in fives (10 marks)}{Ten single strokes}{Ten strokes crossed in fives}{None of these}

\mcqC{Class boundaries of 5-9 is:}{5.5--10.5}{4.5--8.5}{4.5--9.5}{4--10}

\mcqC{Mode of first 20 whole numbers:}{0}{20}{19}{No Mode}

\mcqC{What is the median of the data set 5, 8, 12, 13, 16, 18, 20:}{12}{13}{16}{18}

\mcqC{What is the mode of the data set 10, 15, 15, 20, 20, 20, 25:}{15}{20}{25}{No mode}

\mcqC{What is the term for the most frequently occurring value in a data set:}{Mean}{Median}{Mode}{Range}

\mcqC{What is the formula for the weighted mean?}{Weighted mean $=\dfrac{\sum(wx)}{\sum w}$}{Weighted mean $=(wx)\times \sum w$}{Weighted mean $=(wx)-\sum w$}{Weighted mean $=(wx)+\sum w$}

\mcqC{What is the weighted mean if the weights are 4, 3, 5, 6 and the values are 74, 78, 74, 90:}{80}{85}{90}{95}

\mcqC{What is the term for the average where different observations have different weights:}{Arithmetic mean}{Weighted mean}{Median}{Mode}

\mcqC{What is the average number of packets of medicine distributed by a company if the total number of packets is 114 and the number of areas is 7:}{16}{17}{18}{19}

\mcqC{What is the average budget for the next year if the yearly budgets are 5, 7, 8, 6, 7 million:}{6.6}{7.0}{7.5}{8.0}

\mcqC{Arithmetic mean is denoted by:}{$\tilde{X}$}{$\bar{X}$}{$\hat{X}$}{SD}

\mcqC{What is the average marks of a student if the marks are 84, 91, 72, 68, 87, 78:}{80}{82}{84}{86}

\mcqC{Symbol used for sum:}{$\equiv$}{$\sum$}{\%}{$\sim$}

\mcqC{What is the median of the data set 55, 53, 54, 58, 60, 61, 62, 56, 57, 51, 63:}{56}{57}{58}{59}

\mcqC{Formula used for arithmetic mean for ungrouped data:}{$\frac{n+1}{2}$th term}{$\frac{\sum x}{n}$}{Both a and b}{None of these}

\mcqC{What is the median of the data set 110, 125, 122, 130, 124, 127, 120:}{124}{125}{126}{127}

\mcqC{What is the mode of the data set 88, 70, 72, 125, 95, 81, 90, 95, 90:}{90}{95}{90 and 95}{No mode}

\mcqC{Symbol $\tilde{X}$ is used for:}{Mode}{Median}{Range}{Arithmetic mean}

\mcqC{What is the arithmetic mean of the data set 4, 6, 10, 12, 15, 20, 25, 28, 30:}{16.67}{18.00}{20.00}{22.00}

\mcqC{What is the arithmetic mean of the data set 12, 18, 19, 0, -19, -18, -12:}{0}{1}{2}{3}

\mcqC{$l+\dfrac{f_m-f_1}{(f_m-f_1)+(f_m-f_2)}\times h$ is used for:}{Mode}{Median}{Range}{Arithmetic mean}

\mcqC{Which data takes only some specific values:}{Continuous data}{Discrete data}{Grouped data}{Ungrouped data}

\mcqC{The number of times a value occurs in a data set is called:}{Frequency}{Relative frequency}{Class limit}{Class boundaries}

\mcqC{Midpoint is also known as:}{Mean}{Median}{Class limit}{Class mark}

\mcqC{A frequency polygon is also drawn / constructed by using:}{Histogram}{Bar graph}{Class boundaries}{Class limit}

\mcqC{The difference between the greatest value and the smallest value is called:}{Class limit}{Midpoint}{Relative frequency}{Range}

\mcqC{Measure of central tendency is used to find out the ...... of a data set:}{Class boundaries}{Cumulative frequency}{Middle or centre value}{Frequency}

\mcqC{If the mean of 5, 7, 8, 9 and $x$ is 7.5, what will be the value of $x$:}{10}{8}{8.5}{5.8}

\mcqC{Find the mode of the given data 2, 5, 8, 9, 0, 1, 3, 7 and 10:}{5}{7}{0}{No mode}

\mcqC{In a data set, the value (observation) that appears or occurs most often is called:}{Mean}{Mode}{Median}{Weighted mean}

\mcqC{Find the median of the given data 110, 125, 122, 130, 124 and 120:}{124}{120}{125}{127}

\mcqC{What is the formula to calculate probability:}{$P(A)=\dfrac{n(A)}{n(S)}$}{$P(A)=\dfrac{n(S)}{n(A)}$}{$P(A)=n(A)\times n(S)$}{$P(A)=n(S)-n(A)$}

\mcqC{Which of the following is not an example of an experiment:}{Tossing a coin}{Rolling a dice}{Guessing a number}{Drawing a card from a deck}

\mcqC{What is the probability of getting a head when tossing a fair coin?}{$\frac{1}{2}$}{$\frac{1}{3}$}{$\frac{1}{4}$}{$\frac{1}{6}$}

\mcqC{If a fair dice is rolled, what is the probability of getting a number divisible by 3?}{$\frac{1}{6}$}{$\frac{1}{4}$}{$\frac{1}{3}$}{$\frac{1}{2}$}

\mcqC{What is the sample space when a single coin is tossed:}{$\{A,B\}$}{$\{1,2,3\}$}{$\{1,2,3,4\}$}{$\{H,T\}$}

\mcqC{If a dice is rolled twice, what is the total number of possible outcomes:}{6}{36}{12}{24}

\mcqC{If event A has 3 favorable outcomes, and the sample space contains 6 possible outcomes, what is $P(A)$:}{$\frac{1}{6}$}{$\frac{1}{4}$}{$\frac{1}{3}$}{$\frac{1}{2}$}

\mcqC{The complement of event A is:}{$P(S)$}{$P(A)$}{$P(A')$}{$P(A \text{ and } B)$}

\mcqC{What is the probability of not getting a 6 on a fair dice?}{$\frac{1}{6}$}{$\frac{5}{6}$}{$\frac{1}{3}$}{$\frac{2}{3}$}

\mcqC{If a fair dice is rolled, what is the probability of getting an even number?}{$\frac{1}{2}$}{$\frac{1}{3}$}{$\frac{1}{4}$}{$\frac{1}{6}$}

\mcqC{When a fair dice is rolled twice, what is the probability of getting a sum of 8:}{$\frac{1}{9}$}{$\frac{5}{36}$}{$\frac{1}{36}$}{$\frac{1}{12}$}

\mcqC{What is the probability of getting a double six when rolling two dice:}{$\frac{1}{36}$}{$\frac{1}{18}$}{$\frac{1}{12}$}{$\frac{1}{6}$}

\mcqC{The probability of an event is always:}{Less than 0}{Greater than 1}{Equal to 1}{Between 0 and 1}

\mcqC{The probability of selecting a blue ball from a bag of 5 blue and 8 green balls is:}{$\frac{5}{13}$}{$\frac{8}{13}$}{$\frac{1}{3}$}{$\frac{5}{8}$}




\begin{qenv}{521} What is the probability of getting neither a spade nor a heart when drawing a card from a deck of 52?
\begin{options} \item $\dfrac{26}{52}$ \item $\dfrac{26}{39}$ \item $\dfrac{1}{2}$ \item $\dfrac{1}{4}$ \end{options} \end{qenv}

\begin{qenv}{522} If the probability of an event is $0.3$, the probability of its complement is:
\begin{options} \item $0.4$ \item $0.5$ \item $0.3$ \item $0.7$ \end{options} \end{qenv}

\begin{qenv}{523} What is the probability of drawing a red heart from a deck of 52 playing cards?
\begin{options} \item $\dfrac{1}{13}$ \item $\dfrac{1}{4}$ \item $\dfrac{1}{2}$ \item $\dfrac{1}{26}$ \end{options} \end{qenv}

\begin{qenv}{524} What is the probability of not getting a red ball from a bag containing red and 5 blue balls?
\begin{options} \item $\dfrac{5}{13}$ \item $\dfrac{8}{13}$ \item $\dfrac{5}{8}$ \item $\dfrac{8}{5}$ \end{options} \end{qenv}

\begin{qenv}{525} If 3 dice are rolled, what is the probability of getting at least one 6:
\begin{options} \item $\dfrac{1}{36}$ \item $\dfrac{3}{36}$ \item $\dfrac{1}{6}$ \item $\dfrac{5}{36}$ \end{options} \end{qenv}

\begin{qenv}{526} The probability of getting an even sum when rolling two dice is:
\begin{options} \item $\dfrac{5}{36}$ \item $\dfrac{18}{36}$ \item $\dfrac{9}{36}$ \item $\dfrac{12}{36}$ \end{options} \end{qenv}

\begin{qenv}{527} The probability of drawing a red card from a deck of 52 cards is:
\begin{options} \item $\dfrac{1}{2}$ \item $\dfrac{13}{52}$ \item $\dfrac{26}{52}$ \item $\dfrac{1}{26}$ \end{options} \end{qenv}

\begin{qenv}{528} What is the probability of drawing a white ball from a bag containing 5 white and 7 black balls?
\begin{options} \item $\dfrac{1}{3}$ \item $\dfrac{7}{12}$ \item $\dfrac{1}{2}$ \item $\dfrac{5}{12}$ \end{options} \end{qenv}

\begin{qenv}{529} The relative frequency is:
\begin{options} \item Total number of favorable outcomes \item Total number of possible outcomes \item How often an event occurs \item The number of trials \end{options} \end{qenv}

\begin{qenv}{530} The relative frequency of selecting a red ball from 80 students is:
\begin{options} \item $0.29$ \item $0.19$ \item $0.31$ \item $0.12$ \end{options} \end{qenv}

\begin{qenv}{531} If six dice are rolled 50 times, the expected frequency of a number appearing is given by:
\begin{options} \item $NP(B)$ \item $NP(A)$ \item $NP(X)$ \item None of these \end{options} \end{qenv}

\begin{qenv}{532} The relative frequency of green in a survey of 80 students is:
\begin{options} \item $\dfrac{15}{80}$ \item $0.19$ \item $0.31$ \item $\dfrac{10}{80}$ \end{options} \end{qenv}

\begin{qenv}{533} If a probability is $\dfrac{1}{5}$, the expected frequency in 100 trials is:
\begin{options} \item $20$ \item $50$ \item $10$ \item $25$ \end{options} \end{qenv}

\begin{qenv}{534} Each element of the sample space is called:
\begin{options} \item Event \item Experiment \item Sample point \item Outcomes \end{options} \end{qenv}

\begin{qenv}{535} An outcome which represents how many times we expect something to happen is called:
\begin{options} \item Outcomes \item Favourable outcome \item Sample space \item Sample point \end{options} \end{qenv}

\begin{qenv}{536} Which one tells us how often a specific event occurs relative to the total number of events:
\begin{options} \item Expected frequency \item Sum of relative frequency \item Relative frequency \item Frequency \end{options} \end{qenv}

\begin{qenv}{537} Estimated probability of an event occurring is also known as:
\begin{options} \item Relative frequency \item Expected frequency \item Class boundaries \item Sum of expected frequency \end{options} \end{qenv}

\begin{qenv}{538} The sum of all expected frequencies is equal to the fixed number of:
\begin{options} \item Trials \item Relative frequencies \item Outcomes \item Events \end{options} \end{qenv}

\begin{qenv}{539} Find out the total number of possible sample space when 4 dice are rolled:
\begin{options} \item $6^2$ \item $6^3$ \item $6^4$ \item $6^6$ \end{options} \end{qenv}

\begin{qenv}{540} The chance of occurrence of a particular event is called:
\begin{options} \item Sample space \item Estimated probability \item Probability \item Expected frequency \end{options} \end{qenv}

\begin{qenv}{541} While rolling a pair of dice, what will be the probability of double 2?
\begin{options} \item $\dfrac{1}{6}$ \item $\dfrac{1}{3}$ \item $\dfrac{5}{6}$ \item $\dfrac{1}{36}$ \end{options} \end{qenv}

\begin{qenv}{542} A card is chosen from a pack of 52 playing cards; find the probability of getting no jack and no king:
\begin{options} \item $\dfrac{2}{13}$ \item $\dfrac{11}{13}$ \item $\dfrac{2}{52}$ \item $\dfrac{11}{52}$ \end{options} \end{qenv}

\begin{qenv}{543} An event which will probably occur, having a greater chance to occur, is called:
\begin{options} \item Equally likely event \item Likely event \item Unlikely event \item Certain event \end{options} \end{qenv}


\clearpage

\section*{Part 2 -- Short Questions (Q.1--825)}
\addcontentsline{toc}{section}{Part 2 -- Short Questions (Q.1--825)}

\bigskip
\noindent\textbf{1.} Identify the following as a rational or irrational number: $2.353535\ldots$

\bigskip
\noindent\textbf{2.} Identify the following as a rational or irrational number: $0.\overline{6}$

\bigskip
\noindent\textbf{3.} Identify the following as a rational or irrational number: $2.236067\ldots$

\bigskip
\noindent\textbf{4.} Identify the following as a rational or irrational number: $\sqrt{7}$

\bigskip
\noindent\textbf{5.} Identify the following as a rational or irrational number: $e$

\bigskip
\noindent\textbf{6.} Identify the following as a rational or irrational number: $\pi$

\bigskip
\noindent\textbf{7.} Identify the following as a rational or irrational number: $5+\sqrt{11}$

\bigskip
\noindent\textbf{8.} Identify the following as a rational or irrational number: $\sqrt{3}+\sqrt{13}$

\bigskip
\noindent\textbf{9.} Identify the following as a rational or irrational number: $\dfrac{15}{4}$

\bigskip
\noindent\textbf{10.} Identify the following as a rational or irrational number: $(2-\sqrt{2})(2+\sqrt{2})$

\bigskip
\noindent\textbf{11.} Represent the following number on a number line: $\sqrt{2}$

\bigskip
\noindent\textbf{12.} Represent the following number on a number line: $\sqrt{3}$

\bigskip
\noindent\textbf{13.} Represent the following number on a number line: $4\dfrac{1}{3}$

\bigskip
\noindent\textbf{14.} Represent the following number on a number line: $-2\dfrac{1}{7}$

\bigskip
\noindent\textbf{15.} Represent the following number on a number line: $\dfrac{5}{8}$

\bigskip
\noindent\textbf{16.} Represent the following number on a number line: $2\dfrac{3}{4}$

\bigskip
\noindent\textbf{17.} Express the following as a rational number $\dfrac{p}{q}$, where $p$ and $q$ are integers and $q\neq0$: $0.\overline{4}$

\bigskip
\noindent\textbf{18.} Express the following as a rational number $\dfrac{p}{q}$, where $p$ and $q$ are integers and $q\neq0$: $0.\overline{37}$

\bigskip
\noindent\textbf{19.} Express the following as a rational number $\dfrac{p}{q}$, where $p$ and $q$ are integers and $q\neq0$: $0.\overline{21}$

\bigskip
\noindent\textbf{20.} Name the property used in the following: $(a+4)+b = a+(4+b)$

\bigskip
\noindent\textbf{21.} Name the property used in the following: $\sqrt{2}+\sqrt{3} = \sqrt{3}+\sqrt{2}$

\bigskip
\noindent\textbf{22.} Name the property used in the following: $x-x=0$

\bigskip
\noindent\textbf{23.} Name the property used in the following: $a(b+c)=ab+ac$

\bigskip
\noindent\textbf{24.} Name the property used in the following: $16+0=16$

\bigskip
\noindent\textbf{25.} Name the property used in the following: $100\times1=100$

\bigskip
\noindent\textbf{26.} Name the property used in the following: $4\times(5\times8)=(4\times5)\times8$

\bigskip
\noindent\textbf{27.} Name the property used in the following: $ab=ba$

\bigskip
\noindent\textbf{28.} Name the property used in the following: $-3<-1 \Rightarrow 0<2$

\bigskip
\noindent\textbf{29.} Name the property used in the following: If $a<b$ then $\dfrac{1}{a}>\dfrac{1}{b}$

\bigskip
\noindent\textbf{30.} Name the property used in the following: If $a<b$ then $a+c<b+c$

\bigskip
\noindent\textbf{31.} Name the property used in the following: If $ac<bc$ and $c>0$ then $a<b$

\bigskip
\noindent\textbf{32.} Name the property used in the following: If $ac<bc$ and $c<0$ then $a>b$

\bigskip
\noindent\textbf{33.} Name the property used in the following: Either $a>b$ or $a=b$ or $a<b$

\bigskip
\noindent\textbf{34.} Insert two rational numbers between $\dfrac{1}{3}$ and $\dfrac{1}{4}$.

\bigskip
\noindent\textbf{35.} Insert two rational numbers between $3$ and $4$.

\bigskip
\noindent\textbf{36.} Insert two rational numbers between $\dfrac{3}{5}$ and $\dfrac{4}{5}$.

\bigskip
\noindent\textbf{37.} Define terminating decimal numbers.

\bigskip
\noindent\textbf{38.} Define non-terminating decimal numbers.

\bigskip
\noindent\textbf{39.} Define rational numbers.

\bigskip
\noindent\textbf{40.} Define irrational numbers.

\bigskip
\noindent\textbf{41.} Identify the following decimal number as rational or irrational: $0.35$

\bigskip
\noindent\textbf{42.} Identify the following decimal number as rational or irrational: $0.444\ldots$

\bigskip
\noindent\textbf{43.} Identify the following decimal number as rational or irrational: $3.\overline{5}$

\bigskip
\noindent\textbf{44.} Identify the following decimal number as rational or irrational: $3.36788542\ldots$

\bigskip
\noindent\textbf{45.} Identify the following decimal number as rational or irrational: $1.709975947\ldots$

\bigskip
\noindent\textbf{46.} Represent $\sqrt{5}$ on a number line.

\bigskip
\noindent\textbf{47.} Express the following recurring decimal as the rational number $\dfrac{p}{q}$, where $p$ and $q$ are integers: $0.\overline{5}$

\bigskip
\noindent\textbf{48.} Express the following recurring decimal as the rational number $\dfrac{p}{q}$, where $p$ and $q$ are integers: $0.\overline{93}$

\bigskip
\noindent\textbf{49.} Insert two rational numbers between $2$ and $3$.

\bigskip
\noindent\textbf{50.} Identify the property that justifies the statement: If $a>13$ then $a+2>15$.

\bigskip
\noindent\textbf{51.} Identify the property that justifies the statement: If $3<9$ and $6<12$ then $9<21$.

\bigskip
\noindent\textbf{52.} Identify the property that justifies the statement: If $7>4$ and $5>3$ then $35>12$.

\bigskip
\noindent\textbf{53.} Identify the property that justifies the statement: If $-5<-4 \Rightarrow 20>16$.

\bigskip
\noindent\textbf{54.} Rationalize the denominator of the following: $\dfrac{13}{4+\sqrt{3}}$

\bigskip
\noindent\textbf{55.} Rationalize the denominator of the following: $\dfrac{\sqrt{2}+\sqrt{5}}{\sqrt{3}}$

\bigskip
\noindent\textbf{56.} Rationalize the denominator of the following: $\dfrac{\sqrt{2}-1}{\sqrt{5}}$

\bigskip
\noindent\textbf{57.} Rationalize the denominator of the following: $\dfrac{6-4\sqrt{2}}{6+4\sqrt{2}}$

\bigskip
\noindent\textbf{58.} Rationalize the denominator of the following: $\dfrac{\sqrt{3}-\sqrt{2}}{\sqrt{3}+\sqrt{2}}$

\bigskip
\noindent\textbf{59.} Rationalize the denominator of the following: $\dfrac{4\sqrt{3}}{\sqrt{7}+\sqrt{5}}$

\bigskip
\noindent\textbf{60.} Simplify the following: $\left(\dfrac{81}{16}\right)^{-3/4}$

\bigskip
\noindent\textbf{61.} Simplify the following: $\left(\dfrac{3}{4}\right)^{-2} \div \left(\dfrac{4}{9}\right)^{3} \times \dfrac{16}{27}$

\bigskip
\noindent\textbf{62.} Simplify the following: $(0.027)^{-1/3}$

\bigskip
\noindent\textbf{63.} Simplify the following: $\sqrt[7]{\dfrac{x^{14}\times y^{21}\times z^{35}}{y^{14}\times z^{7}}}$

\bigskip
\noindent\textbf{64.} Simplify the following: $(64)^{-2/3} \div (9)^{-3/2}$

\bigskip
\noindent\textbf{65.} If $x=3+\sqrt{8}$, then find the value of $x+\dfrac{1}{x}$.

\bigskip
\noindent\textbf{66.} If $x=3+\sqrt{8}$, then find the value of $x-\dfrac{1}{x}$.

\bigskip
\noindent\textbf{67.} If $x=3+\sqrt{8}$, then find the value of $x^2+\dfrac{1}{x^2}$.

\bigskip
\noindent\textbf{68.} If $x=3+\sqrt{8}$, then find the value of $x^2-\dfrac{1}{x^2}$.

\bigskip
\noindent\textbf{69.} If $x=3+\sqrt{8}$, then find the value of $\left(x-\dfrac{1}{x}\right)^2$.

\bigskip
\noindent\textbf{70.} Define radicand.

\bigskip
\noindent\textbf{71.} Define surds.

\bigskip
\noindent\textbf{72.} Define a monomial surd.

\bigskip
\noindent\textbf{73.} Define binomial surds.

\bigskip
\noindent\textbf{74.} Define conjugate surds.

\bigskip
\noindent\textbf{75.} Is $0$ a rational number? Explain.

\bigskip
\noindent\textbf{76.} State the trichotomy property of real numbers.

\bigskip
\noindent\textbf{77.} Find two rational numbers between $4$ and $5$.

\bigskip
\noindent\textbf{78.} Simplify the following: $\sqrt[4]{16x^4y^8}$

\bigskip
\noindent\textbf{79.} Simplify the following: $\sqrt[3]{27x^6y^9z^3}$

\bigskip
\noindent\textbf{80.} Simplify the following: $(64)^{-4/3}$

\bigskip
\noindent\textbf{81.} Rationalize the denominator of the following: $\dfrac{3}{\sqrt{5}+\sqrt{2}}$

\bigskip
\noindent\textbf{82.} Rationalize the denominator of the following: $\dfrac{3}{\sqrt{5}-\sqrt{3}}$

\bigskip
\noindent\textbf{83.} The sum of three consecutive integers is forty-two; find the three integers.

\bigskip
\noindent\textbf{84.} The sum of two real numbers is $8$, and their difference is $2$. Find the numbers.

\bigskip
\noindent\textbf{85.} Normal human body temperature is $98.6\,^{\circ}\text{F}$. Convert it into the Celsius and Kelvin scales.

\bigskip
\noindent\textbf{86.} Express the following number in scientific notation: $2{,}000{,}000$

\bigskip
\noindent\textbf{87.} Express the following number in scientific notation: $48{,}900$

\bigskip
\noindent\textbf{88.} Express the following number in scientific notation: $0.0042$

\bigskip
\noindent\textbf{89.} Express the following number in scientific notation: $0.0000009$

\bigskip
\noindent\textbf{90.} Express the following number in scientific notation: $73\times10^3$

\bigskip
\noindent\textbf{91.} Express the following number in scientific notation: $0.65\times10^2$

\bigskip
\noindent\textbf{92.} Express the following number in ordinary notation: $8.04\times10^2$

\bigskip
\noindent\textbf{93.} Express the following number in ordinary notation: $3\times10^5$

\bigskip
\noindent\textbf{94.} Express the following number in ordinary notation: $1.5\times10^{-2}$

\bigskip
\noindent\textbf{95.} Express the following number in ordinary notation: $1.77\times10^{7}$

\bigskip
\noindent\textbf{96.} Express the following number in ordinary notation: $5.5\times10^{-6}$

\bigskip
\noindent\textbf{97.} Express the following number in ordinary notation: $4\times10^{-5}$

\bigskip
\noindent\textbf{98.} The speed of light is approximately $3\times10^{8}$ metres per second. Express it in standard form.

\bigskip
\noindent\textbf{99.} The circumference of the Earth at the equator is about 40075000 metres. Express this number in scientific notation.

\bigskip
\noindent\textbf{100.} The diameter of Mars is $6.7779\times10^{3}$ km. Express this number in standard form.

\bigskip
\noindent\textbf{101.} The diameter of Earth is $1.2756\times10^{4}$ km. Express this number in standard form.

\bigskip
\noindent\textbf{102.} Convert $5.63\times10^{3}$ to ordinary notation.

\bigskip
\noindent\textbf{103.} Convert $6.6\times10^{-5}$ to ordinary notation.

\bigskip
\noindent\textbf{104.} Convert 29,000,000 to scientific notation.

\bigskip
\noindent\textbf{105.} Convert 0.000006 to scientific notation.

\bigskip
\noindent\textbf{106.} Define scientific notation.

\bigskip
\noindent\textbf{107.} Where is scientific notation used?

\bigskip
\noindent\textbf{108.} Convert 78,000,000 to scientific notation.

\bigskip
\noindent\textbf{109.} Convert 0.0000000315 to scientific notation.

\bigskip
\noindent\textbf{110.} Convert $3.47\times10^{6}$ to ordinary notation.

\bigskip
\noindent\textbf{111.} Convert $6.23\times10^{-4}$ to ordinary notation.

\bigskip
\noindent\textbf{112.} Express the following in logarithmic form: $10^3=1000$

\bigskip
\noindent\textbf{113.} Express the following in logarithmic form: $2^8=256$

\bigskip
\noindent\textbf{114.} Express the following in logarithmic form: $3^{-3}=\dfrac{1}{27}$

\bigskip
\noindent\textbf{115.} Express the following in logarithmic form: $20^2=400$

\bigskip
\noindent\textbf{116.} Express the following in logarithmic form: $16^{-\frac{1}{4}}=\dfrac{1}{2}$

\bigskip
\noindent\textbf{117.} Express the following in logarithmic form: $11^2=121$

\bigskip
\noindent\textbf{118.} Express the following in logarithmic form: $p=q^r$

\bigskip
\noindent\textbf{119.} Express the following in logarithmic form: $(32)^{-\frac{1}{5}}=\dfrac{1}{2}$

\bigskip
\noindent\textbf{120.} Express the following in exponential form: $\log_5 125=3$

\bigskip
\noindent\textbf{121.} Express the following in exponential form: $\log_2 16=4$

\bigskip
\noindent\textbf{122.} Express the following in exponential form: $\log_{23}1=0$

\bigskip
\noindent\textbf{123.} Express the following in exponential form: $\log_5 5=1$

\bigskip
\noindent\textbf{124.} Express the following in exponential form: $\log_2 \dfrac{1}{8}=-3$

\bigskip
\noindent\textbf{125.} Express the following in exponential form: $\dfrac{1}{2}=\log_9 3$

\bigskip
\noindent\textbf{126.} Express the following in exponential form: $5=\log_{10}100000$

\bigskip
\noindent\textbf{127.} Express the following in exponential form: $\log_4 \dfrac{1}{16}=-2$

\bigskip
\noindent\textbf{128.} Find the value of $x$ in the following: $\log_x 64=3$

\bigskip
\noindent\textbf{129.} Find the value of $x$ in the following: $\log_5 1= x$

\bigskip
\noindent\textbf{130.} Find the value of $x$ in the following: $\log_x 8=1$

\bigskip
\noindent\textbf{131.} Find the value of $x$ in the following: $\log_{10} x=-3$

\bigskip
\noindent\textbf{132.} Define logarithm.

\bigskip
\noindent\textbf{133.} What do you know about the logarithm of a real number?

\bigskip
\noindent\textbf{134.} Convert $\log_2 8=3$ to exponential form.

\bigskip
\noindent\textbf{135.} Convert $\log_{10} 100=2$ to exponential form.

\bigskip
\noindent\textbf{136.} Find the value of $x$ in the following case: $\log_5 25=x$

\bigskip
\noindent\textbf{137.} Find the value of $x$ in the following case: $\log_2 x=6$

\bigskip
\noindent\textbf{138.} Find the value of $x$ in the following case: $\log_5 25=x$

\bigskip
\noindent\textbf{139.} Convert the following into logarithmic form: $3^4=81$

\bigskip
\noindent\textbf{140.} Convert the following into logarithmic form: $7^0=1$

\bigskip
\noindent\textbf{141.} Find the characteristic of the following number: 5287

\bigskip
\noindent\textbf{142.} Find the characteristic of the following number: 59.28

\bigskip
\noindent\textbf{143.} Find the characteristic of the following number: 0.0567

\bigskip
\noindent\textbf{144.} Find the characteristic of the following number: 234.7

\bigskip
\noindent\textbf{145.} Find the characteristic of the following number: 0.000049

\bigskip
\noindent\textbf{146.} Find the characteristic of the following number: 145000

\bigskip
\noindent\textbf{147.} Find the logarithm of the following number: 43

\bigskip
\noindent\textbf{148.} Find the logarithm of the following number: 579

\bigskip
\noindent\textbf{149.} Find the logarithm of the following number: 1.982

\bigskip
\noindent\textbf{150.} Find the logarithm of the following number: 0.0876

\bigskip
\noindent\textbf{151.} Find the logarithm of the following number: 0.047

\bigskip
\noindent\textbf{152.} Find the logarithm of the following number: 0.000354

\bigskip
\noindent\textbf{153.} Find the value of $x$: $\log x = 0.0065$

\bigskip
\noindent\textbf{154.} Find the value of $x$: $\log x = 1.192$

\bigskip
\noindent\textbf{155.} Define common logarithm.

\bigskip
\noindent\textbf{156.} Define natural logarithm.

\bigskip
\noindent\textbf{157.} Define characteristic.

\bigskip
\noindent\textbf{158.} Define mantissa.

\bigskip
\noindent\textbf{159.} Write the rules for finding the characteristic.

\bigskip
\noindent\textbf{160.} Define anti-logarithm.

\bigskip
\noindent\textbf{161.} Write the difference between common and natural logarithm.

\bigskip
\noindent\textbf{162.} Without using a calculator, evaluate the following: $\log_2 18-\log_2 9$

\bigskip
\noindent\textbf{163.} Without using a calculator, evaluate the following: $\log_2 64+\log_2 2$

\bigskip
\noindent\textbf{164.} Without using a calculator, evaluate the following: $2\log 2+\log 25$

\bigskip
\noindent\textbf{165.} Write the following as a single logarithm: $\dfrac{1}{2}\log 25 + 2\log 3$

\bigskip
\noindent\textbf{166.} Write the following as a single logarithm: $\log 9 - \log \dfrac{1}{3}$

\bigskip
\noindent\textbf{167.} Write the following as a single logarithm: $\log_5 b^2 . \log_a 5^3$

\bigskip
\noindent\textbf{168.} Write the following as a single logarithm: $2\log_3 x+\log_3 y$

\bigskip
\noindent\textbf{169.} Write the following as a single logarithm: $4\log_5 x -\log_5 y+\log_5 z$

\bigskip
\noindent\textbf{170.} Write the following as a single logarithm: $3\ln a+3\ln b - \ln c$

\bigskip
\noindent\textbf{171.} Expand the following using laws of logarithm: $\log\left(\dfrac{11}{5}\right)$

\bigskip
\noindent\textbf{172.} Expand the following using laws of logarithm: $\log_5 \sqrt{8a^6}$

\bigskip
\noindent\textbf{173.} Expand the following using laws of logarithm: $\ln\left(\dfrac{a^2b}{c}\right)$

\bigskip
\noindent\textbf{174.} Expand the following using laws of logarithm: $\log\left(\dfrac{xy}{z}\right)^{\frac{1}{9}}$

\bigskip
\noindent\textbf{175.} Expand the following using laws of logarithm: $\ln \sqrt[3]{16x^3}$

\bigskip
\noindent\textbf{176.} Expand the following using laws of logarithm: $\log_2\left(\dfrac{1-a}{b}\right)^5$

\bigskip
\noindent\textbf{177.} Find the value of $x$ in the following equation: $\log 2 + \log x =1$

\bigskip
\noindent\textbf{178.} Find the value of $x$ in the following equation: $\log_2 x+\log_2 8=5$

\bigskip
\noindent\textbf{179.} Find the value of $x$ in the following equation: $\log(5x-10)=2$

\bigskip
\noindent\textbf{180.} Write the laws of logarithm.

\bigskip
\noindent\textbf{181.} What is the change of base law?

\bigskip
\noindent\textbf{182.} Prove the change of base law.

\bigskip
\noindent\textbf{183.} Prove the product law.

\bigskip
\noindent\textbf{184.} Prove the quotient law.

\bigskip
\noindent\textbf{185.} Prove the power law.

\bigskip
\noindent\textbf{186.} Expand the following using laws of logarithms: $\log_3(20)$

\bigskip
\noindent\textbf{187.} Expand the following using laws of logarithms: $\log_2(9)^5$

\bigskip
\noindent\textbf{188.} Expand the following using laws of logarithms: $\log_3\left(2^{27}\right)$

\bigskip
\noindent\textbf{189.} Expand the following using laws of logarithms: $\log_2\left(\dfrac{x-y}{z}\right)^3$

\bigskip
\noindent\textbf{190.} Expand the following using laws of logarithms: $\log_5\left(\dfrac{xy}{z}\right)^8$

\bigskip
\noindent\textbf{191.} Write the following as a single logarithm: $2\log_3 10 -\log_3 4$

\bigskip
\noindent\textbf{192.} Write the following as a single logarithm: $6\log_3 x + 2\log_3 11$

\bigskip
\noindent\textbf{193.} Express the following number in scientific notation: 0.000567

\bigskip
\noindent\textbf{194.} Express the following number in scientific notation: 734

\bigskip
\noindent\textbf{195.} Express the following number in scientific notation: $0.33\times10^{3}$

\bigskip
\noindent\textbf{196.} Express the following number in ordinary notation: $2.6\times10^{3}$

\bigskip
\noindent\textbf{197.} Express the following number in ordinary notation: $8.794\times10^{-4}$

\bigskip
\noindent\textbf{198.} Express the following number in ordinary notation: $6\times10^{-6}$

\bigskip
\noindent\textbf{199.} Express the following in logarithmic form: $3^7=2187$

\bigskip
\noindent\textbf{200.} Express the following in logarithmic form: $a^b=c$

\bigskip
\noindent\textbf{201.} Express the following in logarithmic form: $(12)^2=144$

\bigskip
\noindent\textbf{202.} Express the following in exponential form: $\log_4 8=x$

\bigskip
\noindent\textbf{203.} Express the following in exponential form: $\log_9 729=3$

\bigskip
\noindent\textbf{204.} Express the following in exponential form: $\log_4 1024=5$

\bigskip
\noindent\textbf{205.} Find the value of $x$ in the following: $\log_9 x=0.5$

\bigskip
\noindent\textbf{206.} Find the value of $x$ in the following: $\left(\dfrac{1}{9}\right)^{3x}=27$

\bigskip
\noindent\textbf{207.} Write the following as a single logarithm: $7\log x - 3\log y^2$

\bigskip
\noindent\textbf{208.} Write the following as a single logarithm: $3\log 4 - \log 32$

\bigskip
\noindent\textbf{209.} Write the following as a single logarithm: $\dfrac{1}{3}(\log_5 8 + \log_5 27) -\log_5 3$

\bigskip
\noindent\textbf{210.} Expand the following using laws of logarithms: $\log(xyz^6)$

\bigskip
\noindent\textbf{211.} Expand the following using laws of logarithms: $\log_3 \sqrt[6]{m^5n^3}$

\bigskip
\noindent\textbf{212.} Expand the following using laws of logarithms: $\log\sqrt{8x^3}$

\bigskip
\noindent\textbf{213.} Find the value of the following with the help of the logarithm table: $\sqrt[3]{68.24}$

\bigskip
\noindent\textbf{214.} Write the following set in set builder notation: $\{1,4,9,16,25,36,\ldots,484\}$

\bigskip
\noindent\textbf{215.} Write the following set in set builder notation: $\{2,4,8,16,32,64,\ldots,156\}$

\bigskip
\noindent\textbf{216.} Write the following set in set builder notation: $\{0,\pm1,\pm2,\ldots,\pm1000\}$

\bigskip
\noindent\textbf{217.} Write the following set in set builder notation: $\{6,12,18,\ldots,120\}$

\bigskip
\noindent\textbf{218.} Write the following set in set builder notation: $\{100,102,104,\ldots,400\}$

\bigskip
\noindent\textbf{219.} Write the following set in set builder notation: $\{1,3,9,27,81,\ldots\}$

\bigskip
\noindent\textbf{220.} Write the following set in set builder notation: $\{1,2,4,5,10,20,25,50,100\}$

\bigskip
\noindent\textbf{221.} Write the following set in set builder notation: $\{5,10,15,\ldots,100\}$

\bigskip
\noindent\textbf{222.} Write the following set in set builder notation: the set of all integers between $-100$ and $1000$.

\bigskip
\noindent\textbf{223.} Write the following set in tabular form: $\{x \mid x \text{ is a multiple of } 3 \land x\leq35\}$

\bigskip
\noindent\textbf{224.} Write the following set in tabular form: $\{x \mid x\in R \land 2x+1=0\}$

\bigskip
\noindent\textbf{225.} Write the following set in tabular form: $\{x \mid x\in P \land x <12\}$

\bigskip
\noindent\textbf{226.} Write the following set in tabular form: $\{x \mid x \text{ is a divisor of } 128\}$

\bigskip
\noindent\textbf{227.} Write the following set in tabular form: $\{x \mid x=2^n, n\in N \land n <8\}$

\bigskip
\noindent\textbf{228.} Write the following set in tabular form: $\{x \mid x\in N \land x+4=0\}$

\bigskip
\noindent\textbf{229.} Write the following set in tabular form: $\{x \mid x\in N \land x = x\}$

\bigskip
\noindent\textbf{230.} Write the following set in tabular form: $\{x \mid x\in Z \land 3x+1=0\}$

\bigskip
\noindent\textbf{231.} Write two proper subsets of the following set: $\{a,b,c\}$

\bigskip
\noindent\textbf{232.} Write two proper subsets of the following set: $\{0,1\}$

\bigskip
\noindent\textbf{233.} Write two proper subsets of the following set: $N$

\bigskip
\noindent\textbf{234.} Write two proper subsets of the following set: $Z$

\bigskip
\noindent\textbf{235.} Write two proper subsets of the following set: $Q$

\bigskip
\noindent\textbf{236.} Write two proper subsets of the following set: $R$

\bigskip
\noindent\textbf{237.} Write two proper subsets of the following set: $\{x \mid x\in Q \land 0<x\leq2\}$

\bigskip
\noindent\textbf{238.} Is there any set which has no proper subset? If so, name that set.

\bigskip
\noindent\textbf{239.} What is the difference between $\{a,b\}$ and $\{\{a,b\}\}$?

\bigskip
\noindent\textbf{240.} What is the number of elements of the power set of the following set: $\{\ \}$?

\bigskip
\noindent\textbf{241.} What is the number of elements of the power set of the following set? $\{0,1\}$

\bigskip
\noindent\textbf{242.} What is the number of elements of the power set of the following set? $\{1,2,3,4,5,6,7\}$

\bigskip
\noindent\textbf{243.} What is the number of elements of the power set of the following set? $\{0,1,2,3,4,5,6,7\}$

\bigskip
\noindent\textbf{244.} What is the number of elements of the power set of the following set? $\{a,\{b,c\}\}$

\bigskip
\noindent\textbf{245.} What is the number of elements of the power set of the following set? $\{\{a,b\},\{b,c\},\{d,e\}\}$

\bigskip
\noindent\textbf{246.} Write down the power set of the following set: $\{9,11\}$

\bigskip
\noindent\textbf{247.} Write down the power set of the following set: $\{+,-,\times,\div\}$

\bigskip
\noindent\textbf{248.} Write down the power set of the following set: $\{\phi\}$

\bigskip
\noindent\textbf{249.} Write down the power set of the following set: $\{a,\{b,c\}\}$

\bigskip
\noindent\textbf{250.} Define a set.

\bigskip
\noindent\textbf{251.} What is the descriptive method?

\bigskip
\noindent\textbf{252.} What is the tabular method?

\bigskip
\noindent\textbf{253.} What is the set-builder method?

\bigskip
\noindent\textbf{254.} What is a singleton set?

\bigskip
\noindent\textbf{255.} Define null set.

\bigskip
\noindent\textbf{256.} How is an empty set written?

\bigskip
\noindent\textbf{257.} Define equal set.

\bigskip
\noindent\textbf{258.} Define equivalent set.

\bigskip
\noindent\textbf{259.} Define subset.

\bigskip
\noindent\textbf{260.} Define proper subset.

\bigskip
\noindent\textbf{261.} Define improper subset.

\bigskip
\noindent\textbf{262.} Define universal set.

\bigskip
\noindent\textbf{263.} Define power set.

\bigskip
\noindent\textbf{264.} Consider the universal set $U=\{x: x \text{ is a multiple of } 2 \text{ and } 0<x\le 30\}$, $A=\{x: x \text{ is a multiple of } 6\}$ and $B=\{x: x \text{ is a multiple of } 8\}$. List all elements of sets $A$ and $B$ in tabular form.

\bigskip
\noindent\textbf{265.} Consider the universal set $U=\{x: x \text{ is a multiple of } 2 \text{ and } 0<x\le 30\}$, $A=\{x: x \text{ is a multiple of } 6\}$ and $B=\{x: x \text{ is a multiple of } 8\}$. Find $A\cap B$.

\bigskip
\noindent\textbf{266.} If $U=\{1,2,3,\ldots,20\}$ and $A=\{1,3,5,\ldots,19\}$, verify the following: $A\cap U=A$.

\bigskip
\noindent\textbf{267.} What is the union of two sets?

\bigskip
\noindent\textbf{268.} What is the intersection of two sets?

\bigskip
\noindent\textbf{269.} Define disjoint sets.

\bigskip
\noindent\textbf{270.} Define overlapping sets.

\bigskip
\noindent\textbf{271.} What is the complement of a set?

\bigskip
\noindent\textbf{272.} What is a Venn diagram?

\bigskip
\noindent\textbf{273.} Who used the Venn diagram first?

\bigskip
\noindent\textbf{274.} What is the inclusion-exclusion principle?

\bigskip
\noindent\textbf{275.} Define cardinality of a set.

\bigskip
\noindent\textbf{276.} For $A=\{1,2,3,4\}$, find the following relation in $A$. State the domain and range of the following relation: $\{(x,y)\mid y=x\}$

\bigskip
\noindent\textbf{277.} For $A=\{1,2,3,4\}$, find the following relation in $A$. State the domain and range of the following relation: $\{(x,y)\mid y+x=5\}$

\bigskip
\noindent\textbf{278.} For $A=\{1,2,3,4\}$, find the following relation in $A$. State the domain and range of the following relation: $\{(x,y)\mid x+y<5\}$

\bigskip
\noindent\textbf{279.} For $A=\{1,2,3,4\}$, find the following relation in $A$. State the domain and range of the following relation: $\{(x,y)\mid x+y>5\}$

\bigskip
\noindent\textbf{280.} Does the given diagram represent a function? If so, of which type? \emph{[Mapping diagram 1 -- see original paper]}

\bigskip
\noindent\textbf{281.} Does the given diagram represent a function? If so, of which type? \emph{[Mapping diagram 2 -- see original paper]}

\bigskip
\noindent\textbf{282.} Does the given diagram represent a function? If so, of which type? \emph{[Mapping diagram 3 -- see original paper]}

\bigskip
\noindent\textbf{283.} Does the given diagram represent a function? If so, of which type? \emph{[Mapping diagram 4 -- see original paper]}

\bigskip
\noindent\textbf{284.} If $g(x)=3x+2$, find: $g(0)$

\bigskip
\noindent\textbf{285.} If $g(x)=3x+2$, find: $g(-3)$

\bigskip
\noindent\textbf{286.} If $g(x)=3x+2$, find: $g\!\left(\dfrac{2}{3}\right)$

\bigskip
\noindent\textbf{287.} If $h(x)=x^2+1$, then find: $h(1)$

\bigskip
\noindent\textbf{288.} If $h(x)=x^2+1$, then find: $h(-4)$

\bigskip
\noindent\textbf{289.} If $h(x)=x^2+1$, then find: $h\!\left(-\dfrac{1}{2}\right)$

\bigskip
\noindent\textbf{290.} Consider the function defined by $f(x)=5x+1$. If $f(x)=32$, find the value of $x$.

\bigskip
\noindent\textbf{291.} Define functions.

\bigskip
\noindent\textbf{292.} How many types of functions are there? Write their names.

\bigskip
\noindent\textbf{293.} What is an into function?

\bigskip
\noindent\textbf{294.} What is an onto function?

\bigskip
\noindent\textbf{295.} What is a one-one function?

\bigskip
\noindent\textbf{296.} What is a bijective function?

\bigskip
\noindent\textbf{297.} Define linear function.

\bigskip
\noindent\textbf{298.} Define quadratic function.

\bigskip
\noindent\textbf{299.} What is a binary relation?

\bigskip
\noindent\textbf{300.} Define domain.

\bigskip
\noindent\textbf{301.} Define range.

\bigskip
\noindent\textbf{302.} Write the following set in tabular form: $\{x \mid x=2n,\ n\in N\}$

\bigskip
\noindent\textbf{303.} Write the following set in tabular form: $\{x \mid x=2m+1,\ m\in N\}$

\bigskip
\noindent\textbf{304.} Write the following set in tabular form: $\{x \mid x=11n,\ n\in W \wedge n<11\}$

\bigskip
\noindent\textbf{305.} Write the following set in tabular form: $\{x \mid x\in E \wedge 4<x<6\}$

\bigskip
\noindent\textbf{306.} Write the following set in tabular form: $\{x \mid x\in O \wedge 5\le x<7\}$

\bigskip
\noindent\textbf{307.} Write the following set in tabular form: $\{x \mid x\in Q \wedge x^2=2\}$

\bigskip
\noindent\textbf{308.} Write the following set in tabular form: $\{x \mid x\in Q \wedge x=-x\}$

\bigskip
\noindent\textbf{309.} Write the following set in tabular form: $\{x \mid x\in R \wedge x\notin Q'\}$

\bigskip
\noindent\textbf{310.} Let $U=\{1,2,3,4,5,6,7,8,9,10\}$, $A=\{2,4,6,8,10\}$, $B=\{1,2,3,4,5\}$ and $C=\{1,3,5,7,9\}$. List the members of the following set: $A'$

\bigskip
\noindent\textbf{311.} Let $U=\{1,2,3,4,5,6,7,8,9,10\}$, $A=\{2,4,6,8,10\}$, $B=\{1,2,3,4,5\}$ and $C=\{1,3,5,7,9\}$. List the members of the following set: $B'$

\bigskip
\noindent\textbf{312.} Let $U=\{1,2,3,4,5,6,7,8,9,10\}$, $A=\{2,4,6,8,10\}$, $B=\{1,2,3,4,5\}$ and $C=\{1,3,5,7,9\}$. List the members of the following set: $A\cup B$

\bigskip
\noindent\textbf{313.} Let $U=\{1,2,3,4,5,6,7,8,9,10\}$, $A=\{2,4,6,8,10\}$, $B=\{1,2,3,4,5\}$ and $C=\{1,3,5,7,9\}$. List the members of the following set: $A-B$

\bigskip
\noindent\textbf{314.} Let $U=\{1,2,3,4,5,6,7,8,9,10\}$, $A=\{2,4,6,8,10\}$, $B=\{1,2,3,4,5\}$ and $C=\{1,3,5,7,9\}$. List the members of the following set: $A\cap C$

\bigskip
\noindent\textbf{315.} Let $U=\{1,2,3,4,5,6,7,8,9,10\}$, $A=\{2,4,6,8,10\}$, $B=\{1,2,3,4,5\}$ and $C=\{1,3,5,7,9\}$. List the members of the following set: $U'$

\bigskip
\noindent\textbf{316.} Using Venn diagrams, if necessary, find the single set equal to the following: $A'$

\bigskip
\noindent\textbf{317.} Using Venn diagrams, if necessary, find the single set equal to the following: $A\cap U$

\bigskip
\noindent\textbf{318.} Using Venn diagrams, if necessary, find the single set equal to the following: $A\cup U$

\bigskip
\noindent\textbf{319.} Using Venn diagrams, if necessary, find the single set equal to the following: $A\cup \phi$

\bigskip
\noindent\textbf{320.} Using Venn diagrams, if necessary, find the single set equal to the following: $\phi\cup\phi$

\bigskip
\noindent\textbf{321.} Consider the set $P=\{x \mid x=5m,\ m\in N\}$ and $Q=\{x \mid x=2m,\ m\in N\}$. Find $P\cap Q$.

\bigskip
\noindent\textbf{322.} If $g(x)=7x-2$, find: $g(0)$

\bigskip
\noindent\textbf{323.} If $g(x)=7x-2$, then find: $g(-1)$

\bigskip
\noindent\textbf{324.} If $g(x)=7x-2$, then find: $g\!\left(-\dfrac{5}{3}\right)$

\bigskip
\noindent\textbf{325.} If $s(x)=8x^2-3$, find: $s(1)$

\bigskip
\noindent\textbf{326.} If $s(x)=8x^2-3$, find: $s(-9)$

\bigskip
\noindent\textbf{327.} If $s(x)=8x^2-3$, find: $s\!\left(\dfrac{7}{2}\right)$

\bigskip
\noindent\textbf{328.} Consider the function defined by $k(x)=7x-5$. If $k(x)=100$, find the value of $x$.

\bigskip
\noindent\textbf{329.} Factorize by identifying common factors: $6x+12$

\bigskip
\noindent\textbf{330.} Factorize by identifying common factors: $15y^2+20y$

\bigskip
\noindent\textbf{331.} Factorize by identifying common factors: $-12x^2-3x$

\bigskip
\noindent\textbf{332.} Factorize by identifying common factors: $4a^2b+8ab^2$

\bigskip
\noindent\textbf{333.} Factorize by identifying common factors: $xy-3x^2+2x$

\bigskip
\noindent\textbf{334.} Factorize by identifying common factors: $3a^2b-9ab^2+15ab$

\bigskip
\noindent\textbf{335.} Factorize and represent pictorially: $5x+15$

\bigskip
\noindent\textbf{336.} Factorize and represent pictorially: $x^2+4x+3$

\bigskip
\noindent\textbf{337.} Factorize and represent pictorially: $x^2+6x+8$

\bigskip
\noindent\textbf{338.} Factorize and represent pictorially: $x^2+4x+4$

\bigskip
\noindent\textbf{339.} Factorize: $x^2+x-12$

\bigskip
\noindent\textbf{340.} Factorize: $x^2+7x+10$

\bigskip
\noindent\textbf{341.} Factorize: $x^2-6x+8$

\bigskip
\noindent\textbf{342.} Factorize: $x^2-x-56$

\bigskip
\noindent\textbf{343.} Factorize: $x^2-10x-24$

\bigskip
\noindent\textbf{344.} Factorize: $y^2+4y-12$

\bigskip
\noindent\textbf{345.} Factorize: $y^2+13y+36$

\bigskip
\noindent\textbf{346.} Factorize: $x^2-x-2$

\bigskip
\noindent\textbf{347.} Factorize: $2x^2+7x+3$

\bigskip
\noindent\textbf{348.} Factorize: $4x^2+13x+3$

\bigskip
\noindent\textbf{349.} Factorize: $3x^2+5x+2$

\bigskip
\noindent\textbf{350.} Define common factor.

\bigskip
\noindent\textbf{351.} Define quadratic expression.

\bigskip
\noindent\textbf{352.} Find the common factor of $x^2+2x$ concretely, pictorially, and symbolically.

\bigskip
\noindent\textbf{353.} Factorize $x^2-5x+4$ concretely, pictorially, and symbolically.

\bigskip
\noindent\textbf{354.} Factorize $x^2-3x-10$ concretely, pictorially, and symbolically.

\bigskip
\noindent\textbf{355.} Factorize: $x^2+9x+14$

\bigskip
\noindent\textbf{356.} Factorize: $x^2-11x+24$

\bigskip
\noindent\textbf{357.} Factorize: $p^2+11p+18$

\bigskip
\noindent\textbf{358.} Factorize: $3x^2-4x-4$

\bigskip
\noindent\textbf{359.} Factorize: $8x^3+12x^2+6x+1$

\bigskip
\noindent\textbf{360.} Factorize: $27a^3+108a^2b+144ab^2+64b^3$

\bigskip
\noindent\textbf{361.} Factorize: $x^3+48x^2y+108xy^2+216y^3$

\bigskip
\noindent\textbf{362.} Factorize: $8x^3-125y^3+150xy^2-60x^2y$

\bigskip
\noindent\textbf{363.} Factorize: $125a^3-1$

\bigskip
\noindent\textbf{364.} Factorize: $64x^3+125$

\bigskip
\noindent\textbf{365.} Factorize: $x^6-27$

\bigskip
\noindent\textbf{366.} Factorize: $1000a^3+1$

\bigskip
\noindent\textbf{367.} Factorize: $343x^3+216$

\bigskip
\noindent\textbf{368.} Factorize: $27-512y^3$

\bigskip
\noindent\textbf{369.} Factorize: $x^3-18x^2+108x-216$

\bigskip
\noindent\textbf{370.} Factorize: $8x^3+27$

\bigskip
\noindent\textbf{371.} Factorize: $x^3-27y^3$

\bigskip
\noindent\textbf{372.} Find HCF by factorization method: $21x^2y,\ 35xy^2$

\bigskip
\noindent\textbf{373.} Find HCF by factorization method: $4x^2-9y^2,\ 2x^2-3xy$

\bigskip
\noindent\textbf{374.} Find HCF by factorization method: $x^3-1,\ x^2+x+1$

\bigskip
\noindent\textbf{375.} Find the LCM of the following expressions by using the prime factorization method: $2a^2b,\ 4ab^2,\ 6ab$

\bigskip
\noindent\textbf{376.} Find the LCM of the following expressions by using the prime factorization method: $x^2+x,\ x^3+x^2$

\bigskip
\noindent\textbf{377.} Find the LCM of the following expressions by using the prime factorization method: $a^2-4a+4,\ a^2-2a$

\bigskip
\noindent\textbf{378.} Define HCF.

\bigskip
\noindent\textbf{379.} Define LCM.

\bigskip
\noindent\textbf{380.} Write the relationship between HCF and LCM.

\bigskip
\noindent\textbf{381.} Find the HCF of $6x^2y,\ 9xy^2$.

\bigskip
\noindent\textbf{382.} Find the LCM of $4x^2y,\ 8x^3y^2$.

\bigskip
\noindent\textbf{383.} Find the square root of the following polynomial by factorization method: $x^2-8x+16$

\bigskip
\noindent\textbf{384.} Find the square root of the following polynomial by factorization method: $9x^2+12x+4$

\bigskip
\noindent\textbf{385.} Find the square root of the following polynomial by factorization method: $36a^2+84a+49$

\bigskip
\noindent\textbf{386.} Find the square root of the following polynomial by factorization method: $64y^2-32y+4$

\bigskip
\noindent\textbf{387.} Factorize the following: $4x^3+18x^2-12x$

\bigskip
\noindent\textbf{388.} Factorize the following: $x^3+64y^3$

\bigskip
\noindent\textbf{389.} Factorize the following: $-x^2-23x-60$

\bigskip
\noindent\textbf{390.} Factorize the following: $x^3y^3-8$

\bigskip
\noindent\textbf{391.} Factorize the following: $2x^2+7x+3$

\bigskip
\noindent\textbf{392.} Factorize the following: $x^4+64$

\bigskip
\noindent\textbf{393.} Find LCM and HCF by prime factorization method: $x^2+8x+16,\ x^2-16$

\bigskip
\noindent\textbf{394.} Find LCM and HCF by prime factorization method: $x^3-9x,\ x^2-4x+3$

\bigskip
\noindent\textbf{395.} Find the square root by factorization and division method of the expression $16x^4+8x^2+1$.

\bigskip
\noindent\textbf{396.} Factorize the following: $x^4+2x^2+9$

\bigskip
\noindent\textbf{397.} Solve and represent the solution on a real line: $12x+30=-6$

\bigskip
\noindent\textbf{398.} Solve and represent the solution on a real line: $\dfrac{x}{3}+6=-12$

\bigskip
\noindent\textbf{399.} Solve and represent the solution on a real line: $\dfrac{x}{2}-\dfrac{3x}{4}=\dfrac{1}{12}$

\bigskip
\noindent\textbf{400.} Solve and represent the solution on a real line: $2=7(2x+4)+12x$

\bigskip
\noindent\textbf{401.} Solve the following inequality and represent the solution on a real line: $x-6\leq-2$

\bigskip
\noindent\textbf{402.} Solve the following inequality and represent the solution on a real line: $-9>-16+x$

\bigskip
\noindent\textbf{403.} Solve the following inequality and represent the solution on a real line: $3+2x\geq3$

\bigskip
\noindent\textbf{404.} Solve the following inequality and represent the solution on a real line: $6(x+10)\leq0$

\bigskip
\noindent\textbf{405.} Shade the solution region for the following linear inequality in the $xy$-plane: $2x+y\leq6$

\bigskip
\noindent\textbf{406.} Shade the solution region for the following linear inequality in the $xy$-plane: $3x+7y\geq21$

\bigskip
\noindent\textbf{407.} Shade the solution region for the following linear inequality in the $xy$-plane: $3x-2y\geq6$

\bigskip
\noindent\textbf{408.} Shade the solution region for the following linear inequality in the $xy$-plane: $5x-4y\leq20$

\bigskip
\noindent\textbf{409.} Shade the solution region for the following linear inequality in the $xy$-plane: $2x+1\geq0$

\bigskip
\noindent\textbf{410.} Shade the solution region for the following linear inequality in the $xy$-plane: $3y-4\leq0$

\bigskip
\noindent\textbf{411.} Define linear equation.

\bigskip
\noindent\textbf{412.} How is a linear equation in one variable solved?

\bigskip
\noindent\textbf{413.} Solve the following equation and represent the solution on a real line: $3x-5=7$

\bigskip
\noindent\textbf{414.} Solve the following equation and represent the solution on a real line: $\dfrac{x-2}{5}-\dfrac{x-4}{2}=2$

\bigskip
\noindent\textbf{415.} Find the solution of $\dfrac{2}{3}x-1<0$ and also represent it on a real line.

\bigskip
\noindent\textbf{416.} Solve the inequality $x+2y<6$.

\bigskip
\noindent\textbf{417.} Solve the following linear inequality in the $xy$-plane: $2x\geq-3$

\bigskip
\noindent\textbf{418.} Solve the following linear inequality in the $xy$-plane: $y\leq2$

\bigskip
\noindent\textbf{419.} Define problem constraints.

\bigskip
\noindent\textbf{420.} Define decision variables.

\bigskip
\noindent\textbf{421.} Define feasible region.

\bigskip
\noindent\textbf{422.} Define feasible solution.

\bigskip
\noindent\textbf{423.} Define non-negative constraints.

\bigskip
\noindent\textbf{424.} Define corner point.

\bigskip
\noindent\textbf{425.} Define vertex.

\bigskip
\noindent\textbf{426.} Define objective function.

\bigskip
\noindent\textbf{427.} Define optimal solution.

\bigskip
\noindent\textbf{428.} Solve and represent the solution on a real line: $\dfrac{x+5}{3}=1-x$

\bigskip
\noindent\textbf{429.} Solve and represent the solution on a real line: $\dfrac{2x+1}{3}+\dfrac{1}{2}=1-\dfrac{x-1}{3}$

\bigskip
\noindent\textbf{430.} Solve and represent the solution on a real line: $3x+7<16$

\bigskip
\noindent\textbf{431.} Find in which quadrant the following angle lies. Write a co-terminal angle for: $65^{\circ}$

\bigskip
\noindent\textbf{432.} Find in which quadrant the following angle lies. Write a co-terminal angle for: $135^{\circ}$

\bigskip
\noindent\textbf{433.} Find in which quadrant the following angle lies. Write a co-terminal angle for: $-40^{\circ}$

\bigskip
\noindent\textbf{434.} Find in which quadrant the following angle lies. Write a co-terminal angle for: $210^{\circ}$

\bigskip
\noindent\textbf{435.} Find in which quadrant the following angle lies. Write a co-terminal angle for: $-150^{\circ}$

\bigskip
\noindent\textbf{436.} Convert the following into degrees, minutes, and seconds: $123.456^{\circ}$

\bigskip
\noindent\textbf{437.} Convert the following into degrees, minutes, and seconds: $58.7891^{\circ}$

\bigskip
\noindent\textbf{438.} Convert the following into degrees, minutes, and seconds: $90.5678^{\circ}$

\bigskip
\noindent\textbf{439.} Convert the following into decimal degrees: $65^{\circ}32'15''$

\bigskip
\noindent\textbf{440.} Convert the following into decimal degrees: $42^{\circ}18'45''$

\bigskip
\noindent\textbf{441.} Convert the following into decimal degrees: $78^{\circ}45'36''$

\bigskip
\noindent\textbf{442.} Convert the following into radians: $36^{\circ}$

\bigskip
\noindent\textbf{443.} Convert the following into radians: $22.5^{\circ}$

\bigskip
\noindent\textbf{444.} Convert the following into radians: $67.5^{\circ}$

\bigskip
\noindent\textbf{445.} Convert the following into degrees: $\dfrac{\pi}{16}$ rad

\bigskip
\noindent\textbf{446.} Convert the following into degrees: $\dfrac{11\pi}{5}$ rad

\bigskip
\noindent\textbf{447.} Convert the following into degrees: $\dfrac{7\pi}{6}$ rad

\bigskip
\noindent\textbf{448.} Find the arc length and area of a sector if $r=6$ cm and the central angle $\theta=\dfrac{\pi}{3}$ radians.

\bigskip
\noindent\textbf{449.} Find the arc length and area of a sector if $r=\dfrac{4.8}{\pi}$ and the central angle $\theta=\dfrac{5\pi}{6}$.

\bigskip
\noindent\textbf{450.} Find the percentage of the area of a sector subtending an angle of $\dfrac{\pi}{8}$ radians.

\bigskip
\noindent\textbf{451.} What is trigonometry?

\bigskip
\noindent\textbf{452.} When will an angle be positive?

\bigskip
\noindent\textbf{453.} When will an angle be negative?

\bigskip
\noindent\textbf{454.} What are co-terminal angles?

\bigskip
\noindent\textbf{455.} Define degree.

\bigskip
\noindent\textbf{456.} Define radian.

\bigskip
\noindent\textbf{457.} Define full circle.

\bigskip
\noindent\textbf{458.} Define half circle.

\bigskip
\noindent\textbf{459.} Define right angle.

\bigskip
\noindent\textbf{460.} Prove that: $l=r\theta$

\bigskip
\noindent\textbf{461.} Prove that: $A=\dfrac{1}{2}r^2\theta$

\bigskip
\noindent\textbf{462.} Convert $73.12^{\circ}$ to degrees, minutes, and seconds.

\bigskip
\noindent\textbf{463.} Convert $109.42^{\circ}$ to degrees, minutes, and seconds.

\bigskip
\noindent\textbf{464.} Convert $45^{\circ}45'45''$ to decimal degrees.

\bigskip
\noindent\textbf{465.} Convert $94^{\circ}27'54''$ to decimal degrees.

\bigskip
\noindent\textbf{466.} Convert radians to degrees: $\dfrac{5\pi}{3}$ rad

\bigskip
\noindent\textbf{467.} Convert radians to degrees: $\dfrac{7\pi}{6}$ rad

\bigskip
\noindent\textbf{468.} Convert radians to degrees: $\dfrac{11\pi}{6}$ rad

\bigskip
\noindent\textbf{469.} Convert radians to degrees: $1.2$ rad

\bigskip
\noindent\textbf{470.} Convert degrees to radians: $15^{\circ}$

\bigskip
\noindent\textbf{471.} Convert degrees to radians: $75^{\circ}$

\bigskip
\noindent\textbf{472.} Convert degrees to radians: $315^{\circ}$

\bigskip
\noindent\textbf{473.} Convert degrees to radians: $15^{\circ}15'$

\bigskip
\noindent\textbf{474.} Find the arc length of a sector with radius $r=10$ cm and central angle $\theta=60^{\circ}$.

\bigskip
\noindent\textbf{475.} Find the area of a sector with radius $r=8$ cm and angle $\theta=45^{\circ}$.

\bigskip
\noindent\textbf{476.} If the arc length of a sector of radius 5 cm is 11 cm, find the angle subtended by the arc in radians and degrees.

\bigskip
\noindent\textbf{477.} If $\theta$ lies in the first quadrant, find the remaining trigonometric ratios of $\theta$: $\sin\theta=\dfrac{2}{3}$

\bigskip
\noindent\textbf{478.} If $\theta$ lies in the first quadrant, find the remaining trigonometric ratios of $\theta$: $\cos\theta=\dfrac{3}{4}$

\bigskip
\noindent\textbf{479.} If $\theta$ lies in the first quadrant, find the remaining trigonometric ratios of $\theta$: $\tan\theta=\dfrac{1}{2}$

\bigskip
\noindent\textbf{480.} If $\theta$ lies in the first quadrant, find the remaining trigonometric ratios of $\theta$: $\sec\theta=3$

\bigskip
\noindent\textbf{481.} If $\theta$ lies in the first quadrant, find the remaining trigonometric ratios of $\theta$: $\cot\theta=\sqrt{\dfrac{3}{2}}$

\bigskip
\noindent\textbf{482.} Prove the following trigonometric identity: $(\sin\theta+\cos\theta)^2=1+2\sin\theta\cos\theta$

\bigskip
\noindent\textbf{483.} Prove the following trigonometric identity: $\dfrac{\cos\theta}{\sin\theta}=\dfrac{1}{\tan\theta}$

\bigskip
\noindent\textbf{484.} Prove the following trigonometric identity: $\cos^2\theta-\sin^2\theta=2\cos^2\theta-1$

\bigskip
\noindent\textbf{485.} Prove the following trigonometric identity: $\cos^2\theta-\sin^2\theta=1-2\sin^2\theta$

\bigskip
\noindent\textbf{486.} Show that: $(\sec^2\theta-1)\cos^2\theta=\sin^2\theta$

\bigskip
\noindent\textbf{487.} Show that: $\tan\theta\cdot\cos\theta=\sec\theta\cdot\text{cosec}\,\theta$

\bigskip
\noindent\textbf{488.} Find the value of the following trigonometric ratio without using a calculator: $\sin30^{\circ}$

\bigskip
\noindent\textbf{489.} Find the value of the following trigonometric ratio without using a calculator: $\cos30^{\circ}$

\bigskip
\noindent\textbf{490.} Find the value of the following trigonometric ratio without using a calculator: $\tan\dfrac{\pi}{6}$

\bigskip
\noindent\textbf{491.} Find the value of the following trigonometric ratio without using a calculator: $\tan60^{\circ}$

\bigskip
\noindent\textbf{492.} Find the value of the following trigonometric ratio without using a calculator: $\sec60^{\circ}$

\bigskip
\noindent\textbf{493.} Find the value of the following trigonometric ratio without using a calculator: $\cos\dfrac{\pi}{3}$

\bigskip
\noindent\textbf{494.} Find the value of the following trigonometric ratio without using a calculator: $\cot 60^{\circ}$

\bigskip
\noindent\textbf{495.} Find the value of the following trigonometric ratio without using a calculator: $\sin 60^{\circ}$

\bigskip
\noindent\textbf{496.} Find the value of the following trigonometric ratio without using a calculator: $\sec 30^{\circ}$

\bigskip
\noindent\textbf{497.} Find the value of the following trigonometric ratio without using a calculator: $\csc 30^{\circ}$

\bigskip
\noindent\textbf{498.} Find the value of the following trigonometric ratio without using a calculator: $\sin 45^{\circ}$

\bigskip
\noindent\textbf{499.} Find the value of the following trigonometric ratio without using a calculator: $\cos\dfrac{\pi}{4}$

\bigskip
\noindent\textbf{500.} Evaluate: $2\sin 60^{\circ}\cos 60^{\circ}$

\bigskip
\noindent\textbf{501.} Evaluate: $2\cos\dfrac{\pi}{3}\sin\dfrac{\pi}{3}$

\bigskip
\noindent\textbf{502.} Evaluate: $2\sin 45^{\circ}+2\cos 45^{\circ}$

\bigskip
\noindent\textbf{503.} Evaluate: $\cos 60^{\circ}\cos 30^{\circ} - \sin 60^{\circ}\sin 30^{\circ}$

\bigskip
\noindent\textbf{504.} Evaluate: $\sin 60^{\circ}\cos 30^{\circ} - \cos 60^{\circ}\sin 30^{\circ}$

\bigskip
\noindent\textbf{505.} Evaluate: $\cos 60^{\circ}\cos 30^{\circ} + \sin 60^{\circ}\sin 30^{\circ}$

\bigskip
\noindent\textbf{506.} Evaluate: $\tan\dfrac{\pi}{6}\cot\dfrac{\pi}{6}+1$

\bigskip
\noindent\textbf{507.} Find the value of $x$, $y$ and $z$ from the following right-angled triangle ($\angle B = 90^{\circ}$, $\angle A = 30^{\circ}$, $AB=y=4$ cm):
\begin{center}
\begin{tikzpicture}[scale=0.9]
\coordinate (A) at (0,0);
\coordinate (B) at (4,0);
\coordinate (C) at (4,2.3);
\draw (A) -- (B) -- (C) -- cycle;
\draw (3.7,0) rectangle (4,0.3);
\node[left] at (2.2,1.35) {$z$};
\node[right] at (4.05,1.15) {$x$};
\node[below] at (2,-0.15) {$y = 4$ cm};
\node[above left] at (0.55,0.15) {$30^{\circ}$};
\node[below left] at (A) {$A$};
\node[below right] at (B) {$B$};
\node[above right] at (C) {$C$};
\end{tikzpicture}
\end{center}

\bigskip
\noindent\textbf{508.} Find the value of $x$, $y$ and $z$ from the following right-angled triangle ($\angle B = 90^{\circ}$, $\angle A = \dfrac{\pi}{4}$, $BC=y=\sqrt{3}$ cm):
\begin{center}
\begin{tikzpicture}[scale=0.9]
\coordinate (A) at (0,0);
\coordinate (B) at (3,0);
\coordinate (C) at (3,3);
\draw (A) -- (B) -- (C) -- cycle;
\draw (2.7,0) rectangle (3,0.3);
\node[left] at (1.5,1.7) {$z$};
\node[right] at (3.05,1.5) {$y=\sqrt{3}$ cm};
\node[below] at (1.5,-0.15) {$x$};
\node[above left] at (0.55,0.15) {$\pi/4$};
\node[below left] at (A) {$A$};
\node[below right] at (B) {$B$};
\node[above right] at (C) {$C$};
\end{tikzpicture}
\end{center}

\bigskip
\noindent\textbf{509.} Find the value of $x$, $y$ and $z$ from the following right-angled triangle ($\angle B = 90^{\circ}$, $\angle A = 45^{\circ}$, $AB=y=4$ cm):
\begin{center}
\begin{tikzpicture}[scale=0.9]
\coordinate (A) at (0,0);
\coordinate (B) at (4,0);
\coordinate (C) at (4,4);
\draw (A) -- (B) -- (C) -- cycle;
\draw (3.7,0) rectangle (4,0.3);
\node[left] at (2.2,2.3) {$z$};
\node[right] at (4.05,2) {$x$};
\node[below] at (2,-0.15) {$y = 4$ cm};
\node[above left] at (0.55,0.15) {$45^{\circ}$};
\node[below left] at (A) {$A$};
\node[below right] at (B) {$B$};
\node[above right] at (C) {$C$};
\end{tikzpicture}
\end{center}

\bigskip
\noindent\textbf{510.} Find the value of $x$, $y$ and $z$ from the following right-angled triangle ($\angle B = 90^{\circ}$, $\angle C = 60^{\circ}$, $CA=z=2$ cm):
\begin{center}
\begin{tikzpicture}[scale=0.9]
\coordinate (C) at (0,0);
\coordinate (B) at (2.6,0);
\coordinate (A) at (2.6,3.2);
\draw (C) -- (B) -- (A) -- cycle;
\draw (2.3,0) rectangle (2.6,0.3);
\node[above left] at (1.2,1.7) {$z=2$ cm};
\node[right] at (2.65,1.6) {$y$};
\node[below] at (1.3,-0.15) {$x$};
\node[above right] at (0.55,0.15) {$60^{\circ}$};
\node[below left] at (C) {$C$};
\node[below right] at (B) {$B$};
\node[above right] at (A) {$A$};
\end{tikzpicture}
\end{center}

\bigskip
\noindent\textbf{511.} Find the unknown side and angles of the following triangle ($\angle B = 90^{\circ}$, $AB=c=\sqrt{13}$ cm, $BC=a=\sqrt{3}$ cm):
\begin{center}
\begin{tikzpicture}[scale=0.9]
\coordinate (A) at (0,0);
\coordinate (B) at (3.6,0);
\coordinate (C) at (3.6,2);
\draw (A) -- (B) -- (C) -- cycle;
\draw (3.3,0) rectangle (3.6,0.3);
\node[left] at (1.8,1.1) {$b$};
\node[right] at (3.65,1) {$a=\sqrt{3}$ cm};
\node[below] at (1.8,-0.15) {$c=\sqrt{13}$ cm};
\node[below left] at (A) {$A$};
\node[below right] at (B) {$B$};
\node[above right] at (C) {$C$};
\end{tikzpicture}
\end{center}

\bigskip
\noindent\textbf{512.} Find the unknown side and angles of the following triangle ($\angle B = 90^{\circ}$, $AB=c=4$ cm, $BC=a=4$ cm):
\begin{center}
\begin{tikzpicture}[scale=0.9]
\coordinate (A) at (0,0);
\coordinate (B) at (3,0);
\coordinate (C) at (3,3);
\draw (A) -- (B) -- (C) -- cycle;
\draw (2.7,0) rectangle (3,0.3);
\node[left] at (1.5,1.7) {$b$};
\node[right] at (3.05,1.5) {$a=4$ cm};
\node[below] at (1.5,-0.15) {$c=4$ cm};
\node[below left] at (A) {$A$};
\node[below right] at (B) {$B$};
\node[above right] at (C) {$C$};
\end{tikzpicture}
\end{center}

\bigskip
\noindent\textbf{513.} Each side of a square field is $60$ m long. Find the lengths of the diagonals of the field.

\bigskip
\noindent\textbf{514.} Calculate the value of $x$ in the following case ($\angle B=90^{\circ}$, $BC=12$ cm, $AC=20$ cm, $AB=x$):
\begin{center}
\begin{tikzpicture}[scale=0.35]
\coordinate (A) at (0,0);
\coordinate (B) at (0,8);
\coordinate (C) at (12,8);
\draw (A) -- (B) -- (C) -- cycle;
\draw (0,7.4) rectangle (0.6,8);
\node[left] at (-0.2,4) {$x$ cm};
\node[above] at (6,8.3) {$12$ cm};
\node[below right] at (6,3.7) {$20$ cm};
\node[below left] at (A) {$A$};
\node[above left] at (B) {$B$};
\node[above right] at (C) {$C$};
\end{tikzpicture}
\end{center}

\bigskip
\noindent\textbf{515.} Calculate the value of $x$ in the following case ($D$ is the apex, $DB \perp AC$, $AD=x$ cm, $DC=5$ cm, $BC=4$ cm, and $m\overline{AB}=m\overline{BC}$):
\begin{center}
\begin{tikzpicture}[scale=0.55]
\coordinate (A) at (0,0);
\coordinate (B) at (4,0);
\coordinate (C) at (8,0);
\coordinate (D) at (4,5);
\draw (A) -- (C);
\draw (A) -- (D) -- (C);
\draw (D) -- (B);
\draw (3.7,0) rectangle (4,0.3);
\node[above left] at (2,2.5) {$x$ cm};
\node[above right] at (6,2.5) {$5$ cm};
\node[below] at (2,-0.2) {};
\node[below] at (6,-0.2) {$4$ cm};
\node[below] at (A) {$A$};
\node[below] at (B) {$B$};
\node[below] at (C) {$C$};
\node[above] at (D) {$D$};
\end{tikzpicture}
\end{center}

\bigskip
\noindent\textbf{516.} Solve triangle $ABC$, when $m\angle A = 60^{\circ}$, $b=5$ cm, $m\angle B = 90^{\circ}$.

\bigskip
\noindent\textbf{517.} Solve triangle $ABC$, when $a=2$ cm, $b=2\sqrt{2}$ cm and $m\angle B = 90^{\circ}$.

\bigskip
\noindent\textbf{518.} The angle of elevation of the top of a flag post from a point on the ground level $40$ m away from the flag post is $60^{\circ}$. Find the height of the post.

\bigskip
\noindent\textbf{519.} A tree is $72$ m high. Find the angle of elevation of its top from a point $100$ m away on the ground level.

\bigskip
\noindent\textbf{520.} A ladder makes an angle of $60^{\circ}$ with the ground and reaches a height of $10$ m along the wall. Find the length of the ladder.

\bigskip
\noindent\textbf{521.} Find the measure of the angle of elevation of the Sun, if a tower $300$ m high casts a shadow $450$ m long.

\bigskip
\noindent\textbf{522.} Convert the given angle from degrees to radians, giving the answer in terms of $\pi$: $255^{\circ}$

\bigskip
\noindent\textbf{523.} Convert the given angle from degrees to radians, giving the answer in terms of $\pi$: $75^{\circ}45'$

\bigskip
\noindent\textbf{524.} Convert the given angle from degrees to radians, giving the answer in terms of $\pi$: $142.5^{\circ}$

\bigskip
\noindent\textbf{525.} Convert the given angle from radians to degrees, giving the answer in degrees and minutes: $\dfrac{17\pi}{24}$

\bigskip
\noindent\textbf{526.} Convert the given angle from radians to degrees, giving the answer in degrees and minutes: $\dfrac{11\pi}{16}$

\bigskip
\noindent\textbf{527.} Prove the following identity: $\dfrac{\sin\theta}{1-\cos\theta} = \dfrac{1+\cos\theta}{\sin\theta}$

\bigskip
\noindent\textbf{528.} Describe the location in the plane of the point $P(x,y)$, for which: $x>0$

\bigskip
\noindent\textbf{529.} Describe the location in the plane of the point $P(x,y)$, for which: $x>0$ and $y>0$

\bigskip
\noindent\textbf{530.} Describe the location in the plane of the point $P(x,y)$, for which: $x=0$

\bigskip
\noindent\textbf{531.} Describe the location in the plane of the point $P(x,y)$, for which: $y=0$

\bigskip
\noindent\textbf{532.} Describe the location in the plane of the point $P(x,y)$, for which: $x\geq0$ and $y<0$

\bigskip
\noindent\textbf{533.} Describe the location in the plane of the point $P(x,y)$, for which: $y=0$, $x=0$

\bigskip
\noindent\textbf{534.} Describe the location in the plane of the point $P(x,y)$, for which: $|x|=|y|$

\bigskip
\noindent\textbf{535.} Describe the location in the plane of the point $P(x,y)$, for which: $x\geq3$

\bigskip
\noindent\textbf{536.} Describe the location in the plane of the point $P(x,y)$, for which: $y>2$

\bigskip
\noindent\textbf{537.} Describe the location in the plane of the point $P(x,y)$, for which: $x$ and $y$ have opposite signs.

\bigskip
\noindent\textbf{538.} Find the distance between the points: $A(6,7)$, $B(0,-2)$

\bigskip
\noindent\textbf{539.} Find the distance between the points: $C(-5,-2)$, $D(3,2)$

\bigskip
\noindent\textbf{540.} Find the distance between the points: $L(0,3)$, $M(-2,-4)$

\bigskip
\noindent\textbf{541.} Find the distance between the points: $P(-8,-7)$, $Q(0,0)$

\bigskip
\noindent\textbf{542.} Find the distance between the two given points (as given in the original figure/table).

\bigskip
\noindent\textbf{543.} Find the midpoint of the line segment joining the two points: $A(3,1)$, $B(-2,-4)$

\bigskip
\noindent\textbf{544.} Find the midpoint of the line segment joining the two points: $A(-8,3)$, $(2,-1)$

\bigskip
\noindent\textbf{545.} Find the midpoint of the line segment joining the two points: $A\left(-\sqrt{5},-\dfrac{1}{3}\right)$, $B(-3\sqrt{5},\ 5)$

\bigskip
\noindent\textbf{546.} Is the following point at a distance of $15$ units from the origin? $(\sqrt{176},\ 7)$

\bigskip
\noindent\textbf{547.} Is the following point at a distance of $15$ units from the origin? $(10,-10)$

\bigskip
\noindent\textbf{548.} Is the following point at a distance of $15$ units from the origin? $(1,15)$

\bigskip
\noindent\textbf{549.} Write briefly about the history of geometry.

\bigskip
\noindent\textbf{550.} Define origin.

\bigskip
\noindent\textbf{551.} Define coordinate axes.

\bigskip
\noindent\textbf{552.} Define quadrants.

\bigskip
\noindent\textbf{553.} Define a two-dimensional rectangular coordinate system.

\bigskip
\noindent\textbf{554.} Write the distance formula.

\bigskip
\noindent\textbf{555.} Write the midpoint formula.

\bigskip
\noindent\textbf{556.} Find the distance between the points: $A(5,6)$, $B(5,-2)$

\bigskip
\noindent\textbf{557.} Show that the points $A(-1,2)$, $B(7,5)$ and $C(2,-6)$ are vertices of a right triangle.

\bigskip
\noindent\textbf{558.} The point $C(-5,3)$ is the centre of a circle and $P(7,-2)$ lies on the circle. What is the radius of the circle?

\bigskip
\noindent\textbf{559.} Find the midpoint of the line segment joining the points $A(2,3)$ and $B(8,7)$.

\bigskip
\noindent\textbf{560.} Find the slope and inclination of the line joining the points: $(-2,4)$; $(5,11)$

\bigskip
\noindent\textbf{561.} Find the slope and inclination of the line joining the points: $(3,-2)$; $(2,7)$

\bigskip
\noindent\textbf{562.} Find the slope and inclination of the line joining the points: $(4,6)$; $(4,8)$

\bigskip
\noindent\textbf{563.} By means of slope, show that the following points lie on the same line: $A(-1,-3)$; $B(1,5)$; $C(2,9)$

\bigskip
\noindent\textbf{564.} By means of slope, show that the following points lie on the same line: $P(4,-5)$; $Q(7,5)$; $R(10,15)$

\bigskip
\noindent\textbf{565.} By means of slope, show that the following points lie on the same line: $L(-4,6)$; $M(3,8)$; $N(10,10)$

\bigskip
\noindent\textbf{566.} By means of slope, show that the following points lie on the same line: $X(a,2b)$; $Y(c,a+b)$; $Z(2c-a,2a)$

\bigskip
\noindent\textbf{567.} Find $k$ so that the line joining $A(7,3)$; $B(k,-6)$ and the line joining $C(-4,5)$; $D(-6,4)$ are parallel.

\bigskip
\noindent\textbf{568.} Find $k$ so that the line joining $A(7,3)$; $B(k,-6)$ and the line joining $C(-4,5)$; $D(-6,4)$ are perpendicular.

\bigskip
\noindent\textbf{569.} Find an equation of the horizontal line through: $(7,-9)$

\bigskip
\noindent\textbf{570.} Find an equation of the vertical line through: $(-5,3)$

\bigskip
\noindent\textbf{571.} Find an equation of the line through $A(-6,5)$ having slope $7$.

\bigskip
\noindent\textbf{572.} Find an equation of the line through $(8,-3)$ having slope $0$.

\bigskip
\noindent\textbf{573.} Find an equation of the line through $(-8,5)$ having an undefined slope.

\bigskip
\noindent\textbf{574.} Find an equation of the line through $(-5,-3)$ and $(9,-1)$.

\bigskip
\noindent\textbf{575.} Find an equation of the line with $y$-intercept $-7$ and slope $-5$.

\bigskip
\noindent\textbf{576.} Find an equation of the line with $x$-intercept $-3$ and $y$-intercept $4$.

\bigskip
\noindent\textbf{577.} Find an equation of the line with $x$-intercept $-9$ and slope $-4$.

\bigskip
\noindent\textbf{578.} Define slope.

\bigskip
\noindent\textbf{579.} Define $x$-intercept.

\bigskip
\noindent\textbf{580.} Define $y$-intercept.

\bigskip
\noindent\textbf{581.} Show that the points $A(-3,6)$, $B(3,2)$ and $C(6,0)$ are collinear.

\bigskip
\noindent\textbf{582.} Show that the triangle with vertices $A(1,1)$, $B(4,5)$ and $C(12,-1)$ is a right triangle.

\bigskip
\noindent\textbf{583.} Find an equation of the line through the points $(-2,1)$ and $(6,-4)$.

\bigskip
\noindent\textbf{584.} Write down an equation of the line which cuts the $x$-axis at $(2,0)$ and the $y$-axis at $(0,-4)$.

\bigskip
\noindent\textbf{585.} Find an equation of the line through the point $P(2,3)$ which forms an isosceles triangle with the coordinate axes in the first quadrant.

\bigskip
\noindent\textbf{586.} Transform the equation $5x-12y+39=0$ into slope-intercept form.

\bigskip
\noindent\textbf{587.} Transform the equation $5x-12y+39=0$ into two-intercept form.

\bigskip
\noindent\textbf{588.} Transform the equation $5x-12y+39=0$ into normal form.

\bigskip
\noindent\textbf{589.} Transform the equation $5x-12y+39=0$ into point-slope form.

\bigskip
\noindent\textbf{590.} Transform the equation $5x-12y+39=0$ into two-point form.

\bigskip
\noindent\textbf{591.} Transform the equation $5x-12y+39=0$ into symmetric form.

\bigskip
\noindent\textbf{592.} If the houses of two friends are represented by coordinates $(2,6)$ and $(9,12)$ on a grid, find the straight-line distance between their houses if the grid units represent kilometres.

\bigskip
\noindent\textbf{593.} Consider a straight trail (represented by a coordinate plane) that starts at point $(5,7)$ and ends at point $(15,3)$. What is the coordinate of the midpoint?

\bigskip
\noindent\textbf{594.} The start and end points of a race track are given by coordinates $(3,9)$ and $(9,13)$. What is the midpoint of the track?

\bigskip
\noindent\textbf{595.} The coordinates of two points on a road are $A(3,4)$ and $B(7,10)$. Find the midpoint of the road.

\bigskip
\noindent\textbf{596.} What is aviation?

\bigskip
\noindent\textbf{597.} What is navigation?

\bigskip
\noindent\textbf{598.} What is latitude?

\bigskip
\noindent\textbf{599.} What is longitude?

\bigskip
\noindent\textbf{600.} On a map, town A is at coordinates $(2,3)$ and town B is at $(-4,-1)$. What is the distance between the two towns?

\bigskip
\noindent\textbf{601.} Suppose two cities, City A and City B, are represented by the coordinates $(3,4)$ and $(7,1)$ on a map. Find the straight-line distance between the two cities.

\bigskip
\noindent\textbf{602.} An engineer is building a bridge between two points on a riverbank. Suppose the coordinates of the two points where the bridge will start and end are $(2,5)$ and $(8,9)$. Find the coordinates of the midpoint, which will represent the centre of the bridge.

\bigskip
\noindent\textbf{603.} A pilot needs to travel from city A $(50,60)$ to city B $(120,150)$. Determine the heading angle the plane should take relative to the east direction.

\bigskip
\noindent\textbf{604.} Abdul Hadi is traveling from point A (Latitude $10^{\circ}$N, Longitude $50^{\circ}$E) to point B (Latitude $20^{\circ}$N, Longitude $60^{\circ}$E). Find the midpoint of his journey in terms of latitude and longitude.

\bigskip
\noindent\textbf{605.} A landscaper is designing a straight pathway from $P(2,3)$ to $Q(8,9)$. What is the length of the pathway?

\bigskip
\noindent\textbf{606.} Find the distance between two points $A(2,3)$ and $B(7,8)$ on a coordinate plane.

\bigskip
\noindent\textbf{607.} Find the midpoint of the line segment joining the points $(4,-2)$ and $(-6,3)$.

\bigskip
\noindent\textbf{608.} Calculate the gradient (slope) of the line passing through the points $(1,2)$ and $(4,6)$.

\bigskip
\noindent\textbf{609.} If two lines are parallel and one line has a gradient of $\dfrac{2}{3}$, what is the gradient of the other line?

\bigskip
\noindent\textbf{610.} An airplane needs to fly from city A at coordinates $(12,5)$ to city B at coordinates $(8,-4)$. Calculate the straight-line distance between these two cities.

\bigskip
\noindent\textbf{611.} In a landscaping project, the path starts at $(2,3)$ and ends at $(10,7)$. Find the midpoint.

\bigskip
\noindent\textbf{612.} For a line with a gradient of $-3$ and a y-intercept of $2$, write the equation of the line in \textbf{slope-intercept} form.

\bigskip
\noindent\textbf{613.} For a line with a gradient of $-3$ and a y-intercept of $2$, write the equation of the line in \textbf{point-slope} form using the point $(1,2)$.

\bigskip
\noindent\textbf{614.} For a line with a gradient of $-3$ and a y-intercept of $2$, write the equation of the line in \textbf{two-point} form using the points $(1,2)$ and $(4,-7)$.

\bigskip
\noindent\textbf{615.} For a line with a gradient of $-3$ and a y-intercept of $2$, write the equation of the line in \textbf{intercepts} form.

\bigskip
\noindent\textbf{616.} For a line with a gradient of $-3$ and a y-intercept of $2$, write the equation of the line in \textbf{symmetric} form.

\bigskip
\noindent\textbf{617.} For a line with a gradient of $-3$ and a y-intercept of $2$, write the equation of the line in \textbf{normal} form.

\bigskip
\noindent\textbf{618.} Write the converse, inverse and contrapositive of the following conditional: $\sim p \rightarrow q$

\bigskip
\noindent\textbf{619.} Write the converse, inverse and contrapositive of the following conditional: $q \rightarrow p$

\bigskip
\noindent\textbf{620.} Write the converse, inverse and contrapositive of the following conditional: $\sim q \rightarrow \sim p$

\bigskip
\noindent\textbf{621.} Write the converse, inverse and contrapositive of the following conditional: $\sim q \rightarrow \sim p$

\bigskip
\noindent\textbf{622.} Write the truth table of the following: $\sim(p\vee q)\vee(\sim q)$

\bigskip
\noindent\textbf{623.} Write the truth table of the following: $\sim(\sim q \vee \sim p)$

\bigskip
\noindent\textbf{624.} Write the truth table of the following: $(p\vee q) \leftrightarrow (p \wedge p)$

\bigskip
\noindent\textbf{625.} Differentiate between a mathematical statement and its proof. Give two examples.

\bigskip
\noindent\textbf{626.} What is the difference between an axiom and a theorem? Give examples of each.

\bigskip
\noindent\textbf{627.} What is the importance of logical reasoning in mathematical proofs? Give an example to illustrate your point.

\bigskip
\noindent\textbf{628.} Indicate whether the following is an axiom, conjecture or theorem, and explain your reasoning: There is exactly one straight line through any two points.

\bigskip
\noindent\textbf{629.} Indicate whether the following is an axiom, conjecture or theorem, and explain your reasoning: Every even number greater than 2 can be written as the sum of two prime numbers.

\bigskip
\noindent\textbf{630.} Indicate whether the following is an axiom, conjecture or theorem, and explain your reasoning: The sum of the angles in a triangle is $180^{\circ}$.

\bigskip
\noindent\textbf{631.} Suppose $x$ is an integer. Then $x$ is odd if and only if $9x+4$ is odd.

\bigskip
\noindent\textbf{632.} Suppose $x$ is an integer. If $x$ is odd, then $7x+5$ is even.

\bigskip
\noindent\textbf{633.} If $x$ and $y$ are positive real numbers and $x^2<y^2$, then $x<y$.

\bigskip
\noindent\textbf{634.} Prove that the sum of the interior angles of a triangle is $180^{\circ}$.

\bigskip
\noindent\textbf{635.} If $a$, $b$ and $c$ are non-zero real numbers, prove that: $\dfrac{a}{b}=\dfrac{c}{d} \Leftrightarrow ad=bc$

\bigskip
\noindent\textbf{636.} If $a$, $b$ and $c$ are non-zero real numbers, prove that: $\dfrac{a}{b}\cdot\dfrac{c}{d}=\dfrac{ac}{bd}$

\bigskip
\noindent\textbf{637.} If $a$, $b$ and $c$ are non-zero real numbers, prove that: $\dfrac{a}{b}+\dfrac{c}{b}=\dfrac{a+c}{b}$

\bigskip
\noindent\textbf{638.} Define induction.

\bigskip
\noindent\textbf{639.} Define deduction.

\bigskip
\noindent\textbf{640.} Define statement.

\bigskip
\noindent\textbf{641.} What is negation?

\bigskip
\noindent\textbf{642.} What is conjunction?

\bigskip
\noindent\textbf{643.} What is disjunction?

\bigskip
\noindent\textbf{644.} Define equivalence.

\bigskip
\noindent\textbf{645.} Define theorem.

\bigskip
\noindent\textbf{646.} What is a conjecture?

\bigskip
\noindent\textbf{647.} Define axiom.

\bigskip
\noindent\textbf{648.} What is Euclid's axiom?

\bigskip
\noindent\textbf{649.} What is Peano's axiom?

\bigskip
\noindent\textbf{650.} What is a deductive proof?

\bigskip
\noindent\textbf{651.} State whether the following statement is true or false: Lahore is the capital of Punjab and Quetta is the capital of Balochistan.

\bigskip
\noindent\textbf{652.} State whether the following statement is true or false: $4<5 \wedge 8<10$

\bigskip
\noindent\textbf{653.} State whether the following statement is true or false: $2+2=3 \wedge 6+6=10$

\bigskip
\noindent\textbf{654.} ``10 is a positive integer'' or ``0 is a rational number.'' Find the truth value of this disjunction.

\bigskip
\noindent\textbf{655.} ``A triangle can have two right angles'' or ``Lahore is the capital of Sindh.'' Find the truth value of this disjunction.

\bigskip
\noindent\textbf{656.} Prove that in any universal set, the empty set $\phi$ is a subset of any set $A$.

\bigskip
\noindent\textbf{657.} Prove the following mathematical statement: If $x$ is an odd integer, then $x^2$ is also an odd integer.

\bigskip
\noindent\textbf{658.} Prove the following mathematical statement: The sum of two odd numbers is an even number.

\bigskip
\noindent\textbf{659.} Find whether the solids are similar. All lengths are in cm. \emph{[Two prisms with dimensions -- see original figure]}

\bigskip
\noindent\textbf{660.} In triangle $ABC$, the sides are given as $m(\seg{AB})=6$ cm, $m(\seg{BC})=9$ cm and $m(\seg{CA})=12$ cm. In triangle $DEF$, the sides are given as $m(\seg{DE})=10.5$ cm, $m(\seg{EF})=15.75$ cm and $m(\seg{FD})=21$ cm. Prove that the triangles are similar.

\bigskip
\noindent\textbf{661.} Find the value of $x$ in the following: \emph{[figure -- see original paper]}

\bigskip
\noindent\textbf{662.} Find the value of $x$ in the following: \emph{[figure -- see original paper]}

\bigskip
\noindent\textbf{663.} Find the value of $x$ in the following: \emph{[figure -- see original paper]}

\bigskip
\noindent\textbf{664.} Find the value of $x$, $y$ and $z$ in the given figure: \emph{[figure -- see original paper]}

\bigskip
\noindent\textbf{665.} If one pair of corresponding sides are parallel to each other, then the triangles so formed as shown in the figure are similar. In the figure, $\seg{AB}$ is parallel to $\seg{CD}$. \emph{[figure]}

\bigskip
\noindent\textbf{666.} The perimeter of a regular octagon is 48 cm. Another octagon has sides that are 1.2 times the sides of the first octagon. What is the length of the side of the second octagon?

\bigskip
\noindent\textbf{667.} Find the ratio of the areas of similar figures if the ratio of their corresponding lengths is $1:3$

\bigskip
\noindent\textbf{668.} Find the ratio of the areas of similar figures if the ratio of their corresponding lengths is $3:4$

\bigskip
\noindent\textbf{669.} Find the ratio of the areas of similar figures if the ratio of their corresponding lengths is $2:7$

\bigskip
\noindent\textbf{670.} Find the ratio of the areas of similar figures if the ratio of their corresponding lengths is $8:9$

\bigskip
\noindent\textbf{671.} The radii of two spheres are in the ratio $3:4$. What is the ratio of their volumes?

\bigskip
\noindent\textbf{672.} Two regular tetrahedrons have volumes in the ratio $8:27$. What is the ratio of their sides?

\bigskip
\noindent\textbf{673.} Two right cones have volumes in the ratio $64:125$. What is the ratio of their heights?

\bigskip
\noindent\textbf{674.} Two right cones have volumes in the ratio $64:125$. What is the ratio of their base areas?

\bigskip
\noindent\textbf{675.} Find the missing value in the following similar solid: \emph{[figure -- see original paper]}

\bigskip
\noindent\textbf{676.} Find the missing value in the following similar solid: \emph{[figure -- see original paper]}

\bigskip
\noindent\textbf{677.} Find the missing value in the following similar solid: \emph{[figure -- see original paper]}

\bigskip
\noindent\textbf{678.} Find the unknown volume in the following similar solid: \emph{[figure -- see original paper]}

\bigskip
\noindent\textbf{679.} What is the sum of the interior angles of a decagon (10-sided polygon)?

\bigskip
\noindent\textbf{680.} What is each exterior angle of a regular pentagon?

\bigskip
\noindent\textbf{681.} Calculate the measure of each interior angle of a regular hexagon.

\bigskip
\noindent\textbf{682.} If the sum of the interior angles of a polygon is $1260^{\circ}$, how many sides does the polygon have?

\bigskip
\noindent\textbf{683.} In a parallelogram $ABCD$, $m(\seg{AB})=10$ cm, $m(\seg{AD})=6$ cm and $m\angle BAD=45^{\circ}$. Calculate the length of diagonal $m(\seg{AC})$.

\bigskip
\noindent\textbf{684.} A rectangular wall is 10 m tall and 120 m wide. How many gallons of paint are needed to cover the wall, if one gallon covers $35\ \text{m}^2$?

\bigskip
\noindent\textbf{685.} A rectangular floor is 12 m by 15 m. How many square tiles, each 1 m by 1 m, are needed to cover the floor?

\bigskip
\noindent\textbf{686.} A rectangular wall has a length of 10 m and a width of 4 m. If 1 litre of paint covers $7\ \text{m}^2$, how many litres of paint are needed to cover the wall?

\bigskip
\noindent\textbf{687.} Find the area of the equilateral triangle $ABC$ of side length $s$ metres.

\bigskip
\noindent\textbf{688.} Sketch the graph of the following linear function: $y=3x-5$

\bigskip
\noindent\textbf{689.} Sketch the graph of the following linear function: $y=-2x+8$

\bigskip
\noindent\textbf{690.} Sketch the graph of the following linear function: $y=0.5x-1$

\bigskip
\noindent\textbf{691.} Plot the graph of the following function: $y=4^x$

\bigskip
\noindent\textbf{692.} Plot the graph of the following function: $y=5^{-x}$

\bigskip
\noindent\textbf{693.} Plot the graph of the following function: $y=\dfrac{1}{x-3},\ x\ne 3$

\bigskip
\noindent\textbf{694.} Plot the graph of the following function: $y=\dfrac{2}{x}+3,\ x\ne 0$

\bigskip
\noindent\textbf{695.} Plot the graph of the following function: $y=x^{\frac{1}{2}}$

\bigskip
\noindent\textbf{696.} Plot the graph of the following function: $y=3x^{\frac{1}{3}}$

\bigskip
\noindent\textbf{697.} Plot the graph of the following function: $y=2x^{-2}$

\bigskip
\noindent\textbf{698.} What are graphs?

\bigskip
\noindent\textbf{699.} Define quadratic function.

\bigskip
\noindent\textbf{700.} Define cubic function.

\bigskip
\noindent\textbf{701.} Define reciprocal function.

\bigskip
\noindent\textbf{702.} Sketch the graph of: $y=2x-1$

\bigskip
\noindent\textbf{703.} Plot the graphs of: $y=x^2$ and $y=-x^2$

\bigskip
\noindent\textbf{704.} Sketch the graph of $y=2x^2-7x-9$ for $-3\leq x\leq6$.

\bigskip
\noindent\textbf{705.} Plot the graph of the following cubic function for $-3\leq x\leq3$: $y=-x^3$

\bigskip
\noindent\textbf{706.} Plot the graph of $y=2x^3-3x^2+x-5$ for $-2\leq x\leq3$.

\bigskip
\noindent\textbf{707.} Sketch the graph of the following reciprocal function: $y=\dfrac{1}{x-0.5},\ x\neq0.5$

\bigskip
\noindent\textbf{708.} Plot the graph of the exponential function $y=2^{x}$, $-6\leq x\leq6$.

\bigskip
\noindent\textbf{709.} Plot the graph of the exponential function $y=e^{x}$.

\bigskip
\noindent\textbf{710.} Plot the graph of $y=x^3$ for $-3\leq x\leq3$.

\bigskip
\noindent\textbf{711.} Plot the graph of $y=x^{-1}$.

\bigskip
\noindent\textbf{712.} Plot the graph of $y=2x^{\frac{1}{5}}$.

\bigskip
\noindent\textbf{713.} A company manufactures school bags. The function of producing $x$ bags is $C(x)=1200+800x$ and the revenue from selling $x$ bags is $R(x)=50x$. Find the break-even point.

\bigskip
\noindent\textbf{714.} The demand and supply functions for a product are given by $P_d=400-5Q$ and $P_s=3Q+24$. Plot the graph of each function over the interval $Q=0$ to $Q=300$.

\bigskip
\noindent\textbf{715.} A company manufactures school bags. The function of producing $x$ bags is $C(x)=1200+800x$ and the revenue from selling $x$ bags is $R(x)=50x$. Determine the profit or loss when 250 bags are sold.

\bigskip
\noindent\textbf{716.} Shahid's salary $S(x)$ in rupees is based on the formula $S(x)=45000+4500x$, where $x$ is the number of years he has been with the company. Sketch and interpret the graph of the salary function for $0\leq x\leq5$.

\bigskip
\noindent\textbf{717.} A company manufactures school bags. The function of producing $x$ bags is $C(x)=1200+800x$ and the revenue from selling $x$ bags is $R(x)=50x$. Plot the graph of both functions and identify the break-even point.

\bigskip
\noindent\textbf{718.} Plot the graph of the following function: $y=3^{-x}$ for $x$ from $-2$ to $4$.

\bigskip
\noindent\textbf{719.} Plot the graph of the following function: $y=\dfrac{2}{x+7},\ x\neq-7$

\bigskip
\noindent\textbf{720.} Plot the graph of the following for $x$ from $-5$ to $5$: $y=x^2-3$

\bigskip
\noindent\textbf{721.} Plot the graph of the following for $x$ from $-5$ to $5$: $y=15-x^2$

\bigskip
\noindent\textbf{722.} What is meant by locus?

\bigskip
\noindent\textbf{723.} What is a scalene triangle?

\bigskip
\noindent\textbf{724.} What is an isosceles triangle?

\bigskip
\noindent\textbf{725.} What is an equilateral triangle?

\bigskip
\noindent\textbf{726.} What is an acute-angled triangle?

\bigskip
\noindent\textbf{727.} What is an obtuse-angled triangle?

\bigskip
\noindent\textbf{728.} What is a right-angled triangle?

\bigskip
\noindent\textbf{729.} What is SSA?

\bigskip
\noindent\textbf{730.} Define perpendicular bisector.

\bigskip
\noindent\textbf{731.} Define point of concurrency.

\bigskip
\noindent\textbf{732.} Define circumcenter.

\bigskip
\noindent\textbf{733.} Define centroid.

\bigskip
\noindent\textbf{734.} What is the angle bisector of a triangle?

\bigskip
\noindent\textbf{735.} Define incenter.

\bigskip
\noindent\textbf{736.} Define altitude.

\bigskip
\noindent\textbf{737.} Define orthocenter.

\bigskip
\noindent\textbf{738.} Define median.

\bigskip
\noindent\textbf{739.} What is loci?

\bigskip
\noindent\textbf{740.} Define circle.

\bigskip
\noindent\textbf{741.} Define angle bisector.


\bigskip
\noindent\textbf{Questions 742--749} refer to the following distribution, which represents the scores achieved by a group of chemistry students in the chemistry laboratory:

\begin{center}
\begin{tabular}{lcccccc}
\toprule
Scores & 24--28 & 29--33 & 34--38 & 39--43 & 44--48 & 49--53 \\
\midrule
No. of students & 3 & 6 & 12 & 23 & 15 & 6 \\
\bottomrule
\end{tabular}
\quad (Total = 65)
\end{center}

\bigskip
\noindent\textbf{742.} What is the upper limit of the last class?

\bigskip
\noindent\textbf{743.} What is the lower limit of the class $39$--$43$?

\bigskip
\noindent\textbf{744.} What is the midpoint of the class $(34$--$38)$?

\bigskip
\noindent\textbf{745.} What are the class frequencies of the classes $29$--$33$ and $44$--$48$?

\bigskip
\noindent\textbf{746.} What is the size of the class limits in the above frequency distribution?

\bigskip
\noindent\textbf{747.} In which class or group does the minimum number of students fall?

\bigskip
\noindent\textbf{748.} What is the lower limit of the class having 15 as its class frequency?

\bigskip
\noindent\textbf{749.} What is the number of students having scores between 24 and 43?

\bigskip
\noindent\textbf{750.} The following data show the number of heads in an experiment of 50 sets of tossing a coin 5 times. Make a discrete frequency distribution from the information:

\begin{center}
\begin{tabular}{ccccccccccccc}
3&3&4&0&5&4&3&3&1&2&4&5&0\\
3&2&4&4&0&0&0&5&5&3&2&1&\\
4&3&2&5&3&2&1&3&5&4&3&2&1\\
3&2&1&3&1&3&1&4&3&2&2&4&\\
\end{tabular}
\end{center}

\bigskip
\noindent\textbf{751.} Find the arithmetic mean in the following: $4,6,10,12,15,20,25,28,30$.

\bigskip
\noindent\textbf{752.} Find the arithmetic mean in the following: $12,18,19,0,-19,-18,-12$.

\bigskip
\noindent\textbf{753.} Find the arithmetic mean in the following: $6.5,11,12.3,9,8.1,16,18,20.5,25$.

\bigskip
\noindent\textbf{754.} Find the arithmetic mean in the following: $8,10,12,14,16,20,22$.

\bigskip
\noindent\textbf{755.} Following are the heights (in inches) of 12 students. Find the median height: $55,53,54,58,60,61,62,56,57,52,51,63$.

\bigskip
\noindent\textbf{756.} Following are the earnings (in Rs) of ten workers: $88,70,72,125,115,95,81,90,95,90$. Calculate the arithmetic mean.

\bigskip
\noindent\textbf{757.} Following are the earnings (in Rs) of ten workers: $88,70,72,125,115,95,81,90,95,90$. Calculate the median.

\bigskip
\noindent\textbf{758.} Following are the earnings (in Rs) of ten workers: $88,70,72,125,115,95,81,90,95,90$. Calculate the mode.

\bigskip
\noindent\textbf{759.} Ten boys work at a petrol pump station. They get weekly wages as follows (in Rs): $4250,4350,4400,4250,4350,4410,4500,4300,4500,4390$. Find the arithmetic mean by the short formula, the median, and the mode of their wages.

\bigskip
\noindent\textbf{760.} The arithmetic mean of 45 numbers is 80. Find their sum.

\bigskip
\noindent\textbf{761.} Five numbers are $1,4,0,7,9$. Find their mean, median, and mode.

\bigskip
\noindent\textbf{762.} A set of data contains the values $148,145,160,157,156,160$. Show that Mode $>$ Median $>$ Mean.

\bigskip
\noindent\textbf{763.} Calculate the arithmetic mean from the following information: $D=X-140$, $\sum D=500$ and $n=10$.

\bigskip
\noindent\textbf{764.} Calculate the arithmetic mean from the following information: $U=\dfrac{x-130}{6}$, $\sum U=-150$ and $n=15$.

\bigskip
\noindent\textbf{765.} Calculate the arithmetic mean from the following information: $D=x-25$, $\sum fD=300$ and $\sum f=20$.

\bigskip
\noindent\textbf{766.} Calculate the arithmetic mean from the following information: $U=\dfrac{x-120}{5}$, $\sum fU=60$ and $\sum f=100$.

\bigskip
\noindent\textbf{767.} For the following data, find the weighted mean:
\begin{center}
\begin{tabular}{lcc}
\toprule
Item & Quantity & Cost of item (in thousands) \\
\midrule
Washing Machine & 5 & 35 \\
Heater & 3 & 5 \\
Stove & 2 & 13 \\
Dispenser & 6 & 18 \\
\bottomrule
\end{tabular}
\end{center}

\bigskip
\noindent\textbf{768.} A company is planning its next year marketing budget across five years: yearly budget (in million) is $5,7,8,6,7$. Find the average budget for the next year.

\bigskip
\noindent\textbf{769.} Ahmad obtained the following marks in a certain examination. Find the weighted mean if weights $5,4,2,3,2,4$ respectively are allotted to the subjects:
\begin{center}
\begin{tabular}{lcccccc}
\toprule
& Urdu & English & Science & Math & Islamiyat & Computer \\
\midrule
Marks & 78 & 65 & 80 & 90 & 85 & 72 \\
\bottomrule
\end{tabular}
\end{center}

\bigskip
\noindent\textbf{770.} The marks of a student in five examinations were $64,75,81,87,90$. Find the arithmetic mean of the marks.

\bigskip
\noindent\textbf{771.} Find the arithmetic mean using the short formula for the runs made by a batsman: Runs $=40,45,50,52,50,60,56,70$.

\bigskip
\noindent\textbf{772.} The following are the scores made by a batsman: $8,12,18,13,16,5,50$. Find the median of the data.

\bigskip
\noindent\textbf{773.} Following are the marks out of 100 obtained by 10 students in English: $23,15,35,48,41,5,8,9,11,51$. Find the median of the data.

\bigskip
\noindent\textbf{774.} The following data describes the marks of a student in different subjects and the weights assigned to these subjects are also given:
\begin{center}
\begin{tabular}{lcccc}
\toprule
Marks $(x)$ & 74 & 78 & 74 & 90 \\
Weights $(w)$ & 4 & 3 & 5 & 6 \\
\bottomrule
\end{tabular}
\end{center}
Find its weighted mean.

\bigskip
\noindent\textbf{775.} Define the following: Frequency distribution.

\bigskip
\noindent\textbf{776.} Define the following: Histogram (unequal class limits).

\bigskip
\noindent\textbf{777.} Define the following: Mean.

\bigskip
\noindent\textbf{778.} Define the following: Median.


\bigskip
\noindent\textbf{Questions 779--781} refer to the following data, which show the marks obtained by 45 students in the monthly test of Biology:

\begin{center}
\begin{tabular}{lcccccc}
\toprule
Marks & 20--24 & 25--29 & 30--34 & 35--39 & 40--44 & 45--49 \\
\midrule
No. of students & 05 & 08 & 12 & 15 & 03 & 02 \\
\bottomrule
\end{tabular}
\end{center}

\bigskip
\noindent\textbf{779.} With reference to the above table, find: the upper class boundary of the 5th class.

\bigskip
\noindent\textbf{780.} With reference to the above table, find: the lower class boundaries of all the classes.

\bigskip
\noindent\textbf{781.} With reference to the above table, find: the midpoint of all the classes.

\bigskip
\noindent\textbf{782.} Given below are marks obtained by 45 students in the monthly test of Biology:
\begin{center}
\begin{tabular}{lcccccc}
\toprule
Marks & 20--24 & 25--29 & 30--34 & 35--39 & 40--44 & 45--49 \\
\midrule
No. of students & 05 & 08 & 12 & 15 & 03 & 02 \\
\bottomrule
\end{tabular}
\end{center}
With reference to the above table, find the following: the class interval with the least frequency.

\bigskip
\noindent\textbf{783.} The marks of a student Saad in six tests were $84, 91, 72, 68, 87, 78$. Find the arithmetic mean of his marks.

\bigskip
\noindent\textbf{784.} Usman rolled a fair die eight times. Each time their sum was recorded as $8,5,6,6,9,4,3,11$. Find the median and mode of the sums.

\bigskip
\noindent\textbf{785.} Arshad rolls a die with sides labelled $L,M,N,O,P,U$. What is the probability that the die lands on a consonant?

\bigskip
\noindent\textbf{786.} Shazia throws a pair of fair dice. What will be the probability of getting: sum of dots is at least $4$.

\bigskip
\noindent\textbf{787.} Shazia throws a pair of fair dice. What will be the probability of getting: product of both dots between $5$ and $10$.

\bigskip
\noindent\textbf{788.} Shazia throws a pair of fair dice. What will be the probability of getting: the difference between both the dots is equal to $4$.

\bigskip
\noindent\textbf{789.} Shazia throws a pair of fair dice. What will be the probability of getting: number at least $5$ on the first die and the number at least $4$ on the second die.

\bigskip
\noindent\textbf{790.} One letter is selected at random from the word ``MATHEMATICS''. Find the probability of getting: a vowel.

\bigskip
\noindent\textbf{791.} One letter is selected at random from the word ``MATHEMATICS''. Find the probability of getting: a consonant.

\bigskip
\noindent\textbf{792.} One letter is selected at random from the word ``MATHEMATICS''. Find the probability of getting: an E.

\bigskip
\noindent\textbf{793.} One letter is selected at random from the word ``MATHEMATICS''. Find the probability of getting: an A.

\bigskip
\noindent\textbf{794.} One letter is selected at random from the word ``MATHEMATICS''. Find the probability of getting: not M.

\bigskip
\noindent\textbf{795.} One letter is selected at random from the word ``MATHEMATICS''. Find the probability of getting: not T.

\bigskip
\noindent\textbf{796.} What is probability?

\bigskip
\noindent\textbf{797.} Write the formula of probability.

\bigskip
\noindent\textbf{798.} Define experiment.

\bigskip
\noindent\textbf{799.} Define outcomes.

\bigskip
\noindent\textbf{800.} Define favourable outcome.

\bigskip
\noindent\textbf{801.} Define sample space.

\bigskip
\noindent\textbf{802.} Define event.

\bigskip
\noindent\textbf{803.} Define sample point.

\bigskip
\noindent\textbf{804.} What is a certain event?

\bigskip
\noindent\textbf{805.} What is an impossible event?

\bigskip
\noindent\textbf{806.} What is a likely event?

\bigskip
\noindent\textbf{807.} What is an unlikely event?

\bigskip
\noindent\textbf{808.} What are equally likely events?

\bigskip
\noindent\textbf{809.} Abdul Rehman rolls a fair die; what is the probability of getting a number divisible by $3$?

\bigskip
\noindent\textbf{810.} If Zeeshan rolled two fair dice, find the probability of getting: even numbers on both dice.

\bigskip
\noindent\textbf{811.} If Zeeshan rolled two fair dice, find the probability of getting: multiples of $3$ on both dice.

\bigskip
\noindent\textbf{812.} If Zeeshan rolled two fair dice, find the probability of getting: an even number on the first die and the number $3$ on the second die.

\bigskip
\noindent\textbf{813.} If Zeeshan rolled two fair dice, find the probability of getting: at least the number $3$ on the first die and number $4$ on the second die.

\bigskip
\noindent\textbf{814.} Zubair rolls a die; what is the probability of not getting the number $6$?

\bigskip
\noindent\textbf{815.} If two fair dice are rolled, what is the probability of getting: not the sum of both dice is $8$.

\bigskip
\noindent\textbf{816.} If two fair dice are rolled, what is the probability of getting: not a double six.

\bigskip
\noindent\textbf{817.} A card is drawn at random from a pack of $52$ playing cards. What is the probability of getting: a card of heart.

\bigskip
\noindent\textbf{818.} A card is drawn at random from a pack of $52$ playing cards. What is the probability of getting: neither spade nor heart.

\bigskip
\noindent\textbf{819.} In $500$ trials of throwing two dice, what is the expected frequency that the sum will be greater than $8$?

\bigskip
\noindent\textbf{820.} What is the expectation of a person who is to get Rs.\ $120$ if he obtains at least $2$ heads in a single toss of three coins?

\bigskip
\noindent\textbf{821.} What is relative frequency?

\bigskip
\noindent\textbf{822.} How is relative frequency calculated?

\bigskip
\noindent\textbf{823.} Define expected frequency.

\bigskip
\noindent\textbf{824.} Find the average number of times of getting $1$ or $6$, when a fair die is rolled $300$ times.

\bigskip
\noindent\textbf{825.} If the probability of a defective bolt is $0.3$, find the number of non-defective bolts in a total of $800$.



\clearpage

\section*{Part 3 -- Long Questions (Q.1--448)}
\addcontentsline{toc}{section}{Part 3 -- Long Questions (Q.1--448)}

\begin{detq}

\item If $a=\dfrac{2}{3}$, $b=\dfrac{3}{2}$, $c=\dfrac{5}{3}$, then verify the distributive property over addition.

\item Simplify the following: $\dfrac{5\cdot(25)^{n+1}-25\cdot(5)^{2n}}{5\cdot(5)^{2n+3}-(25)^{n+1}}$

\item Simplify the following: $\dfrac{3^n \times 9^{n+1}}{3^{n-1}\times 9^{n-1}}$

\item Simplify the following: $\dfrac{(16)^{x+1}+20(4)^{2x}}{2^{x-3}\times 8^{x+2}}$

\item Simplify the following: $\dfrac{5^{n+3}-6\cdot5^{n+1}}{9\times5^n - 2^n\times5^n}$

\item If $x=3+\sqrt{8}$, then find the value of $x^2-\dfrac{1}{x^2}$.

\item Find the rational numbers $p$ and $q$ such that $\dfrac{8-3\sqrt{2}}{4+3\sqrt{2}} = p+q\sqrt{2}$.

\item If $x=3+\sqrt{8}$, then find the value of $x^4+\dfrac{1}{x^4}$.

\item Simplify the following: $\dfrac{(25)^{3/2}\times(243)^{3/5}}{(16)^{5/4}\times(8)^{4/3}}$

\item Simplify the following: $\dfrac{54\times\sqrt[3]{(27)^{2x}}}{9^{x+1}+216(3)^{2x-1}}$

\item Simplify the following: $\left(a^{1/3}+b^{2/3}\right)\times\left(a^{2/3}-a^{1/3}b^{2/3}+b^{4/3}\right)$

\item Simplify the following: $\sqrt{\dfrac{(216)^{2/3}\times(25)^{1/2}}{(0.04)^{-3/2}}}$

\item The diagram shows a right-angled $\triangle ABC$ in which the length of $\overline{AC}$ is $(\sqrt{3}+\sqrt{5})$ cm. The area of $\triangle ABC$ is $(1+\sqrt{15})$ cm$^2$. Find the length $AB$ in the form $(a\sqrt{3}+b\sqrt{15})$ cm, where $a$ and $b$ are integers.
\begin{center}
\begin{tikzpicture}[scale=0.9]
\coordinate (A) at (0,0);
\coordinate (C) at (4.2,0);
\coordinate (B) at (0,2.4);
\draw (A) -- (C) -- (B) -- cycle;
\draw (0,0.3) -- (0.3,0.3) -- (0.3,0);
\node[below] at (2.1,-0.15) {$(\sqrt{3}+\sqrt{5})$ cm};
\node[left] at (-0.1,1.2) {$AB$};
\node[above right] at (2.3,1.4) {$BC$};
\node[below left] at (A) {$A$};
\node[below right] at (C) {$C$};
\node[above left] at (B) {$B$};
\end{tikzpicture}
\end{center}

\item A rectangle has sides of length $\left(2+\sqrt{18}\right)$ m and $\left(5-\dfrac{4}{\sqrt{2}}\right)$ m. Express the area of the rectangle in the form $a+b\sqrt{2}$, where $a$ and $b$ are integers.

\item Find two numbers whose sum is $68$ and difference is $22$.

\item The weather in Lahore was unusually warm during the summer of 2024. The TV news reported temperatures as high as $48^{\circ}$C. Using the formula $F=\dfrac{9}{5}C+32$, find the temperature on the Fahrenheit scale.

\item The sum of the ages of the father and son is $72$ years. Six years ago, the father's age was $2$ times the age of the son. What was the son's age six years ago?

\item Mirha sells a toy for Rs.\ $1520$. What will the selling price be to get a $15\%$ profit?

\item The annual income of Tayyab is Rs.\ $9{,}60{,}000$, while the exempted amount is Rs.\ $1{,}30{,}000$. How much tax would he have to pay at the rate of $0.75\%$?

\item Find the compound markup on Rs.\ $3{,}75{,}000$ for one year at the rate of $14\%$ compounded annually.

\item The sum of three consecutive odd integers is $51$. Find the three integers.

\item Abdullah picked up $96$ balls and placed them into two buckets. One bucket has twenty-eight more balls than the other bucket. How many balls were in each bucket?

\item Salma invested Rs.\ $3{,}50{,}000$ in a bank, which paid simple profit at the rate of $7\dfrac{1}{4}\%$ per annum. After $2$ years, the rate was increased to $8\%$ per annum. Find the amount she had at the end of $7$ years.

\item Hamail purchased a bicycle for Rs.\ $6590$ and sold it for Rs.\ $6850$. Find the profit percentage.

\item Umair bought a book for Rs.\ $850$ and sold it for Rs.\ $720$. What was his loss percentage?

\item Saleem, Nadeem, and Tanveer earned a profit of Rs.\ $4{,}50{,}000$ from a business. If their investments in the business are in the ratio $4:7:14$, find each person's profit.

\item If the simple profit on Rs.\ $6400$ for $12$ years is Rs.\ $3840$, find the rate of profit.

\item Find the value of $x$ in the following: $\log_4 x = \dfrac{3}{2}$

\item Find the value of $x$ in the following: $\log_2 1024 = x$

\item Find the value of $x$: $\log x = -3.434$

\item Find the value of $x$: $\log x = -1.5726$

\item Find the value of $x$: $\log x = 4.3561$

\item Find the value of $x$: $\log x = 2.0184$

\item Find the value of $x$ in the following: $\log x = -2.1234$

\item Without using a calculator, evaluate the following: $\dfrac{1}{3}\log_3 8 - \log_3 18$

\item Without using a calculator, evaluate the following: $\dfrac{1}{3}\log_4 64 + 2\log_5 25$

\item Without using a calculator, evaluate the following: $\log_3 12 + \log_3 0.25$

\item Find the value of $x$ in the following equation: $(81)^x = (243)^{x+2}$

\item Find the value of $x$ in the following equation: $\left(\dfrac{1}{27}\right)^{x-6}=27$

\item Find the value of $x$ in the following equation: $\log_2(x+1)-\log_2(x-4)=2$

\item Find the value of the following with the help of a logarithm table: $\dfrac{3.68\times4.21}{5.234}$

\item Find the value of the following with the help of a logarithm table: $4.67\times2.11\times2.397$

\item Find the value of the following with the help of a logarithm table: $\dfrac{(20.46)^2\times(2.4122)}{754.3}$

\item Find the value of the following with the help of a logarithm table: $\dfrac{\sqrt[3]{9.364}\times21.64}{3.21}$

\item The formula to measure the magnitude of earthquakes is given by $M=\log_{10}\left(\dfrac{A}{A_0}\right)$. If amplitude $(A)$ is $10{,}000$ and reference amplitude $(A_0)$ is $10$, what is the magnitude of the earthquake?

\item Abdullah invested Rs.\ $100{,}000$ in a saving scheme and gains interest at the rate of $5\%$ per annum so that the total value of this investment after $t$ years is Rs.\ $y$. This is modelled by the equation $y=100{,}000(1.05)^t$, $t\geq0$. Find after how many years the investment will be double.

\item Huria is hiking up a mountain where the temperature $(T)$ decreases by $3\%$ (or a factor of $0.97$) for every $100$ metres gained in altitude. The initial temperature $(T_i)$ at sea level is $20^{\circ}$C. Using the formula $T=T_i\times0.97^{h/100}$, calculate the temperature at an altitude $(h)$ of $500$ metres.

\item The decibel scale measures sound intensity using the formula $L=40\log_{10}\left(\dfrac{I}{I_0}\right)$. If a sound has an intensity $(I)$ of $10^6$ times the reference intensity $(I_0)$, what is the sound level in decibels?

\item Find the value of $x$ in the following: $\left(\dfrac{1}{32}\right)^{2x}=64$

\item Find the value of the following with the help of a logarithm table: $319.8\times3.543$

\item Find the value of the following with the help of a logarithm table: $\dfrac{36.12\times750.9}{113.2\times9.98}$

\item In the year 2016, the population of a city was $22$ million and was growing at a rate of $2.5\%$ per year. The function $p(t)=22(1.025)^t$ gives the population in millions, $t$ years after 2016. Use the model to determine in which year the population will reach $35$ million. Round the answer to the nearest year.

\item Consider the universal set $U=\{x: x \text{ is a multiple of } 2 \text{ and } 0<x\leq30\}$, $A=\{x: x \text{ is a multiple of } 6\}$ and $B=\{x: x \text{ is a multiple of } 8\}$: draw a Venn diagram.

\item Let $U=\{x: x \text{ is an integer and } 0<x\leq150\}$, $G=\{x: x=2^m \text{ for integer } m \text{ and } 0\leq m\leq12\}$ and $H=\{x: x \text{ is a square}\}$: find $G\cup H$.

\item Let $U=\{x: x \text{ is an integer and } 0<x\leq150\}$, $G=\{x: x=2^m \text{ for integer } m \text{ and } 0\leq m\leq12\}$ and $H=\{x: x \text{ is a square}\}$: find $G\cup H$.

\end{detq}

\begin{wq}[start=56]

\item Let $U=\{x: x \text{ is an integer and } 0<x\le 150\}$, $G=\{x: x=2^m \text{ for integer } m \text{ and } 0\le m\le 12\}$ and $H=\{x: x \text{ is a square}\}$. Find $G\cap H$.

\item Consider the sets $P=\{x: x \text{ is a prime number and } 0<x\le 20\}$ and $Q=\{x: x \text{ is a divisor of } 210 \text{ and } 0<x\le 20\}$. Find $P\cap Q$.

\item Consider the sets $P=\{x: x \text{ is a prime number and } 0<x\le 20\}$ and $Q=\{x: x \text{ is a divisor of } 210 \text{ and } 0<x\le 20\}$. Find $P\cup Q$.

\item If $U=\{1,2,3,\ldots,20\}$ and $A=\{1,3,5,\ldots,19\}$, verify the following: $A\cup A'=U$

\item If $U=\{1,2,3,\ldots,20\}$ and $A=\{1,3,5,\ldots,19\}$, verify the following: $A\cap A'=\phi$

\item In a class of 55 students, 34 like to play cricket and 30 like to play hockey. Also, each student likes to play at least one of the two games. How many students like to play both games?

\item In a group of 500 employees, 250 can speak Urdu, 150 can speak English, 50 can speak Punjabi, 40 can speak Urdu and English, 30 can speak both English and Punjabi, and 10 can speak Urdu and Punjabi. How many can speak all three languages?

\item In sports events, 19 people wear blue shirts, 15 wear green shirts, 3 wear blue and green shirts, 4 wear a cap and green shirts. The total number of people with either a blue or green shirt or cap is 25. How many people are wearing caps?

\item In a training session, 17 participants have laptops, 11 have tablets, 9 have laptops and tablets, 6 have laptops and books, and 4 have both tablets and books. Eight participants have all three items. The total number of participants with laptops, tablets, or books is 35. How many participants have books?

\item A shopping mall has 150 employees labelled 1 to 150, representing the universal set $U$. The employees fall into the following categories: Set $A$: 40 employees with a salary range of 30k--45k, labelled from 50 to 89. Set $B$: 50 employees with a salary range of 50k--80k, labelled from 101 to 150. Set $C$: 60 employees with a salary range of 100k--150k, labelled from 1 to 49 and 90 to 100. Find $(A'\cup B')\cap C$.

\item Using the same sets as above, find $n\{A\cap(B^c\cap C^c)\}$.

\item In a secondary school with 125 students, each participates in at least one of the following sports: cricket, football, or hockey. 60 students play cricket, 70 play football, 40 play hockey. 25 students play both cricket and football. 15 students play both football and hockey. 10 students play both cricket and hockey. How many students play all three sports?

\item Using the data from the previous question, draw a Venn diagram showing the distribution of sports participants in all the games.

\item A survey was conducted in which 130 people were asked about their favourite foods. The results: 40 people liked nihari, 65 liked biryani, 50 liked korma, 20 liked nihari and biryani, 35 liked biryani and korma, 27 liked nihari and korma, and 12 liked all three foods (nihari, biryani, and korma). At least how many people like nihari, biryani or korma?

\item Using the same survey data, how many people did not like nihari, biryani or korma?

\item Using the same survey data, how many people like only one of the following foods: nihari, biryani, or korma?

\item Using the same survey data, draw a Venn diagram.

\item Consider the adjacent Venn diagram illustrating two non-empty sets, $A$ and $B$. Determine the number of elements common to sets $A$ and $B$. \emph{[Venn diagram -- see original paper]}

\item Consider the adjacent Venn diagram illustrating two non-empty sets, $A$ and $B$. Identify all the elements exclusively in set $B$ and not in set $A$. \emph{[Venn diagram -- see original paper]}

\item Consider the adjacent Venn diagram illustrating two non-empty sets, $A$ and $B$. Calculate the union of sets $A$ and $B$. \emph{[Venn diagram -- see original paper]}

\item Consider the adjacent Venn diagram representing the students enrolled in different courses at an IT institution: $U=\{$Students enrolled in IT institutions$\}$, $A=\{$Students enrolled in Applied Robotics$\}$, $G=\{$Students enrolled in Game Development$\}$, $W=\{$Students enrolled in Web Designing$\}$. How many students enrolled in the Applied Robotics course? \emph{[Venn diagram -- see original paper]}

\item Using the same Venn diagram, determine the total number of students enrolled in Game Development.

\item Using the same Venn diagram, how many students are enrolled in both the Game Development and Web Designing courses?

\item Using the same Venn diagram, identify the students enrolled in Web Designing but not Applied Robotics.

\item Using the same Venn diagram, how many students are enrolled in IT institutions in total?

\item Using the same Venn diagram, how many students enrolled in all three courses?

\item There are 95 secondary school students in a sports club. 58 students join the swimming club, and 50 join the tug-of-war club. How many students participated in both games?

\item Mr.\ Saleem, a school teacher, has a small library in his house containing 150 books. He has two main categories for these books: Islamic and science. He categorized 70 books as Islamic books and 90 books as science books. There are 15 books that belong to neither the Islamic nor the science books category. How many books are classified under both the Islamic and science categories?

\item In a college, 45 teachers teach mathematics, physics, or chemistry. 18 teach mathematics, 12 teach physics, 8 teach chemistry, 6 teach both mathematics and physics, 4 teach both physics and chemistry, 2 teach both mathematics and chemistry. How many teachers teach all three subjects?

\item A survey of 130 customers in a shopping mall was conducted about buying preferences. Results: 57 customers bought garments, 50 bought cosmetics, 46 bought electronics, 31 purchased both garments and cosmetics, 25 purchased both garments and electronics, 21 purchased both cosmetics and electronics, and 12 purchased all three products (garments, cosmetics, and electronics). How many of the customers bought at least one of the products: garments, cosmetics, or electronics?

\item Using the same survey data, how many of the customers bought only one of the products: garments, cosmetics, or electronics?

\item Using the same survey data, how many customers did not buy any of the three products?

\item Given that $f(x)=ax+b+1$, where $a$ and $b$ are constant numbers. If $f(3)=8$ and $f(6)=14$, find the values of $a$ and $b$.

\item Given that $g(x)=ax+b+5$, where $a$ and $b$ are constant numbers. If $g(-1)=0$ and $g(2)=10$, find the values of $a$ and $b$.

\item Consider the function $f(x)=cx^2+d$, where $c$ and $d$ are constant numbers. If $f(1)=6$ and $f(-2)=10$, find the values of $c$ and $d$.

\item Let $c_1,c_2,c_3$ be three children and $m_1,m_2$ be two men such that the father of both $c_1,c_2$ is $m_1$ and the father of $c_3$ is $m_2$. Find the relation $\{(\text{child},\text{father})\}$.

\item Consider $f(x)=ax+b+3$, where $a$ and $b$ are constant numbers. If $f(1)=4$ and $f(5)=9$, find the values of $a$ and $b$.

\item Use a Venn diagram to verify the following: $A-B = A\cup B'$

\item Use a Venn diagram to verify the following: $(A-B)'\cap B = B$

\item Verify the following properties for the sets $A,B,C$ given below: (i) Associativity of union (ii) Associativity of intersection (iii) Distributivity of union over intersection (iv) Distributivity of intersection over union. \\
$A=\{1,2,3,4\}$, $B=\{3,4,5,6,7,8\}$, $C=\{5,6,7,9,10\}$

\item Verify the same properties (associativity of union/intersection, distributivity of union over intersection and intersection over union) for: \\
$A=\phi$, $B=\{0\}$, $C=\{0,1,2\}$

\item Verify the same properties (associativity of union/intersection, distributivity of union over intersection and intersection over union) for: \\
$A=N$, $B=Z$, $C=Q$

\item Verify De Morgan's Laws for the following sets: $U=\{1,2,3,\ldots,20\}$, $A=\{2,4,6,\ldots,20\}$ and $B=\{1,3,5,\ldots,19\}$

\end{wq}



\begin{enumerate}[leftmargin=1.8cm,start=99]

\item From suitable properties of union and intersection, deduce the following result: $A\cap(A\cup B)=A\cup(A\cap B)$

\item From suitable properties of union and intersection, deduce the following result: $A\cup(A\cap B)=A\cap(A\cup B)$

\item Given that $f(x)=ax+b$, where $a$ and $b$ are constant numbers, if $f(-2)=3$ and $f(4)=10$, then find the values of $a$ and $b$.

\item Consider the function $g(x)=mx^2+n$, where $m$ and $n$ are constant numbers. If $g(4)=20$ and $g(0)=5$, find the values of $m$ and $n$.

\item Out of the 180 students who appeared in the annual examination, 120 passed the math test, 90 passed the science test, and 60 passed both the math and science tests. How many did not pass either of the two tests?

\item Out of the 180 students who appeared in the annual examination, 120 passed the math test, 90 passed the science test, and 60 passed both the math and science tests. How many passed the science test but not the math test?

\item Out of the 180 students who appeared in the annual examination, 120 passed the math test, 90 passed the science test, and 60 passed both the math and science tests. How many failed the science test?

\item Let $U=\{1,2,3,4,5,6,7,8,9,10\}$, $A=\{2,4,6,8,10\}$, $B=\{1,2,3,4,5\}$ and $C=\{1,3,5,7,9\}$. List the members of the following set: $A'\cup C'$

\item Let $U=\{1,2,3,4,5,6,7,8,9,10\}$, $A=\{2,4,6,8,10\}$, $B=\{1,2,3,4,5\}$ and $C=\{1,3,5,7,9\}$. List the members of the following set: $A'\cup C$

\item A shopping mall has 100 products from various categories labeled 1 to 100, representing the universal set $U$. The products are categorized as follows: Set A: Electronics, consisting of 30 products labeled from 1 to 30; Set B: Clothing, comprising 25 products labeled from 31 to 55; Set C: Beauty products, comprising 25 products labeled from 76 to 100. Write each set in tabular form, and find the union of all three sets.

\item Out of the 180 students who appeared in the annual examination, 120 passed the math test, 90 passed the science test, and 60 passed both the math and science tests. How many passed either the math or the science test?

\item Factorize: $2x^2+11x+15$

\item Factorize: $3y^2-11y+6$

\item Factorize: $2y^2-5y+2$

\item Factorize: $4z^2-11z+6$

\item Factorize: $6+7x-3x^2$

\item Factorize: $2x^2+17x+26$

\item Factorize the following expression: $4x^4+81y^4$

\item Factorize the following expression: $a^4+64b^4$

\item Factorize the following expression: $x^4+4x^2+16$

\item Factorize the following expression: $x^4-14x^2+1$

\item Factorize the following expression: $x^4-30x^2y^2+9y^4$

\item Factorize the following expression: $x^4+11x^2y^2+y^4$

\item Factorize the following expression: $(x+1)(x+2)(x+3)(x+4)+1$

\item Factorize the following expression: $(x+2)(x-7)(x-4)(x-1)+17$

\item Factorize the following expression: $(2x^2+7x+3)(2x^2+7x+5)+1$

\item Factorize the following expression: $(3x^2+5x+3)(3x^2+5x+5)-3$

\item Factorize the following expression: $(x+1)(x+2)(x+3)(x+6)-3x^2$

\item Factorize the following expression: $(x+1)(x-1)(x+2)(x-2)+13x^2$

\item Factorize: $a^4+64$

\item Factorize: $(x^2+5x+4)(x^2+5x+6)-3$

\item Factorize: $(x-2)(x+2)(x+1)(x-4)+2x^2$

\item Factorize: $8x^3+60x^2+150x+125$

\item Factorize: $x^4+x^2+25$

\item Factorize: $x^4+y^4$

\item Find the HCF by factorization method: $a^3+2a^2-3a,\ 2a^3+5a^2-3a$

\item Find the HCF by factorization method: $t^2+3t-4,\ t^2+5t+4,\ t^2-1$

\item Find the HCF by factorization method: $x^2+15x+56,\ x^2+5x-24,\ x^2+8x$

\item Find the HCF of the following expressions by using the division method: $27x^3+9x^2-3x-9,\ 3x-2$

\item Find the HCF of the following expressions by using the division method: $x^3-9x^2+21x-15,\ x^2-4x+3$

\item Find the HCF of the following expressions by using the division method: $2x^3+2x^2+2x+2,\ 6x^3+12x^2+6x+12$

\item Find the HCF of the following expressions by using the division method: $2x^3-4x^2+6x,\ x^3-2x,\ 3x^2-6x$

\item Find the LCM of the following expressions by using the prime factorization method: $x^4-16,\ x^3-4x$

\item Find the LCM of the following expressions by using the prime factorization method: $16-4x^2,\ x^2+x-6,\ 4-x^2$

\item The HCF of two polynomials is $y-7$ and their LCM is $y^3-10y^2+11y+70$. If one of the polynomials is $y^2-5y-14$, find the other.

\item The LCM and HCF of two polynomials $p(x)$ and $q(x)$ are $36x^3(x+a)(x^3-a^3)$ and $x^2(x-a)$ respectively. If $p(x)=4x^2(x^2-a^2)$, find $q(x)$.

\item The HCF and LCM of two polynomials is $(x+a)$ and $12x^2(x+a)(x^2-a^2)$ respectively. Find the product of the two polynomials.

\item Find the HCF by factorization method: $x^2-27,\ x^2+6x-27,\ x^2-9$

\item Find the HCF of $6x^3-17x^2-5x+6$ and $6x^3-5x^2-3x+2$ by using the division method.

\item Find the LCM of the polynomials $x^2-3x+2$, $x^2-1$ and $x^2-5x+4$

\item The LCM and HCF of two polynomials are $x^3-10x^2+11x+70$ and $x-7$. If one polynomial is $x^2-12x+35$, find the other polynomial.

\item The LCM of $x^2y+xy^2$ and $x^2+xy$ is $xy(x+y)$. Find the HCF.

\item Find the square root of the following polynomial by factorization method: $200t^2-120t+18$

\item Find the square root of the following polynomial by factorization method: $40x^2+120x+90$

\item Find the square root of the following polynomial by division method: $4x^4-28x^3+37x^2+42x+9$

\item The potential energy $V(x)$ in an electric field varies as a cubic function of distance $x$, given by $V(x)=2x^3-6x^2+4x$. Determine where the potential energy is zero.

\item In structural engineering, the deflection $Y(x)$ of a beam is given by $Y(x)=2x^2-8x+6$. This equation gives the vertical deflection at any point $x$ along the beam. Find the points of zero deflection.

\item Find the square root of the following polynomial by division method: $121x^4-198x^3-183x^2+216x+144$

\item Find the square root of the following polynomial by division method: $x^4-10x^3y+27x^2y^2-10xy^3+y^4$

\item Find the square root of the following polynomial by division method: $x^4-10x^3y+27x^2y^2-10xy^3+y^4$

\item An investor's return $R(x)$ in rupees after investing $x$ thousand rupees is given by the quadratic expression $R(x)=-x^2+6x-8$. Factorize the expression and find the investment levels that result in zero return.

\item A company's profit $P(x)$ in rupees from selling $x$ units of a product is modeled by the cubic expression $P(x)=x^3-15x^2+75x-125$. Find the break-even point(s), where the profit is zero.

\item Find the square root of the expression $36x^4-36x^2+9$

\item Find the square root of the polynomial $x^4-12x^3+42x^2-36x+9$

\item Cost function for producing a part is modeled by $C(x)=5x^2-25x+30$, where $x$ is the width of the component and $C(x)$ is the cost. Find the value of $x$ where $C(x)$ is minimum.

\item A company's profit $P(x)$ is modeled by the quadratic equation $P(x)=-5x^2+50x-120$, where $x$ represents the number of units produced and $P(x)$ represents the profit in dollars. Find how many units should be produced to maximize profit.

\item The potential energy $U(x)$ of a particle moving in a cubic potential is expressed as $U(x)=x^3-6x^2+12x-8$. Factorize the expression to find the points where the energy is minimized.

\item Factorize the following: $(x+3)(x+4)(x+5)(x+6)-360$

\item Factorize the following: $(x^2+6x+3)(x^2+6x-9)+36$

\item Find the LCM and HCF by prime factorization method: $4x^3+12x^2,\ 8x^2+16x$

\item Find the LCM and HCF by prime factorization method: $x^3+3x^2-4x,\ x^2-x-6$

\item Huria is analyzing the total cost of her loan, modeled by the expression $C(x)=x^2-8x+15$, where $x$ represents the number of years. What is the optimal repayment period for Huria's loan?

\item Solve and represent the solution on a real line: $\dfrac{2x-1}{3}-\dfrac{3x}{4}=\dfrac{5}{6}$

\item Solve and represent the solution on a real line: $\dfrac{-5x}{10}=9-\dfrac{10}{5}x$

\item Solve the following inequality and represent the solution on a real line: $\dfrac{5}{3}x-\dfrac{3}{4}<\dfrac{-1}{12}$

\item Solve the following inequality and represent the solution on a real line: $\dfrac{1}{4}x-\dfrac{1}{2}\leq-1+\dfrac{1}{2}x$

\item Indicate the solution region of the following linear inequalities by shading: $2x-3y\leq6,\ 2x-3y\leq12$

\item Indicate the solution region of the following linear inequalities by shading: $x+y\geq5,\ -y+x\leq1$

\item Indicate the solution region of the following linear inequalities by shading: $3x+7y\geq21,\ x-y\leq2$

\item Indicate the solution region of the following linear inequalities by shading: $4x-3y\leq12,\ x\geq-\dfrac{3}{2}$

\item Indicate the solution region of the following linear inequalities by shading: $3x+7y\geq21,\ y\leq4$

\item Indicate the solution region of the following linear inequalities by shading: $5x+7y\leq35,\ x-2y\leq2$

\item Find the solution region by drawing the graph of the system of inequalities: $x-2y\leq6,\ 2x+y\geq2$

\item Find the minimum and maximum values of $z=3x+y$ subject to the constraints: $3x+5y\geq15;\ x+6y\geq9;\ x\geq0;\ y\geq0$

\item Maximize $f(x,y)=2x+5y$ subject to the constraints: $2y-x\leq8,\ x-y\leq4,\ x\geq0;\ y\geq0$

\item Maximize $f(x,y)=x+3y$ subject to the constraints: $2x+5y\leq30,\ 5x+4y\leq20,\ x\geq0;\ y\geq0$

\item Maximize $z=2x+3y$ subject to the constraints: $2x+y\leq4;\ 4x-y\leq4;\ x\geq0;\ y\geq0$

\item Minimize $z=2x+y$ subject to the constraints: $x+y\geq3;\ 7x+5y\leq35;\ x\geq0;\ y\geq0$

\item Maximize the function defined as $f(x,y)=2x+3y$ subject to the constraints: $2x+y\leq8;\ x+2y\leq14;\ x\geq0;\ y\geq0$

\item Shade the feasible region and find the corner points for the following system of inequalities: $x-y\leq3,\ x+2y\leq6,\ x\geq0,\ y\geq0$

\item Find the maximum and minimum values of the function defined as $f(x,y)=2x+3y$ subject to the constraints: $x-y\leq2,\ x+y\leq4,\ x\geq0,\ y\geq0$

\item Solve and represent the solution on a real line: $5(x-3)\geq26x-(10x+4)$

\item Find the solution region of the following linear inequalities: $3x-4y\leq12;\ 3x+2y\geq3$

\item Find the solution region of the following linear inequalities: $2x+y\leq4;\ x+2y\leq6$

\item Find the maximum value of $g(x,y)=x+4y$ subject to the constraints $x+y\leq4,\ x\geq0$ and $y\geq0$.

\item Find the minimum value of $f(x,y)=3x+5y$ subject to the constraints $x+3y\geq3,\ x+y\geq2,\ x\geq0,\ y\geq0$

\item If the central angle of a sector is $60^{\circ}$ and the radius of the circle is 12 cm, find the area of the sector and the percentage of the total area of the circle it represents.

\item A circular sector of radius $r=12$ cm has an angle of $150^{\circ}$. This sector is cut out and then bent to form a cone. What is the slant height and the radius of the base of this cone? (Hint: Arc length of sector = circumference of cone.)

\item Prove the following trigonometric identity: $\dfrac{\sin\theta}{\text{cosec}\,\theta}+\dfrac{\cos\theta}{\sec\theta}=1$

\end{enumerate}



\noindent\textbf{198.} Prove the following trigonometric identity:
\[
\frac{1-\sin\theta}{\cos\theta} = \frac{\cos\theta}{1+\sin\theta}
\]

\bigskip
\noindent\textbf{199.} Prove the following trigonometric identity:
\[
(\sec\theta - \tan\theta)^2 = \frac{1-\sin\theta}{1+\sin\theta}
\]

\bigskip
\noindent\textbf{200.} Prove the following trigonometric identity:
\[
(\tan\theta + \cot\theta)^2 = \sec^2\theta \, \csc^2\theta
\]

\bigskip
\noindent\textbf{201.} Prove the following trigonometric identity:
\[
\frac{\tan\theta + \sec\theta - 1}{\tan\theta - \sec\theta + 1} = \tan\theta + \sec\theta
\]

\bigskip
\noindent\textbf{202.} Prove the following trigonometric identity:
\[
\sin^3\theta - \cos^3\theta = (\sin\theta - \cos\theta)(1+\sin\theta\cos\theta)
\]

\bigskip
\noindent\textbf{203.} Prove the following trigonometric identity:
\[
\sin^6\theta - \cos^6\theta = (\sin^2\theta - \cos^2\theta)(1-\sin^2\theta\cos^2\theta)
\]

\bigskip
\noindent\textbf{204.} Show that:
\[
\frac{1}{\csc\theta - \cot\theta} - \frac{1}{\sin\theta} = \frac{1}{\sin\theta} - \frac{1}{\csc\theta + \cot\theta}
\]

\bigskip
\noindent\textbf{205.} Show that: $\sin^6\theta + \cos^6\theta = 1 - 3\sin^2\theta\cos^2\theta$

\bigskip
\noindent\textbf{206.} If $\tan\theta = \dfrac{3}{4}$, find the remaining trigonometric ratios, when $\theta$ lies in the first quadrant.

\bigskip
\noindent\textbf{207.} Evaluate: $\sin 60^{\circ}\cos 30^{\circ} + \cos 60^{\circ}\sin 30^{\circ}$

\bigskip
\noindent\textbf{208.} If $\sin\dfrac{\pi}{4}$ and $\cos\dfrac{\pi}{4}$ are equal to $\dfrac{1}{\sqrt{2}}$ each, then find the value of the following: $2\sin 45^{\circ} - 2\cos 45^{\circ}$

\bigskip
\noindent\textbf{209.} If $\sin\dfrac{\pi}{4}$ and $\cos\dfrac{\pi}{4}$ are equal to $\dfrac{1}{\sqrt{2}}$ each, then find the value of the following: $3\cos 45^{\circ} + 4\sin 45^{\circ}$

\bigskip
\noindent\textbf{210.} If $\sin\dfrac{\pi}{4}$ and $\cos\dfrac{\pi}{4}$ are equal to $\dfrac{1}{\sqrt{2}}$ each, then find the value of the following: $5\cos 45^{\circ} - 3\sin 45^{\circ}$

\bigskip
\noindent\textbf{211.} Solve the following triangle when $m\angle B = 90^{\circ}$: $m\angle C = 60^{\circ}$, $c = 3\sqrt{3}$ cm

\bigskip
\noindent\textbf{212.} Solve the following triangle when $m\angle B = 90^{\circ}$: $a = 12$ cm, $C = 6$ cm

\bigskip
\noindent\textbf{213.} Solve the following triangle when $m\angle B = 90^{\circ}$: $b = 10$ cm, $a = 6$ cm

\bigskip
\noindent\textbf{214.} Calculate the length $x$ in the adjoining figure. (Triangle $ABD$ with $C$ on $AD$: $AB = 10$~cm, $BD \perp CD$ with $CD = 8$~cm, $BC = x$, $AD = 17$~cm.)

\bigskip
\noindent\textbf{215.} Solve the following triangle when $m\angle B = 90^{\circ}$: $m\angle C = 45^{\circ}$, $a = 8$ cm

\bigskip
\noindent\textbf{216.} Solve the following triangle when $m\angle B = 90^{\circ}$: $m\angle A = 60^{\circ}$, $c = 4$ cm

\bigskip
\noindent\textbf{217.} Solve the following triangle when $m\angle B = 90^{\circ}$: $m\angle A = 60^{\circ}$, $c = 4$ cm

\bigskip
\noindent\textbf{218.} Let $Q$ and $R$ be two points on the same bank of a canal. The point $P$ is placed on the other bank, straight across from point $R$. Find the width of the canal and the angle $PQR$. \textit{(figure given)}

\bigskip
\noindent\textbf{219.} A ladder is placed along a wall such that the foot of the ladder is 2~m away from the wall. If the length of the ladder is 8~m, find the height of the wall.

\bigskip
\noindent\textbf{220.} The diagonal of a rectangular field $ABCD$ is $(x+9)$~m and the sides are $(x+7)$~m and $x$~m. Find the value of $x$.

\bigskip
\noindent\textbf{221.} Solve triangle $ABC$, in which $m\angle B = 90^{\circ}$, $m\angle A = 30^{\circ}$, $a = 2$.

\bigskip
\noindent\textbf{222.} Solve triangle $ABC$, when $a = \sqrt{2}$~cm, $c = 1$~cm, $m\angle B = 90^{\circ}$.

\bigskip
\noindent\textbf{223.} An isosceles triangle has a vertical angle of $120^{\circ}$ and a base 10~cm long. Find the length of its altitude.

\bigskip
\noindent\textbf{224.} A lighthouse tower is 150~m high from the sea level. The angle of depression from the top of the tower to a ship is $60^{\circ}$. Find the distance between the ship and the tower.

\bigskip
\noindent\textbf{225.} The measure of the angle of elevation of the top of a pole is $15^{\circ}$ from a point on the ground. On walking 100~m towards the pole, the measure of the angle is found to be $30^{\circ}$. Find the height of the pole.

\bigskip
\noindent\textbf{226.} The measure of the angle of elevation of the top of a cliff is $25^{\circ}$. On walking 100~m towards the cliff, the measure of the angle of elevation of the top is $45^{\circ}$. Find the height of the cliff.

\bigskip
\noindent\textbf{227.} From the top of a hill 300~m high, the measure of the angle of depression of a point on the nearer shore of a river is $70^{\circ}$, and the measure of the angle of depression of a point directly across the river is $50^{\circ}$. Find the width of the river. How far is the river from the foot of the hill?

\bigskip
\noindent\textbf{228.} A kite has 120~m of string attached to it when at an angle of elevation of $50^{\circ}$. How far is it above the hand holding it? (Assume that the string is stretched.)

\bigskip
\noindent\textbf{229.} From the top of a lookout tower, the angle of depression of a building on the ground level is $45^{\circ}$. How far is a man on the ground from the tower, if the height of the tower is 30~m?

\bigskip
\noindent\textbf{230.} The angle of elevation of the top of a pole 40~m high is $60^{\circ}$ when seen from a point on the ground level. Find the distance of the point from the foot of the pole.

\bigskip
\noindent\textbf{231.} Prove the following identity:
\[
\sin\theta(\csc\theta - \sin\theta) = \frac{1}{\sec^2\theta}
\]

\bigskip
\noindent\textbf{232.} Prove the following identity:
\[
\frac{\csc\theta - \sec\theta}{\csc\theta + \sec\theta} = \frac{1-\cos\theta}{1+\tan\theta}
\]

\bigskip
\noindent\textbf{233.} Prove the following identity:
\[
\frac{\cos\theta+\sin\theta}{\cos\theta-\sin\theta} + \frac{\cos\theta-\sin\theta}{\cos\theta+\sin\theta} = \frac{2}{1-2\sin^2\theta}
\]

\bigskip
\noindent\textbf{234.} Prove the following identity:
\[
\frac{1+\cos\theta}{1-\cos\theta} = (\csc\theta + \cot\theta)^2
\]

\bigskip
\noindent\textbf{235.} If $\tan\theta = \dfrac{3}{\sqrt{2}}$, then find the remaining trigonometric ratios when $\theta$ lies in the first quadrant.

\bigskip
\noindent\textbf{236.} Show that the points $A(0,2)$, $B(\sqrt{3}, 1)$ and $C(0,-2)$ are vertices of a right triangle.

\bigskip
\noindent\textbf{237.} Show that the points $A(3,1)$, $B(-2,-3)$ and $C(2,2)$ are vertices of an isosceles triangle.

\bigskip
\noindent\textbf{238.} Show that the points $A(5,2)$, $B(-2,3)$, $C(-3,-4)$ and $D(4,-5)$ are vertices of a parallelogram.

\bigskip
\noindent\textbf{239.} Find $h$ such that the points $A(\sqrt{3}, -1)$, $B(0,2)$ and $C(h,-2)$ are vertices of a right triangle with the right angle at vertex $A$.

\bigskip
\noindent\textbf{240.} Find $h$ such that $A(-1,h)$, $B(3,2)$ and $C(7,3)$ are collinear.

\bigskip
\noindent\textbf{241.} The points $A(-5,-2)$ and $B(5,-4)$ are ends of a diameter of a circle. Find the centre and radius of the circle.

\bigskip
\noindent\textbf{242.} Find $h$ such that the points $A(h,1)$, $B(2,7)$ and $C(-6,-7)$ are vertices of a right triangle with the right angle at vertex $A$.

\bigskip
\noindent\textbf{243.} A quadrilateral has the points $A(9,3)$, $B(-7,7)$, $C(-3,-7)$ and $D(5,-5)$ as its vertices. Find the midpoints of its sides. Show that the figure formed by joining the midpoints consecutively is a parallelogram.

\bigskip
\noindent\textbf{244.} Find the distance between the points $C(-4,-2)$ and $D(0,9)$.

\bigskip
\noindent\textbf{245.} Using slopes, show that the triangle with vertices $A(6,1)$, $B(2,7)$ and $C(-6,-7)$ is a right triangle.

\bigskip
\noindent\textbf{246.} Two pairs of points are given. Find whether the two lines determined by these points are: (i) parallel (ii) perpendicular (iii) none. \\
$(1,-2), (2,4)$ and $(4,1), (-8,2)$

\bigskip
\noindent\textbf{247.} Two pairs of points are given. Find whether the two lines determined by these points are: (i) parallel (ii) perpendicular (iii) none. \\
$(-3,4), (6,2)$ and $(4,5), (-2,-7)$

\bigskip
\noindent\textbf{248.} Find an equation of the perpendicular bisector of the segment joining the points $A(3,5)$ and $B(9,8)$.

\bigskip
\noindent\textbf{249.} Find an equation of the line through $(-4,-6)$ and perpendicular to a line having slope $\dfrac{-3}{2}$.

\bigskip
\noindent\textbf{250.} Find an equation of the line through $(11,-5)$ and parallel to a line with slope $-24$.

\bigskip
\noindent\textbf{251.} Convert the following equation into slope-intercept form, two-intercept form and normal form:
\[
2x - 4y + 11 = 0
\]

\bigskip
\noindent\textbf{252.} Convert the following equation into slope-intercept form, two-intercept form and normal form:
\[
4x + 7y - 2 = 0
\]

\bigskip
\noindent\textbf{253.} Convert the following equation into slope-intercept form, two-intercept form and normal form:
\[
15y - 8x + 3 = 0
\]

\bigskip
\noindent\textbf{254.} Check whether the following two lines are: (i) parallel (ii) perpendicular (iii) neither parallel nor perpendicular:
\[
2x + y - 3 = 0 \, ; \quad 4x + 2y + 5 = 0
\]

\bigskip
\noindent\textbf{255.} Check whether the following two lines are: (i) parallel (ii) perpendicular (iii) neither parallel nor perpendicular:
\[
3y = 2x + 5 \, ; \quad 3x + 2y - 8 = 0
\]

\bigskip
\noindent\textbf{256.} Check whether the following two lines are: (i) parallel (ii) perpendicular (iii) neither parallel nor perpendicular:
\[
4y + 2x - 1 = 0 \, ; \quad x - 2y - 7 = 0
\]

\bigskip
\noindent\textbf{257.} Find an equation of the line through $(-4,7)$ and parallel to the line $2x - 7y + 4 = 0$.

\bigskip
\noindent\textbf{258.} Find an equation of the line through $(5,-8)$ and perpendicular to the join of $A(-15,-8)$, $B(10,7)$.

\bigskip
\noindent\textbf{259.} Find the equation of the straight line if its slope is 2 and $y$-intercept is 5.

\bigskip
\noindent\textbf{260.} Find the equation of the straight line if it is perpendicular to a line with slope $-6$ and its $y$-intercept is $\dfrac{4}{3}$.

\bigskip
\noindent\textbf{261.} Write down an equation of the straight line passing through $(5,1)$ and parallel to a line passing through the points $(0,1)$, $(7,15)$.

\bigskip
\noindent\textbf{262.} The length of the perpendicular from the origin to a line is 5 units, and the inclination of this perpendicular is $120^{\circ}$. Find the slope and $y$-intercept of the line.

\bigskip
\noindent\textbf{263.} An architect is designing a park with two buildings located at $(10,8)$ and $(4,3)$ on the grid. Calculate the straight-line distance between the buildings. Assume the coordinates are in meters.

\bigskip
\noindent\textbf{264.} A delivery driver needs to calculate the distance between two delivery locations. One location is at $(7,2)$ and the other at $(12,10)$ on the city grid map, where each unit represents a kilometre. What is the distance between the two locations?

\bigskip
\noindent\textbf{265.} A ship is navigating from port $A$ located at $(12^{\circ}N, 65^{\circ}W)$ to port $B$ at $(20^{\circ}N, 45^{\circ}W)$. If the ship travels along the shortest path on the surface of the Earth, calculate the straight-line distance between the points.

\bigskip
\noindent\textbf{266.} Farah is fencing around a rectangular field with corners at $(0,0)$, $(0,5)$, $(8,5)$ and $(8,0)$. How much fencing material will she need to cover the entire perimeter of the field?

\bigskip
\noindent\textbf{267.} An airplane is flying from city $X$ at $(40^{\circ}N, 100^{\circ}W)$ to city $Y$ at $(50^{\circ}N, 80^{\circ}W)$. Using coordinate geometry, calculate the shortest distance between these two cities.

\bigskip
\noindent\textbf{268.} A land surveyor is marking out a rectangular plot of land with corners at $(3,1)$, $(3,6)$, $(8,6)$ and $(8,1)$. Calculate the perimeter.

\bigskip
\noindent\textbf{269.} A landscaper needs to install a fence around a rectangular garden. The garden has its corners at the coordinates $A(0,0)$, $B(5,0)$, $C(5,3)$ and $D(0,3)$. How much fencing is required?

\bigskip
\noindent\textbf{270.} A landscaper is designing a triangular garden with corners at points $A(2,3)$, $B(5,7)$ and $C(6,2)$. Calculate the lengths of the sides of the triangle.

\bigskip
\noindent\textbf{271.} Find the equation of the line in the form $y = mx + c$ that passes through the points $(3,7)$ and $(5,11)$.

\bigskip
\noindent\textbf{272.} A drone is flying from point $(2,3)$ to point $(10,15)$ on the grid. Calculate the gradient of the line along which the drone is flying and the total distance travelled.

\bigskip
\noindent\textbf{273.} Formulate a simple deductive proof of the following algebraic expression; prove that the L.H.S. is equal to the R.H.S.:
\[
(x-4)^2 + 9 = x^2 - 8x + 25
\]

\bigskip
\noindent\textbf{274.} Formulate a simple deductive proof of the following algebraic expression; prove that the L.H.S. is equal to the R.H.S.:
\[
(x+1)^2 - (x-1)^2 = 4x
\]

\bigskip
\noindent\textbf{275.} Formulate a simple deductive proof of the following algebraic expression; prove that the L.H.S. is equal to the R.H.S.:
\[
(x+5)^2 - (x-5)^2 = 20x
\]

\bigskip
\noindent\textbf{276.} Prove the following by justifying each step:
\[
\frac{4+16x}{4} = 1 + 4x
\]

\bigskip
\noindent\textbf{277.} Prove the following by justifying each step:
\[
\frac{6x^2 + 18x}{3x^2 - 9} = \frac{2x}{x-3}
\]

\bigskip
\noindent\textbf{278.} Prove the following by justifying each step:
\[
\frac{x^2 + 7x + 10}{x^2 - 3x - 9} = \frac{x+5}{x-5}
\]

\bigskip
\noindent\textbf{279.} Prove the following statement: If $x$ is an odd integer, then show that $x^2 - 4x + 6$ is odd.

\bigskip
\noindent\textbf{280.} Prove the following statement: If $x$ is an even integer, then show that $x^2 + 2x + 4$ is even.

\bigskip
\noindent\textbf{281.} Prove that for any two non-empty sets $A$ and $B$, $(A \cap B)' = A' \cup B'$.

\bigskip
\noindent\textbf{282.} Construct the truth table of: $[(p \rightarrow q) \wedge p]$ and $[(p \rightarrow q) \wedge p] \rightarrow q$.

\bigskip
\noindent\textbf{283.} Prove that for any two non-empty sets $A$ and $B$, $(A \cup B)' = A' \cap B'$.

\bigskip
\noindent\textbf{284.} Prove that
\[
\frac{a}{b} + \frac{c}{d} = \frac{ad+bc}{bd}
\]
where $a, b, c$ and $d$ are non-zero real numbers.

\bigskip
\noindent\textbf{285.} Prove that: $(x+1)^2 + 7 = x^2 + 2x + 8$

\bigskip
\noindent\textbf{286.} Prove that
\[
\frac{45x+15}{15} = 3x+1
\]
by justifying each step.



\begin{wq}[start=287]

\item A plank is placed straight upstairs; it is 20 cm wide and 16 cm deep. A rectangular box of height 8 cm and width $x$ cm is placed on a stair under the plank. Find the value of $x$. \emph{[figure -- see original paper]}

\item A man who is 1.8 m tall casts a shadow of 0.76 m in length. If at the same time a telephone pole casts a 3 m shadow, find the height of the pole.

\item Draw an isosceles trapezoid $ABCD$ where $\seg{AB}\parallel \seg{CD}$ and $m(\seg{AB})>m(\seg{CD})$. Draw diagonals $\seg{AC}$ and $\seg{BD}$, intersecting at $E$. Prove that $\triangle ABE$ is similar to $\triangle CDE$. If $m(\seg{AB})=8$ cm, $m(\seg{CD})=4$ cm, and $m(\seg{AE})=3$ cm, find the length of $\seg{CE}$.

\item A regular dodecagon has its side lengths decreased by a factor of $\dfrac{1}{\sqrt{2}}$. If the perimeter of the original dodecagon is 72 cm, what is the side length of the scaled dodecagon?

\item In the given figure, $\triangle ABC \sim \triangle DEF$, with $m(\seg{AB})=12$ cm, $m(\seg{AC})=20$ cm, $m(\seg{BC})=16$ cm. In $\triangle DEF$, $m(\seg{DE})=6$ cm. Find $m(\seg{DF})$ and $m(\seg{EF})$. \emph{[figure -- see original paper]}

\item Find whether the parallelograms are similar, given that one of the angles between sides is $45^{\circ}$ in both parallelograms.

\item In the triangles $XBC$ and $XDE$, find the value of $x$ and $y$. \emph{[figure -- see original paper]}

\item The quadrilateral $ABCD$ has side lengths $m(\seg{AB})=5$ cm, $m(\seg{BC})=8$ cm, $m(\seg{CD})=10$ cm, $m(\seg{AD})=12$ cm, and its angles are $m\angle A=90^{\circ}$, $m\angle B=120^{\circ}$ and $m\angle C=90^{\circ}$. Quadrilateral $EFGH$ has side lengths $m(\seg{EF})=10$ cm, $m(\seg{FG})=16$ cm, $m(\seg{GH})=20$ cm, $m(\seg{EH})=24$ cm and its angles are $m\angle E=90^{\circ}$, $m\angle F=120^{\circ}$ and $m\angle H=60^{\circ}$. Prove that quadrilateral $ABCD$ is similar to quadrilateral $EFGH$.

\item Given that the area of $\triangle ABC = 36\ \text{cm}^2$ and $m(\seg{AB})=6$ cm, $m(\seg{BD})=4$ cm, find the area of $\triangle ADE$.

\item Given that the area of $\triangle ABC = 36\ \text{cm}^2$ and $m(\seg{AB})=6$ cm, $m(\seg{BD})=4$ cm, find the area of trapezium $BCED$.

\item Given that $\triangle ABC$ and $\triangle DEF$ are similar, with a scale factor of $k=3$. If the area of $\triangle ABC$ is $50\ \text{cm}^2$, find the area of triangle $\triangle DEF$.

\item Quadrilaterals $ABCD$ and $EFGH$ are similar, with a scale factor of $k=\dfrac{1}{4}$. If the area of quadrilateral $ABCD$ is $64\ \text{cm}^2$, find the area of quadrilateral $EFGH$.

\item The areas of two similar triangles are $16\ \text{cm}^2$ and $25\ \text{cm}^2$. What is the ratio of a pair of corresponding sides?

\item The areas of two similar triangles are $144\ \text{cm}^2$ and $81\ \text{cm}^2$. If the base of the larger triangle is 30 cm, find the corresponding base of the smaller triangle.

\item A regular heptagon is inscribed in a larger regular heptagon, and each side of the larger heptagon is 1.7 times the side of the smaller heptagon. If the area of the smaller heptagon is $100\ \text{cm}^2$, find the area of the larger heptagon.

\item The ratio of the heights of two similar cylindrical water tanks is $5:3$.
\begin{enumerate}[label=(\arabic*)]
\item If the surface area of the larger tank is 250 square metres, find the surface area of the smaller tank.
\item If the volume of the smaller tank is 270 cubic metres, find the volume of the larger tank.
\end{enumerate}

\item The ratio of the corresponding lengths of two similar conical cans is $3:2$.
\begin{enumerate}[label=(\roman*)]
\item The larger conical can has a surface area of $96\ \text{m}^2$. Find the surface area of the smaller conical can.
\item The smaller conical can has a volume of $240\ \text{m}^3$. Find the volume of the larger conical can.
\end{enumerate}

\item The ratio of the corresponding lengths of two similar cylindrical cans is $3:2$. The larger cylindrical can has a surface area of 67.5 square metres. Find the surface area of the smaller cylindrical can.

\item The ratio of the corresponding lengths of two similar cylindrical cans is $3:2$. The smaller cylindrical can has a volume of 132 cubic metres. Find the volume of the larger cylindrical can.

\item A solid cone $C$ is cut into two pieces $A$ and $B$ with sloping edges 6 cm and 4 cm. Find the ratio of the diameters of the bases of cones $A$ and $C$. \emph{[figure -- see original paper]}

\item Using the same cone $C$ (cut into pieces $A$ and $B$ with sloping edges 6 cm and 4 cm), find the ratio of the volumes of cones $A$ and $C$.

\item Using the same cone $C$, find the ratio of the areas of the bases of cones $A$ and $C$.

\item Using the same cone $C$, if the volume of cone $A$ is $72\ \text{cm}^3$, find the volume of solid $B$.

\item The mass of a sack of rice is 50 kg and its height is 4 m. Find the mass of a similar sack of rice with a height of 6 m.

\item A shape is created by cutting a square in half diagonally and then attaching a right-angled triangle to the hypotenuse of each half. Explain why this shape can tessellate, and calculate the interior angle of the new shape.

\item A tessellation is created by repeatedly reflecting a basic shape. The basic shape is a right-angled triangle with sides of lengths 3, 4 and 5 units. Find the minimum number of reflections needed to create a tessellation that covers a square with an area of 3600 square units.

\item A tessellation is created using regular hexagons. Each hexagon has a side length of 5 cm. Find the total area of the tessellation if it consists of 25 hexagons, and the total perimeter of the outer edge of the tessellation, assuming it forms a perfect hexagon.

\item In a parallelogram $ABCD$, if $m\angle DAB=70^{\circ}$, find the measures of all other angles in the parallelogram.

\item A window has a trapezoidal shape with parallel sides of 3 m and 1.5 m and a height of 2 m. Find the area of the window.

\item Falak plans to tile a square patio with an area of 100 square metres. He decides to use both square tiles and triangular tiles, each with an area of 0.25 square metres. If 60\% of the tiles will be square and 40\% will be triangular, how many tiles of each shape are needed?

\item Ali wants to create a floor design that uses regular hexagons (each with a side length of 1 metre) and equilateral triangles (each with a side length of 1 metre) to cover a rectangular area measuring 10 m by 5 m. Find how many hexagons and triangles Ali will need to complete the tessellation.

\item A parallelogram-shaped room has a base of 10 metres and a height of 8 m. Babar wants to carpet the room using rolls that cover $20\ \text{m}^2$ each. How many rolls of carpet does he need?

\item A tessellation is created using a combination of regular pentagons and decagons. Find the sum of the angles at a vertex where a pentagon and a decagon meet.

\item Plot the graph of the following cubic function: $y=x^3+2x^2-5x-6$; $-3.5\le x\le 2.5$

\item Plot the graph of the following quadratic function: $y=x^2+x-2$

\item Plot the graph of the following cubic function: $y=x^3+3x^2+2x$; $-2.5\le x\le 0.5$

\item Plot the graph of the following quadratic function: $y=5x^2-2x-3$

\item Plot the graph of $y=2x^2-4x+3$ for $x$ from $-1$ to $3$. Draw the tangent at $(2,3)$ and find the gradient.

\item Plot the graph of $y=3x^2+x+1$ and draw the tangent at $(1,5)$. Also find the gradient of the tangent line at this point.

\item The strength of students in a school was 1000 in 2016. If the strength decays according to the equation $S=1000e^{-t}$, where $S$ is the number of students at time $t$, graph the given equation from $t=0$ (in 2016) to $t=9$ (in 2025).

\item Using the equation from the previous question, estimate from the graph the student strength in 2019 and 2023.

\item A newspaper agency has a fixed cost of Rs.\ 70 per edition and marginal printing and distribution costs of Rs.\ 40 per copy. The profit function is $p(x)=0.10x+70$, where $x$ is the number of newspapers. Plot the graph and find the profit for 500 newspapers.

\item Ali manufactures expensive shirts for sale to a school. Its cost (in rupees) for $x$ shirts is $C(x)=1500+10x+0.2x^2$, $0\le x\le 150$. Plot the graph and find the cost of 200 shirts.

\item The population of a village was 753 in 2010. If the population grows according to the equation $p=753e^{0.03t}$, where $p$ is the number of persons in the population at time $t$, graph the population equation for $t=0$ (in 2010) to $t=30$ (in 2040).

\item Using the equation from the previous question, estimate from the graph the population (i) in 2020 and (ii) in 2030.

\item Sketch the graph of $y=x^2-3x-2$ for values of $x$ from $-8$ to $8$, draw a tangent line at $x=-6$ and determine the gradient.

\item Majid's salary $S(x)$ in rupees is based on the formula $S(x)=25000+1500x$, where $x$ is the number of years he has worked. Sketch and interpret the graph of the salary function for $0\le x\le 10$.

\item A company manufactures footballs. The cost of manufacturing $x$ footballs is $C(x)=90{,}000+600x$. The revenue from selling is $R(x)=1{,}800x$. Find the break-even point and determine the profit or loss when 200 footballs are sold. Draw the graph of both functions and identify the break-even point.

\item Sales of a new magazine are expected to grow according to the equation $S=200000(1-e^{-0.05t})$, where $t$ is given in weeks. Plot the graph of sales for the first 50 weeks.

\item Using the equation from the previous question, calculate the number of magazines sold when $t=5$ and $t=35$.

\item Plot the graph of $y=\dfrac{1}{2}(x+4)(x-1)(x-3)$ from $-5$ to $4$.

\item The supply and demand functions for a particular market are given by $P_s=Q^2+5$ and $P_d=Q^2-10Q$, where $P$ represents price and $Q$ represents quantity. Sketch the graph of each function over the interval $Q=-20$ to $Q=20$.

\item A television manufacturer makes 40-inch LEDs. The cost of manufacturing $x$ LEDs is $C(x)=60{,}000+250x$ and the revenue from selling $x$ LEDs is $R(x)=1200x$. Find the break-even point and find the profit or loss when 100 LEDs are sold. Identify the break-even point graphically.

\item Verify that the angle bisectors of $\triangle ABC$ are concurrent with the following measurements: $m(\seg{AB})=4.5$ cm, $m\angle A=45^{\circ}$ and $m(\seg{AC})=5.3$ cm.

\item Verify that the angle bisectors of $\triangle ABC$ are concurrent with the following measurements: $m(\seg{AB})=6$ cm, $m\angle A=150^{\circ}$ and $m\angle B=60^{\circ}$.

\item Construct $\triangle ABC$ with the given measurements and verify that the perpendicular bisectors of the triangle are concurrent: $m(\seg{AB})=5$ cm, $m(\seg{BC})=6$ cm and $m(\seg{AC})=7$ cm.

\item Construct $\triangle ABC$ with the given measurements and verify that the perpendicular bisectors of the triangle are concurrent: $m(\seg{AB})=7.1$ cm, $m\angle B=135^{\circ}$ and $m(\seg{BC})=6.5$ cm.

\item Construct $\triangle LMN$ with the following measurements and verify that the medians of the triangle are concurrent: $m(\seg{LM})=4.9$ cm, $m\angle L=51^{\circ}$ and $m\angle M=38^{\circ}$.

\item Construct $\triangle LMN$ with the following measurements and verify that the medians of the triangle are concurrent: $m(\seg{MN})=4.8$ cm, $m\angle N=30^{\circ}$ and $m(\seg{LM})=8.1$ cm.

\item Verify that the angle bisectors of $\triangle ABC$ are concurrent with the following measurements: $m(\seg{AB})=4.5$ cm, $m\angle A=45^{\circ}$ and $m(\seg{AC})=5.3$ cm.

\item Given the measurements of $\triangle DEF$: $m(\seg{DE})=4.8$ cm, $m(\seg{EF})=4$ cm and $m\angle E=45^{\circ}$, draw the altitudes of $\triangle DEF$ and find the orthocenter.

\item Construct the following triangle and find whether there exists any ambiguous case: $\triangle BCD$; $m(\seg{BC})=5$ cm, $m\angle B=62^{\circ}$ and $m(\seg{CD})=4.7$ cm.

\item Construct the following triangle and find whether there exists any ambiguous case: $\triangle KLM$; $m(\seg{LM})=6$ cm, $m\angle M=42^{\circ}$ and $m(\seg{LN})=5$ cm.

\item Construct a triangle of sides 5.3 cm, 5.9 cm and 6.2 cm.

\item Construct a triangle $BCD$ in which the measures of two sides are 5.5 cm and 4.2 cm and the measure of their included angle is $60^{\circ}$.

\item Draw a triangle $CDE$ when $m(\seg{DE})=4.3$ cm, $m\angle D=30^{\circ}$ and $m\angle E=120^{\circ}$.

\item Construct a triangle $DEF$ when $m(\seg{DE})=6$ cm, $m\angle D=110^{\circ}$ and $m(\seg{EF})=9$ cm.

\item Construct triangles $DEF$ and $DEF'$ when $m(\seg{DE})=6$ cm, $m\angle D=30^{\circ}$ and $m(\seg{EF})=3.6$ cm.

\item In the above example, repeat the construction if we take $m(\seg{EF})=3$ cm.

\item In the above example, repeat the construction if we take $m(\seg{EF})=2.5$ cm.

\end{wq}



\noindent\textbf{357.} Draw the perpendicular bisectors of the triangle $EFG$ with $m\overline{EF}=5$~cm, $m\overline{FG}=2.5$~cm and $m\overline{EG}=4.3$~cm.

\bigskip
\noindent\textbf{358.} Draw the angle bisector of a triangle $FGH$ if: $m\overline{FG}=5.2$, $m\overline{GH}=4.1$ and $m\angle FGH = 120^{\circ}$.

\bigskip
\noindent\textbf{359.} Construct triangle $GHI$ in which $m\overline{GH}=5.7$, $m\angle G = 68^{\circ}$ and $m\angle G = 50^{\circ}$. Prove that the altitudes of $\triangle GHI$ are concurrent.

\bigskip
\noindent\textbf{360.} Construct $\triangle CDE$ with $m\overline{CD}=7.6$~cm, $m\angle D = 45^{\circ}$ and $m\overline{DE}=5.9$~cm. Draw the locus of all points which are equidistant from $C$ and $D$. Mark the point $X$ where the two loci intersect.

\bigskip
\noindent\textbf{361.} Two points $A$ and $B$ are 8.2~cm apart. Construct the locus of points 5~cm from point $A$.

\bigskip
\noindent\textbf{362.} Construct the locus of points 2.2~cm from line segment $CD$ of measure 5.7~cm.

\bigskip
\noindent\textbf{363.} Construct an angle $ABC = 105^{\circ}$. Construct the locus of a point $P$ which moves such that it is equidistant from $\overline{BA}$ and $\overline{BC}$.

\bigskip
\noindent\textbf{364.} Two points $E$ and $F$ are 5.4~cm apart. Construct the locus of a point $P$ which moves such that it is equidistant from $E$ and $F$.

\bigskip
\noindent\textbf{365.} An island has two main cities $A$ and $B$, 8~km apart. Kashif lives on the island exactly 6.8~km from city $A$ and exactly 7.3~km from city $B$. Mark with a cross the point on the island where Kashif could live.

\bigskip
\noindent\textbf{366.} Construct $\triangle CDE$ with $m\overline{CD}=7.6$~cm, $m\angle D = 45^{\circ}$ and $m\overline{DE}=5.9$~cm. Draw the locus of all points which are equidistant from $\overline{CD}$ and $\overline{CE}$.

\bigskip
\noindent\textbf{367.} Construct a triangle $LMN$ with $m\overline{LM}=7$~cm, $m\angle L = 70^{\circ}$ and $m\angle M = 45^{\circ}$. Find a point within triangle $LMN$ which is equidistant from $L$ and $M$, and 3~cm from $L$.

\bigskip
\noindent\textbf{368.} Construct a right-angled triangle $RST$ with $m\overline{RS}=6.8$~cm, $m\angle S = 90^{\circ}$ and $m\overline{ST}=7.5$~cm. Find a point within triangle $RST$ which is equidistant from $\overline{RS}$ and $\overline{RT}$, and 4.5~cm from $R$.

\bigskip
\noindent\textbf{369.} Construct a rectangle $UVWX$ with $m\overline{UV}=7.2$~cm and $m\overline{VW}=5.6$~cm. Draw the locus of points at a distance of 2~cm from $\overline{UV}$ and 3.5~cm from $W$.

\bigskip
\noindent\textbf{370.} Two cell towers are located at points $A$ and $B$ on a coordinate plane. A GPS-enabled device, positioned somewhere on the plane, receives signals from both towers. To ensure accurate navigation, the device is placed equidistant from both towers to estimate its position. Draw this locus of navigation.

\bigskip
\noindent\textbf{371.} Epidemiologists use loci to determine infection zones, especially for contagious diseases, to predict the spread and take containment measures. In the case of a disease outbreak, authorities might determine a quarantine zone within 10~km of the infection source. Draw the locus of all points 10~km from the source, defining the quarantine area to monitor and control the disease's spread.

\bigskip
\noindent\textbf{372.} There is a treasure buried somewhere on an island. The treasure is 24~km from $A$ and equidistant from $B$ and $C$. Using a scale of 1~cm to represent 10~km, find where the treasure could be buried.

\bigskip
\noindent\textbf{373.} There is an apple tree at a distance of 90~m from a banana tree in the garden of Sara's house. Sara wants to plant a mango tree $M$ which is 64~m from the apple tree and between 54~m and 82~m from the banana tree. Using a scale of 1~cm to represent 10~m, find the points where the mango tree should be planted.

\bigskip
\noindent\textbf{374.} Construct a rectangle $ABCD$ with $m\overline{AB}=5$~cm and $m\overline{BC}=3.2$~cm. Draw the locus of all points which are at a distance of 3.1~cm from point $A$. Label the point $P$ inside the rectangle which is 3.1~cm from point $A$ and equidistant from $A$ and $B$.

\bigskip
\noindent\textbf{375.} Construct a rectangle $ABCD$ with $m\overline{AB}=5$~cm and $m\overline{BC}=3.2$~cm. Draw the locus of all points which are equidistant from $A$ and $B$. Label the point $P$ inside the rectangle which is 3.1~cm from point $A$ and equidistant from $A$ and $B$.

\bigskip
\noindent\textbf{376.} Construct an isosceles triangle $DEF$ with vertical angle $80^{\circ}$ at $E$ and $m\overline{EF}=m\overline{DE}=4.8$~cm. Draw the locus of all points which are at a distance of 2.8~cm from point $E$. Label the point $X$ inside the triangle which is 2.8~cm from point $E$ and equidistant from $\overline{DE}$ and $\overline{EF}$.

\bigskip
\noindent\textbf{377.} Construct an isosceles triangle $DEF$ with vertical angle $80^{\circ}$ at $E$ and $m\overline{EF}=m\overline{DE}=4.8$~cm. Draw the locus of all points which are equidistant from $\overline{DE}$ and $\overline{EF}$. Label the point $X$ inside the triangle which is 2.8~cm from point $E$ and equidistant from $\overline{DE}$ and $\overline{EF}$.

\bigskip
\noindent\textbf{378.} Construct a right-angled triangle with side lengths 6~cm, 8~cm and 10~cm.

\bigskip
\noindent\textbf{379.} Construct a triangle $ABC$ with $m\overline{AB}=5.3$~cm, $m\angle A = 30^{\circ}$ and $m\angle B = 120^{\circ}$. Draw the locus of all points which are equidistant from $A$ and $B$.

\bigskip
\noindent\textbf{380.} Construct a triangle with $m\overline{DE}=7.3$~cm, $m\angle D = 42^{\circ}$ and $m\overline{EF}=5.4$~cm.

\bigskip
\noindent\textbf{381.} Construct a triangle $XYZ$ with $m\overline{YX}=8$~cm, $m\overline{YZ}=7$~cm and $m\overline{XZ}=6.5$~cm. Draw the locus of all points which are equidistant from $\overline{XY}$ and $\overline{XZ}$.

\bigskip
\noindent\textbf{382.} Construct a triangle $FGH$ such that $m\overline{FG}=m\overline{GH}=6.4$~cm, $m\angle G = 122^{\circ}$. Draw the locus of all points which are equidistant from $F$ and $G$.

\bigskip
\noindent\textbf{383.} Construct a triangle $FGH$ such that $m\overline{FG}=m\overline{GH}=6.4$~cm, $m\angle G = 122^{\circ}$. Draw the locus of all points which are equidistant from $\overline{FG}$ and $\overline{GH}$.

\bigskip
\noindent\textbf{384.} Construct a triangle $FGH$ such that $m\overline{FG}=m\overline{GH}=6.4$~cm, $m\angle G = 122^{\circ}$. Draw the locus of all points and mark the point where the two loci intersect.

\bigskip
\noindent\textbf{385.} Two houses $Q$ and $R$ are 73~m apart. Using a scale of 1~cm to represent 10~m, construct the locus of a point $P$ which moves such that it is at a distance of 32~m from $Q$.

\bigskip
\noindent\textbf{386.} Two houses $Q$ and $R$ are 73~m apart. Using a scale of 1~cm to represent 10~m, construct the locus of a point $P$ which moves such that it is at a distance of 48~m from the line joining $Q$ and $R$.

\bigskip
\noindent\textbf{387.} A field is in the form of a rectangle $ABCD$ with $m\overline{AB}=70$~m and $m\overline{BC}=60$~m. Construct the rectangle $ABCD$ using a scale of 1~cm to represent 10~m. Show the region inside the field which is less than 30~m from $C$ and farther than 25~m from $\overline{AB}$.

\bigskip
\noindent\textbf{388.} From the table given below, draw a frequency polygon on histogram for the given frequency distribution:

\begin{center}
\begin{tabular}{|c|c|c|c|c|c|c|}
\hline
Weight (kg) & 10--14 & 15--19 & 20--24 & 25--29 & 30--34 & 35--39 \\
\hline
Frequency ($f$) & 06 & 17 & 23 & 30 & 22 & 13 \\
\hline
\end{tabular}
\end{center}

\bigskip
\noindent\textbf{389.} The marks obtained by the students of Grade 10 in a mathematics test were grouped into the following frequency distribution:

\begin{center}
\begin{tabular}{|c|c|c|c|c|c|c|}
\hline
Marks & 35--37 & 38--44 & 45--54 & 55--61 & 62--67 & 68--72 \\
\hline
Frequency & 1 & 2 & 16 & 13 & 9 & 3 \\
\hline
\end{tabular}
\end{center}

Draw a histogram for the above distribution.

\bigskip
\noindent\textbf{390.} Make a frequency polygon on histogram for the following grouped data:

\begin{center}
\begin{tabular}{|c|c|c|c|c|c|c|}
\hline
Class limits & 5--8 & 8--12 & 12--20 & 20--25 & 25--27 & 27--32 \\
\hline
Frequency ($f$) & 2 & 12 & 25 & 32 & 14 & 5 \\
\hline
\end{tabular}
\end{center}

\bigskip
\noindent\textbf{391.} For a school staff, the following expenditure (rupees in hundreds) is required for the repair of chairs:
\[
\begin{array}{cccccccccc}
145 & 152 & 153 & 156 & 158 & 160 & 146 & 152 & 155 & 159 \\
161 & 163 & 165 & 147 & 148 & 151 & 154 & 156 & 158 & 160 \\
144 & 167 & 151 & 150 & 152 & 149 & 145 & 153 & 152 & 155
\end{array}
\]
Prepare a frequency distribution by the tally-bar method using 3 as the size of the class limits, and write down the frequencies of the last three classes.

\bigskip
\noindent\textbf{392.} Given below are the weights (in kg) of 30 students of a high school:
\[
\begin{array}{cccccccccc}
30 & 33 & 24 & 21 & 15 & 39 & 37 & 44 & 42 & 33 \\
33 & 28 & 29 & 32 & 31 & 28 & 26 & 32 & 34 & 35 \\
38 & 36 & 41 & 30 & 35 & 41 & 23 & 26 & 18 & 34
\end{array}
\]
Taking 5 as the size of the class limit, prepare a frequency table and construct a frequency polygon.

\bigskip
\noindent\textbf{393.} A group of Grade 10 students obtained the following marks out of 100 marks in an English test:
\[
\begin{array}{cccccccccccc}
58 & 59 & 58 & 33 & 40 & 58 & 45 & 46 & 43 & 45 & 45 \\
50 & 52 & 49 & 50 & 57 & 52 & 55 & 49 & 50 & 62 & 49 \\
48 & 44 & 42 & 47 & 46 & 47 & 46 & 53 & 40 & 44 &
\end{array}
\]
Classify the data into a frequency distribution by the direct method, taking 6 as the size of the class limit. Also find the class limit with the least class frequency, and construct a histogram for the data.

\bigskip
\noindent\textbf{394.} Following are the number of telephone calls made in a week to 30 teachers of a high school:
\[
\begin{array}{cccccccccccccccc}
5 & 8 & 11 & 25 & 13 & 16 & 20 & 17 & 15 & 16 & 30 & 21 & 14 & 18 & 19 \\
6 & 22 & 26 & 15 & 19 & 35 & 29 & 31 & 23 & 25 & 20 & 10 & 9 & 7 & 26
\end{array}
\]
Construct a frequency distribution with the number of classes.

\bigskip
\noindent\textbf{395.} The frequency distribution of ages (in years) of 76 members of a locality is available. Draw a histogram for this data.

\begin{center}
\begin{tabular}{|c|c|c|c|c|c|c|c|}
\hline
Class limits & 2--4 & 4--9 & 9--12 & 12--17 & 17--20 & 20--27 & 27--30 \\
\hline
Frequency ($f$) & 7 & 10 & 18 & 20 & 10 & 7 & 4 \\
\hline
\end{tabular}
\end{center}

\bigskip
\noindent\textbf{396.} The following are the marks obtained by 30 students out of 100 in Mathematics at their final examination. Construct a frequency polygon for the following frequency table:

\begin{center}
\begin{tabular}{|c|c|c|c|c|c|c|}
\hline
Marks & 45--49 & 50--54 & 55--59 & 60--64 & 65--69 & 70--74 \\
\hline
Frequency & 2 & 9 & 4 & 7 & 3 & 5 \\
\hline
\end{tabular}
\end{center}

\bigskip
\noindent\textbf{397.} The marks obtained by students in English are given below:

\begin{center}
\begin{tabular}{|c|c|c|c|c|c|}
\hline
Marks obtained & 15--19 & 20--24 & 25--29 & 30--34 & 35--39 \\
\hline
Frequency & 9 & 18 & 35 & 17 & 5 \\
\hline
\end{tabular}
\end{center}

Find the arithmetic mean of the marks by the direct and short formula.

\bigskip
\noindent\textbf{398.} The marks obtained by students in English are given below:

\begin{center}
\begin{tabular}{|c|c|c|c|c|c|}
\hline
Marks obtained & 15--19 & 20--24 & 25--29 & 30--34 & 35--39 \\
\hline
Frequency & 9 & 18 & 35 & 17 & 5 \\
\hline
\end{tabular}
\end{center}

Find the median of the marks.

\bigskip
\noindent\textbf{399.} Given below is a frequency distribution:

\begin{center}
\begin{tabular}{|c|c|c|c|c|c|}
\hline
Class interval & 5--9 & 10--14 & 15--19 & 20--24 & 25--29 \\
\hline
Frequency & 1 & 8 & 18 & 11 & 2 \\
\hline
\end{tabular}
\end{center}

Find the mode of the frequency distribution.

\bigskip
\noindent\textbf{400.} The monthly attendance of 10 students for their lunch in the hostel is recorded as: $21, 15, 16, 18, 14, 17, 15, 12, 13, 11$. Find the median and mode of the attendance. Also find the mean if $D = A - 20$.

\bigskip
\noindent\textbf{401.} On a prize distribution day, 50 students brought pocket money as under:

\begin{center}
\begin{tabular}{|c|c|c|c|c|c|}
\hline
Rupees & 5--10 & 10--15 & 15--20 & 20--25 & 25--30 \\
\hline
Frequency ($f$) & 12 & 9 & 18 & 7 & 4 \\
\hline
\end{tabular}
\end{center}

Find the median and mode of the above data.

\bigskip
\noindent\textbf{402.} The children Haris, Maham and Minal made the following scores in a game conducted by a group of teachers in the school:

\begin{center}
\begin{tabular}{|l|c|c|c|c|c|}
\hline
Haris scores & 50 & 55 & 70 & 85 & 90 \\
Maham scores & 75 & 60 & 60 & 45 & 53 \\
Minal scores & 80 & 77 & 66 & 42 & 48 \\
\hline
\end{tabular}
\end{center}

It is decided that the candidate who gets the highest average score will be awarded Rs.~1000. Who will get the awarded amount?

\bigskip
\noindent\textbf{403.} As partners, Hafsa and Fatima took part in a quiz programme. They made the following numbers of points: $45, 51, 58, 61, 74, 48, 46$ and $50$. Compute the average number of points using deviation $D = x - 58$.

\bigskip
\noindent\textbf{404.} A person purchased the following food items:

\begin{center}
\begin{tabular}{|l|c|c|}
\hline
Food item & Quantity (in kg) & Cost per kg (in Rs.) \\
\hline
Rice & 10 & 96 \\
Flour & 12 & 48 \\
Ghee & 4 & 190 \\
Sugar & 3 & 49 \\
Mutton & 2 & 650 \\
\hline
\end{tabular}
\end{center}

What is the weighted mean cost of food items per kg?

\bigskip
\noindent\textbf{405.} The arithmetic mean of the ages of 20 boys is 13 years, 4 months and 5 days. Find the sum of their ages. If one of the boys is exactly 15 years old, what is the average age of the remaining boys?

\bigskip
\noindent\textbf{406.} Given below is the frequency distribution derived by making a substitution as $D = X - 20$. Calculate the arithmetic mean:

\begin{center}
\begin{tabular}{|c|c|c|c|c|c|c|c|}
\hline
$D$ & $-6$ & $-4$ & $-2$ & 0 & 2 & 4 & 6 \\
\hline
$f$ & 1 & 3 & 6 & 20 & 26 & 12 & 2 \\
\hline
\end{tabular}
\end{center}

\bigskip
\noindent\textbf{407.} The marks of Jamal in Mathematics in eight monthly tests were $75, 76, 80, 80, 82, 82, 82, 85$. Find the mode of the marks.

\bigskip
\noindent\textbf{408.} Ten students were asked about the number of questions they solved out of 20 questions last week. The records were: $13, 14, 15, 11, 16, 10, 19, 20, 18, 17$. Find the mode of the data.

\bigskip
\noindent\textbf{409.} A survey was conducted among 15 students of a school, asking about their favourite colour. The responses were: Purple, yellow, purple, yellow, yellow, red, blue, green, yellow, yellow, red, blue, yellow, purple, green. Find the mode of the data.

\bigskip
\noindent\textbf{410.} Ten students each from Grade V, Section A and Section B, of a well-reputed school were taken randomly. Their weights were measured in kg and recorded as given below:

\begin{center}
\begin{tabular}{|l|c|c|c|c|c|c|c|c|c|c|}
\hline
Weights (kg) Section A & 30 & 28 & 32 & 29.5 & 35 & 34 & 31 & 33 & 40 & 37.5 \\
\hline
Weights (kg) Section B & 35 & 31.5 & 34.5 & 35 & 32.8 & 38 & 39.5 & 36 & 36.5 & 34 \\
\hline
\end{tabular}
\end{center}

Find the average weight for Section A and Section B.



% ---------------------------------------------------
\noindent\textbf{411.} A government allocates funds of Rs.~200{,}000 to five sectors of a school:
(i) school library: Rs.~35{,}000 \quad (ii) Sports facilities: Rs.~25{,}000 \quad (iii) Parking Area: Rs.~40{,}000 \quad (iv) Room renovation: Rs.~45{,}000 \quad (v) Furniture: Rs.~55{,}000. \\
Find the average of the fund allocation in each sector of a school.

\bigskip
\noindent\textbf{412.} Given below are the marks out of 100 obtained by 100 students in an examination. Find the average marks of the students:

\begin{center}
\begin{tabular}{|c|c|c|c|c|c|c|}
\hline
Marks & 30--35 & 35--40 & 40--45 & 45--50 & 50--55 & 55--60 \\
\hline
No.\ of students & 14 & 16 & 18 & 23 & 18 & 11 \\
\hline
\end{tabular}
\end{center}

\bigskip
\noindent\textbf{413.} Deviations from 12.5 of ten different values are
\[
6,\ -2,\ 3.5,\ 9,\ 8.7,\ -5.5,\ 14,\ 11.3,\ -6.8,\ -4.2
\]
Find the arithmetic mean.

\bigskip
\noindent\textbf{414.} The heights (in inches) of 200 students are recorded in the following frequency distribution. Find the mean height of the students by the short formula:

\begin{center}
\begin{tabular}{|c|c|c|c|c|c|c|c|c|c|c|}
\hline
Height ($x$) (inches) & 51 & 52 & 53 & 54 & 55 & 56 & 57 & 58 & 59 & 60 \\
\hline
Frequency ($f$) & 2 & 5 & 8 & 24 & 55 & 45 & 38 & 16 & 6 & 1 \\
\hline
\end{tabular}
\end{center}

\bigskip
\noindent\textbf{415.} The heights of 100 athletes, measured to the nearest inch, are given in the following table. Find the median:

\begin{center}
\begin{tabular}{|c|c|c|c|c|c|c|c|c|c|}
\hline
Heights (in) & 62.5--63.5 & 63.5--64.5 & 64.5--65.5 & 65.5--66.5 & 66.5--67.5 & 67.5--68.5 & 68.5--69.5 & 69.5--70.5 & 70.5--71.5 \\
\hline
No.\ of Students & 4 & 6 & 10 & 20 & 30 & 13 & 12 & 3 & 2 \\
\hline
\end{tabular}
\end{center}

\bigskip
\noindent\textbf{416.} Following are the weights (in kg) of 50 men. Find the median weight.

\begin{center}
\begin{tabular}{|c|c|c|c|c|c|}
\hline
Weights (kg) & 110--114 & 115--119 & 120--124 & 125--129 & 130--134 \\
\hline
No.\ of men ($f$) & 5 & 12 & 23 & 6 & 4 \\
\hline
\end{tabular}
\end{center}

\bigskip
\noindent\textbf{417.} Following are the heights (in inches) of 40 students in Grade 8. Find the mode of the data.

\begin{center}
\begin{tabular}{|c|c|c|c|c|c|c|}
\hline
Heights (inches) & 48--50 & 50--52 & 52--54 & 54--56 & 56--58 & 58--60 \\
\hline
No.\ of students ($f$) & 5 & 7 & 10 & 9 & 6 & 3 \\
\hline
\end{tabular}
\end{center}

\bigskip
\noindent\textbf{418.} A medicine company started marketing a sample medicine in seven different areas of a city. The company distributed the packets of medicine in each area of the city, and the weight of each area was based on the demand of the medicine. Find the mean and weighted mean of the given data.

\begin{center}
\begin{tabular}{|c|c|c|}
\hline
Area & Number of packets & Weights (kg) \\
\hline
A & 15 & 5 \\
B & 25 & 4 \\
C & 18 & 3 \\
D & 23 & 4 \\
E & 15 & 2 \\
F & 10 & 1 \\
G & 8  & 2 \\
\hline
\end{tabular}
\end{center}

\bigskip
\noindent\textbf{419.} A toy factory sold toys in a month. Consider the following data:

\begin{center}
\begin{tabular}{|c|c|c|c|c|c|}
\hline
Class limits & 10--20 & 20--30 & 30--40 & 40--50 & 50--60 \\
\hline
$f$ & 15 & 28 & 45 & 29 & 20 \\
\hline
\end{tabular}
\end{center}

\begin{enumerate}[label=(\roman*)]
\item Calculate mean, median and mode for the number of toys sold by the factory.
\item Also tell the modal class of the distribution.
\end{enumerate}

\bigskip
\noindent\textbf{420.} Following are the weights of 40 students recorded to the nearest lb.

\[
\begin{array}{cccccccccc}
138 & 164 & 150 & 132 & 144 & 125 & 149 & 157 & 146 & 158 \\
140 & 147 & 136 & 148 & 152 & 144 & 168 & 126 & 138 & 176 \\
163 & 119 & 154 & 165 & 146 & 173 & 142 & 147 & 135 & 153 \\
140 & 135 & 161 & 145 & 135 & 142 & 150 & 156 & 145 & 128
\end{array}
\]

Make a frequency table taking size of class limit as 10. Also draw a frequency polygon of the given data.

\bigskip
\noindent\textbf{421.} The adjoining distribution shows the maximum load (in kg) supported by certain ropes. Find the mean load using the short method:

\begin{center}
\begin{tabular}{|c|c|c|c|c|c|c|}
\hline
Max-Load (kg) & 93--97 & 98--102 & 103--107 & 108--112 & 113--117 & 118--122 \\
\hline
No.\ of ropes & 2 & 5 & 8 & 12 & 6 & 2 \\
\hline
\end{tabular}
\end{center}

\bigskip
\noindent\textbf{422.} Given below is a frequency distribution. Draw a frequency polygon and histogram for the distribution:

\begin{center}
\begin{tabular}{|c|c|c|c|c|c|c|}
\hline
Class limits & 5--9 & 10--14 & 15--19 & 20--24 & 25--29 & 30--34 \\
\hline
Frequency & 1 & 8 & 18 & 11 & 2 & 5 \\
\hline
\end{tabular}
\end{center}

\bigskip
\noindent\textbf{423.} For the following data, find the weighted mean:

\begin{center}
\begin{tabular}{|l|c|c|}
\hline
Item & Quantity & Cost of item (Rs.) \\
\hline
Chair & 20 & 500 \\
Table & 20 & 400 \\
Black board & 10 & 750 \\
Tube light & 25 & 230 \\
Cupboard & 09 & 950 \\
\hline
\end{tabular}
\end{center}

\bigskip
\noindent\textbf{424.} Two partners Mr.\ Aslam and Mrs.\ Kalsoom run a company. In the following data the weekly wages (in Rs.) of employees who work in the company are given:

\begin{center}
\begin{tabular}{|c|c|c|c|c|c|}
\hline
Wages (Rs.) & 600--700 & 700--800 & 800--900 & 900--1000 & 1000--1100 \\
\hline
Employees & 3 & 5 & 7 & 21 & 11 \\
\hline
\end{tabular}
\end{center}

Find the arithmetic mean, median, and standard deviation of the given table.

\bigskip
\noindent\textbf{425.} From the table given below, draw a frequency polygon on histogram for the given frequency distribution:

\begin{center}
\begin{tabular}{|c|c|c|c|c|c|c|}
\hline
Weight (kg) & 50--56 & 57--59 & 60--64 & 65--72 & 73--75 & 76--80 \\
\hline
Frequency ($f$) & 25 & 32 & 40 & 30 & 15 & 8 \\
\hline
\end{tabular}
\end{center}

\bigskip
\noindent\textbf{426.} A principal of a school allocates funds of Rs.~50{,}000 to five different sectors:
(i) Chairs Rs.~1500, \quad (ii) Tables Rs.~12{,}000, \quad (iii) Black boards: Rs.~6{,}000, \quad (iv) Room renovation: Rs.~10{,}000, \quad (v) Gardening: Rs.~7{,}000. \\
Find the average of funds allocation in each sector of the school.

\bigskip
\noindent\textbf{427.} Aslam rolled a dice. What is the probability of getting the number 3 or 4? Also find the probability of not getting the number 3 or 4.

\bigskip
\noindent\textbf{428.} Abdul Hadi labelled cards from 1 to 30 and put them in a box. He selects a card at random. What is the probability that the selected card contains:
\begin{enumerate}[label=(\roman*)]
\item the number 25
\item a number between 17 to 22
\item a number at least 20
\item a number not 27 and not 29
\item a number not between 12 to 15
\end{enumerate}

\bigskip
\noindent\textbf{429.} The probability that Ayesha will pass the examination is 0.85. What will be the probability that Ayesha will not pass the examination?

\bigskip
\noindent\textbf{430.} A card is selected at random from a well-shuffled pack of 52 playing cards. What will be the probability of selecting:
\begin{enumerate}[label=(\roman*)]
\item a queen
\item neither a queen nor a jack
\end{enumerate}

\bigskip
\noindent\textbf{431.} A card is chosen at random from a pack of 52 playing cards. Find the probability of getting:
\begin{enumerate}[label=(\roman*)]
\item a jack
\item no diamond
\end{enumerate}

\bigskip
\noindent\textbf{432.} Taabish tossed a fair coin and rolled a fair dice once. Find the probability of the following events:
\begin{enumerate}[label=(\roman*)]
\item tail on coin and at least 4 on dice.
\item head on coin and the number 2, 3 on dice.
\item head and tail on coin and the number 6 on dice.
\item not tail on coin and the number 5 on dice.
\item not head on coin and the number 5 and 2 on dice.
\end{enumerate}

\bigskip
\noindent\textbf{433.} Let $A$, $B$ and $C$ be missiles fired at a target. If the probabilities of hitting the target are
\[
P(A) = \frac{1}{4}, \qquad P(B) = \frac{3}{7}, \qquad P(C) = \frac{5}{9},
\]
respectively, find the probabilities of:
\begin{enumerate}[label=(\roman*)]
\item missile $A$ does not hit the target.
\item missile $B$ does not hit the target.
\item missile $C$ does not hit the target.
\end{enumerate}

\bigskip
\noindent\textbf{434.} A basket contains 5 blue balls and 8 green balls. Find the probability of selecting at random:
\begin{enumerate}[label=(\roman*)]
\item a blue ball
\item a green ball
\item not a green ball
\end{enumerate}

\bigskip
\noindent\textbf{435.} A researcher collected data on the number of deaths from horse-kicks in Russian Army corps over the years. The table is as follows:

\begin{center}
\begin{tabular}{|c|c|c|c|c|c|c|c|}
\hline
No.\ of Deaths & 0 & 1 & 2 & 3 & 4 & 5 & 6 \\
\hline
Frequency & 60 & 50 & 87 & 40 & 32 & 15 & 10 \\
\hline
\end{tabular}
\end{center}

Find the relative frequency of the given data.

\bigskip
\noindent\textbf{436.} The frequency of defective products in 750 samples is shown in the following table. Find the relative frequency for the given table.

\begin{center}
\begin{tabular}{|c|c|c|c|c|c|c|c|c|c|}
\hline
No.\ of Defective per sample & 0 & 1 & 2 & 3 & 4 & 5 & 6 & 7 & 8 \\
\hline
No.\ of Sample & 120 & 40 & 94 & 85 & 105 & 50 & 40 & 66 & 50 \\
\hline
\end{tabular}
\end{center}

\bigskip
\noindent\textbf{437.} A quiz competition on general knowledge is conducted. The number of correct answers out of 5 questions for 100 sets of questions is given below.

\begin{center}
\begin{tabular}{|c|c|c|c|c|c|c|}
\hline
$x$ & 0 & 1 & 2 & 3 & 4 & 5 \\
\hline
$f$ & 10 & 23 & 15 & 25 & 18 & 9 \\
\hline
\end{tabular}
\end{center}

Find the relative frequencies for the given data.

\bigskip
\noindent\textbf{438.} A survey was conducted among 108 students of a class, asking about their favourite food. The responses are as under:

\begin{center}
\begin{tabular}{|l|c|c|c|c|c|}
\hline
Name of food item & Biryani & Fresh Juice & Chicken & Bar.B.Q & Sweets \\
\hline
No.\ of students & 40 & 07 & 21 & 15 & 25 \\
\hline
\end{tabular}
\end{center}

\begin{enumerate}[label=(\roman*)]
\item How many percentage of students like biryani?
\item How many percentage of students like chicken?
\item Which food is the least liked by the students?
\item Which food is the most preferred by the students?
\end{enumerate}

\bigskip
\noindent\textbf{439.} Find the expected frequencies of the given data if the experiment is repeated 200 times.

\begin{center}
\begin{tabular}{|c|c|c|c|c|c|c|c|}
\hline
$x$ & 0 & 1 & 2 & 3 & 4 & 5 & 6 \\
\hline
$P(x)$ & 0.11 & 0.21 & 0.17 & 0.18 & 0.09 & 0.17 & 0.07 \\
\hline
\end{tabular}
\end{center}

\bigskip
\noindent\textbf{440.} The probability of getting 5 sixes while tossing six dice is $\dfrac{2}{8}$. The dice are rolled 200 times. How many times would you expect it to show 5 sixes?

\bigskip
\noindent\textbf{441.} Find the relative frequency of the given data.

\begin{center}
\begin{tabular}{|c|c|c|c|c|c|c|c|}
\hline
$x$ & 2 & 3 & 4 & 5 & 6 & 7 & 8 \\
\hline
$f$ & 3 & 5 & 6 & 9 & 10 & 8 & 2 \\
\hline
\end{tabular}
\end{center}

\bigskip
\noindent\textbf{442.} Abdul Rehman obtained different marks in different subjects out of 100 marks. The details are as under:

\begin{center}
\begin{tabular}{|l|c|c|c|c|c|c|}
\hline
Subjects & English & Urdu & Islamiyat & Mathematics & Science & Computer Science \\
\hline
Marks obtained & 75 & 80 & 72 & 95 & 81 & 85 \\
\hline
\end{tabular}
\end{center}

Find the relative frequency of the above given data.

\bigskip
\noindent\textbf{443.} Six fair dice are rolled 50 times. The probability of occurrence of different numbers of sixes is given below. Find the expected frequency of the following data:

\begin{center}
\begin{tabular}{|c|c|c|c|c|c|c|c|}
\hline
$x$ & 0 & 1 & 2 & 3 & 4 & 5 & 6 \\
\hline
$P(x)$ & 0.09 & 0.10 & 0.12 & 0.24 & 0.10 & 0.20 & 0.15 \\
\hline
\end{tabular}
\end{center}

Find the expected frequency of occurrence of each six.

\bigskip
\noindent\textbf{444.} An urn contains 10 red balls, 5 green balls and 8 blue balls. Find the probability of selecting at random:
\begin{enumerate}[label=(\roman*)]
\item a green ball
\item a red ball
\item a blue ball
\item not a red ball
\item not a green ball
\end{enumerate}

\bigskip
\noindent\textbf{445.} Three coins are tossed together. What is the probability of getting:
\begin{enumerate}[label=(\roman*)]
\item exactly three heads
\item at least two tails
\item not at least two heads
\item not exactly two heads
\end{enumerate}

\bigskip
\noindent\textbf{446.} A card is drawn from a well-shuffled pack of 52 playing cards. What will be the probability of getting:
\begin{enumerate}[label=(\roman*)]
\item a king or jack of red colour
\item not a 2 of club and spade
\end{enumerate}

\bigskip
\noindent\textbf{447.} Six coins are tossed 600 times. The number of occurrences of tails is recorded and shown in the table given below:

\begin{center}
\begin{tabular}{|c|c|c|c|c|c|c|c|}
\hline
No.\ of Tails & 0 & 1 & 2 & 3 & 4 & 5 & 6 \\
\hline
Frequency & 110 & 90 & 105 & 80 & 76 & 123 & 16 \\
\hline
\end{tabular}
\end{center}

Find the relative frequency of the given table.

\bigskip
\noindent\textbf{448.} From a lot containing 25 items, 8 items are defective. Find the relative frequency of non-defective items; also find the expected frequency of non-defective items.


\end{document}
